\documentclass[12pt,a4paper]{article}
\usepackage[a4paper,margin=25mm]{geometry}
\usepackage[T1]{fontenc}
\usepackage{lmodern,microtype,setspace}
\usepackage{amsmath,amssymb,amsthm,mathtools,bm}
\usepackage{booktabs,tabularx,array,threeparttable,longtable,pdflscape,needspace}
\usepackage[round,authoryear]{natbib}
\usepackage[hidelinks]{hyperref}
\hypersetup{pdftitle={Supplementary material: Squared assignment correlations and randomisation tests of weak null hypotheses under covariate-adaptive allocation},pdfauthor={Masahiro Kojima}}
\newtheorem{theorem}{Theorem}
\newtheorem{proposition}{Proposition}
\newtheorem{corollary}{Corollary}
\newtheorem{lemma}{Lemma}
\theoremstyle{definition}
\newtheorem{assumption}{Assumption}
\newtheorem{definition}{Definition}
\theoremstyle{remark}\newtheorem{remark}{Remark}
\renewcommand{\thetheorem}{S\arabic{theorem}}
\renewcommand{\theproposition}{S\arabic{proposition}}
\renewcommand{\thecorollary}{S\arabic{corollary}}
\renewcommand{\thelemma}{S\arabic{lemma}}
\renewcommand{\theassumption}{S\arabic{assumption}}
\renewcommand{\thedefinition}{S\arabic{definition}}
\renewcommand{\theremark}{S\arabic{remark}}
\numberwithin{equation}{section}
\newcommand{\E}{\operatorname{E}}
\newcommand{\Pbb}{\mathbb P}
\newcommand{\Pp}{\mathbb P}
\newcommand{\Var}{\operatorname{Var}}
\newcommand{\Cov}{\operatorname{Cov}}
\newcommand{\Rcal}{\mathcal R}
\newcommand{\diag}{\operatorname{diag}}
\newcommand{\sg}{\operatorname{sgn}_0}
\newcommand{\R}{\mathbb R}
\newcommand{\Z}{\mathbb Z}
\newcommand{\one}{\mathbf1}
\newcommand{\norm}[1]{\left\lVert#1\right\rVert}
\newcommand{\trans}{^{\mathsf T}}
\setlength{\emergencystretch}{2em}
\allowdisplaybreaks[2]
\clubpenalty=10000
\widowpenalty=10000
\begin{document}
\setstretch{1.72}
\begin{center}
{\Large\bfseries Supplementary material for\\
Squared assignment correlations and randomisation tests\\ of weak null hypotheses under covariate-adaptive allocation}
\end{center}
\renewcommand{\thesection}{\Alph{section}}
This supplement gives the covariance constructions and proofs for the repeated-sampling and conditional randomisation results, verifies the allocation limits for stochastic Pocock--Simon minimisation, and records the complete numerical results. Appendix~H gives additional allocation examples, including stratified permuted blocks and stratified Efron biased coins. Appendix~I proves the results on the sign of the variance difference under minimisation (Section 4.3 of the main text) and describes the deterministic computations. Appendix~J reports the null-value study with heterogeneous effects, including the corrected test, under the main-study biased coin and under a strong biased coin. The minimisation-specific allocation theorems use the original equal-weight range method, including overall balance. Binary minimisation factors are assumed in that verification; the inferential arguments are stated for any two-arm allocation design satisfying the general condition below. Appendix~K reports the extended paired evaluation of the covariance reconstruction. The statistical endpoints may be continuous or binary.

\section{Covariance construction and allocation conditions}
\label{sec:supp-covariance}

\subsection{Exact projection and covariance completion}
Let $\bm H$ be the $n\times J$ joint-stratum indicator matrix and let
\[
\bm N=\bm H^\top\bm H=\diag(n_1,\ldots,n_J),
\qquad
\bm P_H=\bm H\bm N^+\bm H^\top,
\qquad
\bm M_H=\bm I_n-\bm P_H.
\]
For a centred fixed score $\bm r$, define
\[
\bm D=\bm H^\top\bm Z,
\qquad
\bar{\bm r}=\bm N^+\bm H^\top\bm r,
\qquad
\bm e=\bm M_H\bm r.
\]

\begin{proposition}\label{prop:supp-decomp}
For $T_n(\bm r)=\bm Z^\top\bm r/2$, the score satisfies
\[
\bm r=\bm H\bar{\bm r}+\bm e,
\qquad
\bm H^\top\bm e=\bm0,
\]
and
\begin{equation}\label{eq:supp-decomp}
T_n(\bm r)
=\frac12\bar{\bm r}^{\top}\bm D
+\frac12\bm Z^\top\bm e.
\end{equation}
\end{proposition}

\begin{proof}
Because $\bm N\bm N^+\bm H^\top\bm r=\bm H^\top\bm r$,
\[
\bm H^\top\bm e
=\bm H^\top\bm r-\bm N\bm N^+\bm H^\top\bm r
=\bm0.
\]
Substitution of $\bm r=\bm H\bar{\bm r}+\bm e$ into $T_n(\bm r)=\bm Z^\top\bm r/2$ gives \eqref{eq:supp-decomp}.
\end{proof}

Let $\bm\Omega$ be a positive-semidefinite approximation to $\Cov_{\Rcal}[\bm D]$ for the realised counts, with zero rows and columns for empty joint strata. Consider
\begin{equation}\label{eq:supp-covclass}
\bm\Sigma_Z(\kappa)
=\bm H\bm N^+\bm\Omega\bm N^+\bm H^\top
+\kappa\bm M_H.
\end{equation}

\begin{proposition}\label{prop:supp-completion}
For every $\kappa\ge0$, $\bm\Sigma_Z(\kappa)$ is positive semidefinite and $\bm H^\top\bm\Sigma_Z(\kappa)\bm H=\bm\Omega$ on the nonempty joint strata. If $J_+$ is the number of nonempty strata and $n>J_+$, the unique member of \eqref{eq:supp-covclass} that has trace $n$ and is a scalar multiple of the identity on the residual subspace has
\begin{equation}\label{eq:supp-kappa}
\kappa^{\rm raw}
=\frac{n-\operatorname{tr}(\bm N^+\bm\Omega)}{n-J_+}.
\end{equation}
If $\kappa^{\rm raw}<0$, no matrix in this isotropic family can be both positive semidefinite and have trace $n$; $\max(\kappa^{\rm raw},0)$ is the scalar projection onto the admissible set. If $n=J_+$, then $\bm M_H=\bm0$, the residual subspace is trivial and the value of $\kappa$ is immaterial; trace matching is then possible only if $\operatorname{tr}(\bm N^+\bm\Omega)=n$.
\end{proposition}

\begin{proof}
Both terms in \eqref{eq:supp-covclass} are positive semidefinite when $\kappa\ge0$. Because $\bm\Omega$ is supported on the nonempty strata,
\[
\bm H^\top\bm H\bm N^+\bm\Omega\bm N^+\bm H^\top\bm H
=\bm N\bm N^+\bm\Omega\bm N^+\bm N
=\bm\Omega,
\]
and $\bm H^\top\bm M_H\bm H=\bm0$. The trace identities $\operatorname{tr}\{\bm H\bm N^+\bm\Omega\bm N^+\bm H^\top\}=\operatorname{tr}(\bm N^+\bm\Omega)$ and $\operatorname{tr}(\bm M_H)=n-J_+$ give \eqref{eq:supp-kappa} when $n>J_+$. Uniqueness follows because isotropy on the range of $\bm M_H$ leaves only the scalar $\kappa$. When $n=J_+$, every nonempty stratum contains exactly one participant, so $\bm P_H=\bm I_n$ and $\bm M_H=\bm0$.
\end{proof}

For $\bm r=\bm H\bar{\bm r}+\bm e$,
\begin{equation}\label{eq:supp-quadratic}
\frac14\bm r^\top\bm\Sigma_Z(\kappa)\bm r
=\frac14\{\bar{\bm r}^\top\bm\Omega\bar{\bm r}
+\kappa\bm e^\top\bm e\}.
\end{equation}
This construction is a finite-sample working approximation that reproduces the estimated covariance of the joint-stratum imbalances and, when $\kappa^{\rm raw}\ge0$, matches the total trace. It is not intended to recover the full finite-sample covariance matrix. The large-sample results below require only consistency of the induced quadratic form for the score under consideration.

\subsection{General allocation condition for the inference results}\label{sec:supp-general-allocation}
Let $S_i\in\{1,\ldots,J\}$ be independent and identically distributed profiles with fixed $J$ and probabilities $\pi_s>0$, let $\bm h(S_i)$ be the corresponding indicator vector, and write $\bm\Pi=\diag(\bm\pi)$ and $\bm\Sigma_S=\bm\Pi-\bm\pi\bm\pi^\top$. For a two-arm sequential randomisation design, generate an observed allocation sequence $\bm Z^{(0)}$ and two additional sequences $\bm Z^{(1)},\bm Z^{(2)}$ independently conditional on the same ordered profile sequence, using the same prespecified rule. The design uses only the ordered profiles and its own random numbers, and its conditional path law is invariant under interchange of the treatment labels (label symmetry). The random numbers are independent of the remaining participant variables conditional on the complete profiles, as specified in Assumption~\ref{ass:supp-participants}. Define
\[
 \bm R_n=n^{-1/2}\sum_{i=1}^n\{\bm h(S_i)-\bm\pi\},\qquad
 \bm D_n^{(m)}=\sum_{i=1}^n\bm h(S_i)Z_i^{(m)},
\]
\[
 \bm G_n^{(ab)}=n^{-1/2}\sum_{i=1}^n\bm h(S_i)Z_i^{(a)}Z_i^{(b)},\qquad
 \bm K_n^{(012)}=n^{-1/2}\sum_{i=1}^n\bm h(S_i)Z_i^{(0)}Z_i^{(1)}Z_i^{(2)}.
\]

\begin{assumption}\label{ass:supp-general-allocation}
Retain the fixed finite profile space, independent and identically distributed profiles, conditional independence of allocation repetitions, and label symmetry specified above. For finite positive-semidefinite matrices $\bm\Sigma_D$ and $\bm\Psi_D$,
\begin{align*}
 &\left(\bm R_n,\frac{\bm D_n^{(0)}}{\sqrt n},\frac{\bm D_n^{(1)}}{\sqrt n},
 \frac{\bm D_n^{(2)}}{\sqrt n},\bm G_n^{(01)},\bm G_n^{(02)},\bm G_n^{(12)}\right)\\
 &\hspace{12mm}\Rightarrow
 \mathcal N\!\left(0,\diag(\bm\Pi-\bm\pi\bm\pi^\top,
 \bm\Sigma_D,\bm\Sigma_D,\bm\Sigma_D,
 \bm\Psi_D,\bm\Psi_D,\bm\Psi_D)\right).
\end{align*}
The corresponding second moments converge. The normalised fourth moments of the imbalance and product-process blocks are uniformly bounded, and $\bm K_n^{(012)}=O_p(1)$.
\end{assumption}

The inference arguments in Appendices B and C use Assumption~\ref{ass:supp-general-allocation}, the structural allocation properties just specified and the participant conditions; they do not otherwise use minimisation. The following subsection verifies the condition for the stochastic Pocock--Simon range method. Independent assignments, fixed-size stratified permuted blocks and stratified Efron biased coins are treated in Appendix~H.

\subsection{Verification for the Pocock--Simon range method}\label{sec:supp-allocation}
\begin{assumption}\label{ass:supp-design}
There are $F\ge1$ fixed binary minimisation factors. The profiles $S_i$ are independent and identically distributed over the $J=2^F$ combinations, with probabilities $\pi_s>0$. Allocation starts from overall and marginal balance. The range method sums the absolute overall imbalance and the absolute imbalances at the incoming factor levels, with equal weights. The sign with the smaller candidate sum is selected with the fixed probability $p_{\rm bc}\in(1/2,1)$; a tie is resolved with probability $1/2$. Each copy is driven by fresh independent Uniform$(0,1)$ random numbers, independent of the complete profile sequence; the streams are independent across copies.
\end{assumption}
Write $\bm h(s)$ for the joint-stratum indicator and $\bm\Pi=\diag(\bm\pi)$. For copy $m$, put
\[
 \bm D_n^{(m)}=\sum_{i=1}^n\bm h(S_i)Z_i^{(m)},\quad
 \bm R_n=n^{-1/2}\sum_{i=1}^n\{\bm h(S_i)-\bm\pi\},
\]
\[
 \bm G_n^{(ab)}=n^{-1/2}\sum_{i=1}^n\bm h(S_i)Z_i^{(a)}Z_i^{(b)},\quad
 \bm K_n^{(012)}=n^{-1/2}\sum_{i=1}^n\bm h(S_i)Z_i^{(0)}Z_i^{(1)}Z_i^{(2)}.
\]
The following proof uses a sign-coded profile $U_i$ in one-to-one correspondence with $S_i$. Thus $h(U_i)=\bm h(S_i)$, $\Pi=\bm\Pi$, and the unbold symbols $D_n,R_n,G_n$ in the allocation argument denote the same vectors as their bold counterparts above.

Let $F\ge1$ be fixed, let $\mathcal U=\{-1,1\}^F$, and let $U_i\in\mathcal U$ be iid with $\Pp(U_i=u)=\pi_u>0$. Write $J=2^F$, $d=F+1$, and
\[
 a(u)=(1,u_1,\ldots,u_F)\trans,\qquad
 \mathsf L=[a(u):u\in\mathcal U],\qquad
 \Pi=\diag(\pi_u),\qquad K=\mathsf L\Pi \mathsf L\trans.
\]
Full support implies that $K$ is positive definite: if $v\trans K v=0$, then the affine function $v_0+\sum_jv_ju_j$ vanishes on every vertex of the cube, so $v=0$.

Let $Z_i\in\{-1,1\}$ denote the treatment sign. Define the contrast state
\[
 X_n=(X_{n0},X_{n1},\ldots,X_{nF})\trans
       =\sum_{i=1}^n Z_i a(U_i),\qquad X_0=0.
\]
Here $X_{n0}$ is overall imbalance and $X_{nj}$ is the difference between the two level-specific imbalances for factor $j$. Consequently, the imbalance at the incoming level $u_j$ is $(X_{n0}+u_jX_{nj})/2$. This is an invertible reparameterisation of the feasible marginal-imbalance state; the allocation rule is not changed.

Write $p=p_{\rm bc}\in(1/2,1)$, $q=1-p$, and $\eta=2p-1$. Throughout, $\sg(t)$ is $-1$, $0$ or $1$ according as $t<0$, $t=0$ or $t>0$. Given $X_{i-1}=x$ and $U_i=u$, apply the comparison of the range method, using probability $p$ for the sign with the smaller candidate score and probability $1/2$ at a tie. Define
\begin{equation}\label{eq:fx}
 f_x(u)=\sg\left\{\sg(x_0)+\sum_{j=1}^F\sg(x_0+u_jx_j)\right\}.
\end{equation}
Then
\begin{equation}\label{eq:mean}
 X_i=x+Z_i a(u),\qquad \E[Z_i\mid x,u]=-\eta f_x(u).
\end{equation}
This follows because $|b+1|-|b-1|=2\sg(b)$ for integer $b$. The identity applies because the imbalances are integer-valued.

Let $P$ be the state transition kernel on the class reachable from $0$. The same construction with $k$ independent treatment-randomisation streams and one common iid profile stream defines $P_k$. All state transitions have their reverse with positive probability. Thus the reachable class is irreducible. Its period is two: $X_{n0}$ changes parity at each step, and a profile followed by the same profile with all treatment signs reversed gives a two-step return. The two-step skeleton is irreducible on each parity class and has a self-loop at every state. Any admissible path in that skeleton can therefore be padded to a multiple of any fixed positive integer. Hence every even-step skeleton is irreducible and aperiodic on its reachable even class.

\subsubsection{Drift bounds for the range method}\label{sec:quadratic}\label{sec:two}
Define
\begin{equation}\label{eq:QH}
 \mathcal V_0(x)=x\trans K^{-1}x,\qquad \mathcal V(x)=1+\mathcal V_0(x),\qquad
 \mathcal D(x)=x\trans K^{-1}\sum_{u\in\mathcal U}\pi_u a(u)f_x(u).
\end{equation}
The quadratic Lyapunov function is used only in the proof; the allocation criterion remains the range method based on absolute imbalances.

\begin{lemma}\label{lem:drift}
For every reachable state,
\begin{equation}\label{eq:quadratic-drift}
 P\mathcal V_0(x)-\mathcal V_0(x)=d-2\eta \mathcal D(x).
\end{equation}
Also $\mathcal D(-x)=\mathcal D(x)$ and the one-step change of $\sqrt{\mathcal V(x)}$ is bounded by a constant depending only on $\pi$ and $F$.
\end{lemma}
\begin{proof}
Expand $\mathcal V_0(x+Za)-\mathcal V_0(x)=2Zx\trans K^{-1}a+a\trans K^{-1}a$ and use \eqref{eq:mean}. The final term averages to $\operatorname{tr}(K^{-1}K)=d$. Equation \eqref{eq:fx} gives $f_{-x}=-f_x$, proving evenness. Finally,
\[
 |\sqrt{1+\mathcal V_0(x+Za)}-\sqrt{1+\mathcal V_0(x)}|
 \le \sqrt{a\trans K^{-1}a}
 \le C_\Delta:=\max_{u\in\mathcal U}\sqrt{a(u)\trans K^{-1}a(u)}.
\]
The first inequality is the reverse triangle inequality after adjoining a coordinate equal to one.
\end{proof}

The following two bounds will be sufficient:
\begin{align}
 \mathcal D(x)&\ge c_0|x_0|\quad\text{for every reachable }x, \label{eq:H1}\\
 \mathcal D(0,x_1,\ldots,x_F)&\ge c_1\sum_{j=1}^F|x_j|\quad\text{whenever }x_0=0, \label{eq:H2}
\end{align}
where $c_0,c_1>0$. We first derive these bounds for all positive two-factor laws, and then give a finite, exactly checkable sufficient condition for higher dimensions. These bounds are derived below for the admissible profile laws.

For two binary factors, these bounds hold for every strictly positive joint profile distribution.
\begin{theorem}\label{thm:two}
If $F=2$ and all four probabilities $\pi_{uv}$ are strictly positive, then \eqref{eq:H1}--\eqref{eq:H2} hold with $c_0=1$ and $c_1=1/2$. Independence of the two factors is not required.
\end{theorem}
\begin{proof}
For a single profile $U=(U_1,U_2)^\top$, let
\[
 \gamma=(\gamma_0,\gamma_1,\gamma_2)\trans
     =K^{-1}\E[a(U)U_1U_2],\qquad
 w_{uv}=\frac{\pi_{uv}^{-1}}{\sum_{r,t\in\{-1,1\}}\pi_{rt}^{-1}}.
\]
The weighted linear projection of $uv$ on $(1,u,v)$ satisfies
\begin{equation}\label{eq:gamma}
 \gamma=-\sum_{u,v}w_{uv}(uv,v,u)\trans.
\end{equation}
To verify this without any independence assumption, the proposed coefficients give the residual
\[
 uv-\gamma_0-\gamma_1u-\gamma_2v=4w_{uv}uv.
\]
Its $\pi$-weighted inner products with $1,u,v$ vanish, since $\pi_{uv}w_{uv}$ is constant. Hence it is the unique regression residual. In particular $|\gamma_j|<1$, and for $s,t\in\{-1,1\}$,
\begin{equation}\label{eq:widentity}
 st\gamma_0+t\gamma_1+s\gamma_2=1-4w_{s,t}.
\end{equation}
Indeed, substituting \eqref{eq:gamma} into the left side gives $-3w_{s,t}+\sum_{(u,v)\ne(s,t)}w_{uv}$.

By evenness of $\mathcal D$, consider first $x_0=a>0$. Put $R=|x_1|$ and $C=|x_2|$, with $s=\sg(x_1)$ and $t=\sg(x_2)$ when the corresponding coordinate is nonzero. There are three cases.

\emph{Case 1: $f_x$ is constant.} If at least one of $R,C$ is less than $a$, or if both equal $a$, all values of the expression defining $f_x$ are positive. Thus $f_x=1$ and $\mathcal D(x)=a$, since $K^{-1}\E[a(U)]$ is the first coordinate vector.

\emph{Case 2: $R>a$ and $C>a$.} Directly on the four profiles,
\[
 f_x(u,v)=\tfrac12(1+su+tv-stuv).
\]
Using the linearity of projection and \eqref{eq:gamma},
\[
 \mathcal D(x)=\tfrac12\{(1-st\gamma_0)a+(1-t\gamma_1)R+(1-s\gamma_2)C\}.
\]
The coefficients of $R,C$ are positive. Replacing both by $a$ and applying \eqref{eq:widentity} therefore gives
\[
 \mathcal D(x)\ge\tfrac a2\{3-st\gamma_0-t\gamma_1-s\gamma_2\}
       =a(1+2w_{s,t})\ge a.
\]

\emph{Case 3: exactly one of $R,C$ is greater than $a$ and the other equals $a$.} The value of $f_x$ is zero at $(-s,-t)$ and one elsewhere, so
\[
 f_x(u,v)=\tfrac14(3+su+tv-stuv).
\]
Consequently,
\[
 \mathcal D(x)=\tfrac14\{(3-st\gamma_0)a+(1-t\gamma_1)R+(1-s\gamma_2)C\}
      \ge a(1+w_{s,t})\ge a.
\]
These cases exhaust $a>0$.

At $x_0=0$, if both coordinates are nonzero,
\[
 f_x(u,v)=\sg(su+tv)=\tfrac12(su+tv),\qquad
 \mathcal D(x)=\tfrac12(R+C).
\]
If only one is nonzero, $f_x=su$ or $tv$ and $\mathcal D(x)=R$ or $C$. At zero $\mathcal D=0$. This proves both bounds, including $\mathcal D\ge0$ everywhere.
\end{proof}


\subsubsection{A finite sufficient condition for any fixed number of binary factors}\label{sec:finite-condition}
The following sufficient condition uses expectations with exactly $2^F$ terms.

For $c\in\{-2,-1,0,1,2\}^F$, set
\begin{align}
 f_c(u)&=\sg\left\{1+\sum_{j=1}^F\sg(1+u_jc_j)\right\},\\
 \xi(c)&=K^{-1}\sum_u\pi_u a(u)f_c(u),\\
 L(c)&=\xi_0(c)+\sum_{j:c_j\ne0}\sg(c_j)\xi_j(c)
                  -\sum_{j:c_j=0}|\xi_j(c)|. \label{eq:finite-condition-base}
\end{align}
For $e\in\{-1,0,1\}^F\setminus\{0\}$, set
\[
 g_e(u)=\sg\left(\sum_je_ju_j\right),\qquad
 \zeta(e)=K^{-1}\sum_u\pi_u a(u)g_e(u).
\]
\begin{definition}\label{def:finite-condition}
A full-support profile law belongs to $\mathcal P_F$ if
\begin{align}
 c_0&:=\min_c L(c)>0,\label{eq:C0}\\
 \sg(c_j)\xi_j(c)&\ge0\quad\text{for all }c\text{ and }j\text{ with }|c_j|=2,\label{eq:C1}\\
 c_1&:=\min_{e,j:e_j\ne0} e_j\zeta_j(e)>0.\label{eq:C2}
\end{align}
Membership is checked by evaluating these finite sums in the profile probabilities.
\end{definition}
\begin{assumption}\label{ass:supp-profile}
Either $F=2$, with an arbitrary full-support joint profile law, or $\pi\in\mathcal P_F$ as defined above. No independence of the factors is imposed. The condition is checked on the profile probabilities before applying the limit theorem.
\end{assumption}


\begin{proposition}\label{prop:finite-condition}
Every $\pi\in\mathcal P_F$ satisfies \eqref{eq:H1}--\eqref{eq:H2} with the constants in \eqref{eq:C0} and \eqref{eq:C2}.
\end{proposition}
\begin{proof}
For $x_0>0$, assign category $c_j=-2,-1,0,1,2$ according as $x_j/x_0$ is less than $-1$, equal to $-1$, strictly between $-1$ and $1$, equal to $1$, or greater than $1$. This classification preserves every sign in \eqref{eq:fx}, so $f_x=f_c$ and $\mathcal D(x)=\xi(c)\trans x$.

For $|c_j|=2$, write $x_j=\sg(c_j)(x_0+t_j)$, $t_j>0$. For $|c_j|=1$, $x_j=c_jx_0$. For $c_j=0$, $|x_j|<x_0$. Thus
\[
 \mathcal D(x)\ge x_0L(c)+\sum_{j:|c_j|=2}t_j\sg(c_j)\xi_j(c)\ge c_0x_0.
\]
The use of $-|\xi_j|$ takes the infimum over an entire interior interval, not just a representative state. The ray conditions handle unbounded coordinates. Evenness handles $x_0<0$.

For $x_0=0$, use $e_j=\sg(x_j)$. Then $f_x=g_e$ and
\[
 \mathcal D(x)=\sum_{j:e_j\ne0}|x_j|e_j\zeta_j(e)\ge c_1\sum_j|x_j|.
\]
The zero state is immediate.
\end{proof}

For the numerical settings, the accompanying script uses Gauss--Jordan inversion in rational arithmetic. It clears denominators once and checks every inequality using integers. Probabilities expressed as decimals below mean the corresponding exact terminating rationals. No numerical tolerance, simulation, or state-space truncation enters these checks.

\Needspace{9\baselineskip}
For $F=5$ there are $5^5=3125$ positive-overall sign chambers, $2F5^{F-1}=6250$ ray inequalities, and $2F3^{F-1}=810$ nonzero-coordinate inequalities on the zero-overall hyperplane. The results are
\begin{center}
\small
\begin{tabular}{@{}>{\raggedright\arraybackslash}p{61mm}rrr@{}}
\toprule
Independent marginal probabilities & $\min L(c)$ & Minimum ray & $\min e_j\zeta_j(e)$\\
\midrule
$(.50,.40)$ & $1$ & $0$ & $1/2$\\
$(.50,.40,.30,.20,.10)$ & $1$ & $0$ & $134/625$\\
$(.466,.592,.287,.154,.300)$ & $1$ & $0$ & $159670787/625000000$\\
\bottomrule
\end{tabular}
\end{center}
Complete exact evaluations of the finite inequalities and all input probabilities are supplied in the accompanying computational files. The listed finite decimals define the model distributions used in these checks. An independent implementation uses a separate symbolic matrix inversion and explicit chamber-vertex enumeration.

The minimum ray coefficient among nonconstant sign functions is $3/2500$ for the main five-factor law and $3381147/1250000000$ for the trial-based law. Every zero ray coefficient belongs to a constant sign function. Such a function is identically one, so its coefficient vector is exactly $(1,0,\ldots,0)\trans$ for every positive profile law.

\begin{corollary}\label{cor:open}
Each of the two five-factor laws in the preceding table has an open neighbourhood in the full strictly positive probability simplex contained in $\mathcal P_5$. In particular, the theorem is not restricted to independent factors.
\end{corollary}
\begin{proof}
There are finitely many inequalities. The matrix $K^{-1}$ and hence all their left sides are continuous in $\pi$ on the full-support simplex. The base, nonconstant-ray, and zero-overall inequalities have strictly positive minima. Constant-sign zero rays remain identically zero for every $\pi$. The stated neighbourhood follows.
\end{proof}

As a further check, for either five-factor product law $\pi^0$, let $\mu_j=\E_{\pi^0}U_j$ and define
\[
 \pi_u=\pi^0_u\{1+\tfrac1{10}(u_1-\mu_1)(u_2-\mu_2)\}.
\]
This is a positive correlated law preserving every univariate margin. Exact checks give $c_0=1$ for both perturbations and
\[
 c_{1,\mathrm{main}}=\frac{165433}{773750},\qquad
 c_{1,\mathrm{trial}}=\frac{1911843769570152559}{7556024433125000000}.
\]
These two distributions provide explicit correlated examples satisfying the finite sufficient condition.

\noindent The finite condition is sufficient, not necessary. For example, the independent five-factor law with all prevalences $.01$ fails its base inequality; no assertion about instability follows from that failure.

\begin{proposition}\label{prop:alloc-box}
For independent binary factors with prevalences $p_1,\ldots,p_F$, every component of $\xi(c)$ and $\zeta(e)$ is multi-affine in the prevalences. If all $2^F$ corners of a closed prevalence box satisfy the finite condition with common constants $c_0,c_1>0$, then every point of that box satisfies it with the same constants. If all nonconstant-sign ray coefficients at its corners are strictly positive, every law in the box also has a full-simplex neighbourhood in $\mathcal P_F$.
\end{proposition}
\begin{proof}
For any fixed sign function $f$, independence gives its affine regression coefficients explicitly:
\[
 \xi_j=\tfrac12\{E[f(U)\mid U_j=1]-E[f(U)\mid U_j=-1]\},\qquad
 \xi_0=E f(U)-\sum_j(2p_j-1)\xi_j.
\]
Thus $\xi_j$ is multi-affine in all prevalences other than $p_j$, and $\xi_0$ is multi-affine in all prevalences. The same argument applies to $g_e$ and $\zeta(e)$. On a prevalence box, a multi-affine vector equals a convex combination of its corner vectors, with common tensor-product interpolation weights. The ray and zero-overall inequalities are linear in the corresponding vectors. The base functional $L$ is concave, being linear minus a sum of absolute values. Its value at an interpolated vector is therefore at least the interpolated corner lower bound. This proves the finite condition throughout the box. Strict positivity of all nonconstant rays is also preserved; constant-sign zero rays are identities for arbitrary positive laws. Continuity in the full profile simplex proves the neighbourhood assertion.
\end{proof}

The production trial code computes the stroke-severity and age prevalences from truncated-normal tails and uses ratios $117/408$ and $63/408$ for the next two factors; it does not assign all five displayed three-decimal values literally. To certify the model actually defined by those formulas, the companion verification script checks all 32 corners of each of the following boxes:
\begin{align*}
 \mathcal B_{\rm main}={}&[.499,.501]\times[.399,.401]\times[.299,.301]
                       \times[.199,.201]\times[.099,.101],\\
 \mathcal B_{\rm trial}={}&[.465,.467]\times[.591,.593]\times[.286,.288]
                        \times[.153,.155]\times[.299,.301].
\end{align*}
The exact common constants are
\[
 c_{0,\rm main}=c_{0,\rm trial}=1,\qquad
 c_{1,\rm main}=\frac{106599553003}{500000000000},\quad
 c_{1,\rm trial}=\frac{63596532191}{250000000000}.
\]
All ray coefficients are nonnegative at every corner, with strict positivity for nonconstant sign functions. Thus both boxes, including neighbourhoods allowing correlated profiles, are admissible.

The trial formula is also enclosed using rational arithmetic rather than an unquantified numerical approximation. Let $q=\Phi^{-1}(.75)$ and
\[
 t(c;\mu,s,l,h)=\frac{\Phi((h-\mu)/s)-\Phi((c-\mu)/s)}
                         {\Phi((h-\mu)/s)-\Phi((l-\mu)/s)}.
\]
The code's model probabilities are
\begin{align*}
 p_1&=t(17.5;17,[20-(201\cdot13+207\cdot12)/408]/(2q),5,30),\\
 p_2&=t(69.5;72.5,16.5/(2q),18,100),\quad
 p_3=39/136,\quad p_4=21/136,\quad p_5=.3.
\end{align*}
Exact interval calculations give
\[
 .466244729212<p_1<.466244729215,\qquad
 .591867113705<p_2<.591867113709.
\]
These bounds and the three rational probabilities place the production model strictly inside $\mathcal B_{\rm trial}$. The interval verification uses the alternating series for $\arctan$ with Machin's identity $\pi=16\arctan(1/5)-4\arctan(1/239)$, integer square-root enclosures, and the integrated power series for the normal density. Its bracket $.67448975019<q<.67448975020$ is verified by bounding the two normal cumulative probabilities on either side of $.75$. For the normal series it retains 141 terms and bounds the alternating remainder after the terms are decreasing. Every comparison uses rational numbers. These calculations certify the probability formulas, not the implementation of a floating-point special-function library; the probability box leaves an explicit margin around those formulas.

\subsubsection{Return times and joint allocation limits}\label{sec:stability}\label{sec:clt}
The next argument obtains the return moments directly from the two drift bounds.

\begin{lemma}\label{lem:polynomial}
Assume \eqref{eq:H1}--\eqref{eq:H2}. For every integer $r\ge2$, there exist an even integer $T_r$, constants $\delta_r>0$, $b_r<\infty$, and a finite set $C_r$ such that
\begin{equation}\label{eq:polynomial}
 P^{T_r}\mathcal V^r(x)\le \mathcal V^r(x)-\delta_r\mathcal V(x)^{r-1}+b_r\one_{C_r}(x).
\end{equation}
All constants may depend on $\pi,p,F,r$.
\end{lemma}
\begin{proof}
For one step, $|\mathcal V(X_1)-\mathcal V(x)|\le K_1\sqrt{\mathcal V(x)}$ for a finite $K_1$. The binomial expansion of the integer power, together with \eqref{eq:quadratic-drift}, therefore gives a finite constant $C_r'>0$ such that
\begin{equation}\label{eq:power-one}
 P\mathcal V^r(x)-\mathcal V^r(x)
 \le \mathcal V(x)^{r-1}\{C_r'-2r\eta \mathcal D(x)\}.
\end{equation}
Indeed, the linear term is $r\mathcal V^{r-1}(d-2\eta \mathcal D)$ and each term of order $j\ge2$ is bounded in absolute value by a constant times $\mathcal V^{r-j/2}\le \mathcal V^{r-1}$.

Fix an even $T$ temporarily. Lemma~\ref{lem:drift} gives
$|\sqrt{\mathcal V(X_j)}-\sqrt{\mathcal V(x)}|\le C_\Delta T$ for $j\le T$ on every path. Hence, for all sufficiently large $\mathcal V(x)$,
\[
 \tfrac12\mathcal V(x)^{r-1}\le \mathcal V(X_j)^{r-1}\le 2\mathcal V(x)^{r-1},\qquad j\le T.
\]
Since $\mathcal D\ge0$, summing \eqref{eq:power-one} and using the tower property gives
\begin{equation}\label{eq:telescope}
 P^T\mathcal V^r(x)-\mathcal V^r(x)
 \le \mathcal V(x)^{r-1}\left\{2C_r'T-r\eta\sum_{j=0}^{T-1}\E_x\mathcal D(X_j)\right\}.
\end{equation}

If $|x_0|\ge T$, the unit changes of the overall coordinate imply, on every path,
\[
 \sum_{j=0}^{T-1}\mathcal D(X_j)\ge c_0\sum_{j=0}^{T-1}(|x_0|-j)
 \ge c_0T(T+1)/2.
\]
Choose the even integer $T=T_r$ sufficiently large that
\[
 r\eta c_0T(T+1)/2\ge 2C_r'T+1.
\]
This makes the braces in \eqref{eq:telescope} at most $-1$ in this case.

If $|x_0|<T$, put $m=|x_0|\in\{0,\ldots,T-1\}$. Prescribe the first $m$ treatment signs to move the overall coordinate toward zero. Each prescribed sign has conditional probability at least $q>0$, whatever the incoming profiles. Thus this event has probability at least $q^m\ge q^{T-1}$; for $m=0$ it is the whole sample space. At time $m$ on this event, $X_{m0}=0$.

Let $\lambda_{\min}(K)>0$. Bounded coordinate increments and norm equivalence give
\[
 \sum_{j=1}^F|X_{mj}|
 \ge\sqrt{\lambda_{\min}(K)}\sqrt{\mathcal V(x)-1}-(F+1)T.
\]
For fixed $T$ and sufficiently large $\mathcal V(x)$, the right side exceeds $c_*\sqrt{\mathcal V(x)}$ for some $c_*>0$. By \eqref{eq:H2}, and by nonnegativity of $\mathcal D$ off the prescribed event,
\[
 \sum_{j=0}^{T-1}\E_x\mathcal D(X_j)
 \ge q^{T-1}c_1c_*\sqrt{\mathcal V(x)}.
\]
The braces in \eqref{eq:telescope} are again at most $-1$ once $\mathcal V(x)$ is large enough. The two cases cover every integer overall coordinate.

Thus $P^T\mathcal V^r-\mathcal V^r\le -\mathcal V^{r-1}$ outside a sublevel set of $\mathcal V$. This set is finite because $K$ is positive definite and the state space is a lattice. The $T$-step expectation is finite on this set because increments are bounded. Absorbing its maximum positive discrepancy into $b_r$ proves \eqref{eq:polynomial}, with $\delta_r=1$ or any smaller positive constant.
\end{proof}

\begin{lemma}\label{lem:alloc-return}
Let $P_0$ be an irreducible, aperiodic Markov kernel on a countable subset of a finite-dimensional lattice, with bounded increments. Suppose $W\ge1$ is finite and proper and, for an integer $r\ge2$,
\[
 P_0W\le W-cW^{1-1/r}+b\mathbf1_C
\]
for a finite set $C$ and $c>0$. Then every accessible singleton has finite $r$th hitting and return moments from every fixed state. The bounds are uniform over a family of kernels with common transition support if the displayed constants and finite set are common, $W$ is uniformly bounded on finite sets, the increment bound is common, and each fixed finite admissible path has a probability bounded away from zero throughout the family.
\end{lemma}
\begin{proof}
Replace $c$ by $\min(c,1/2)$ if necessary and enlarge $C$ to contain the target state. Then $0<cW^{-1/r}<1$ everywhere, with the same choice of $c$ for the uniform version. Let $F_j=W^{j/r}$ for $0\le j\le r$, with $F_0=1$. Jensen's inequality and concavity imply, off $C$,
\[
 P_0F_j\le(P_0W)^{j/r}
 \le W^{j/r}(1-cW^{-1/r})^{j/r}
 \le F_j-a_jF_{j-1},\qquad a_j=jc/r.
\]
For $\tau_C=\inf\{n\ge0:X_n\in C\}$, stopping and telescoping at $\tau_C\wedge N$, then applying monotone convergence, gives
\[
 E_x\sum_{t<\tau_C}F_{j-1}(X_t)\le F_j(x)/a_j.
\]
In particular $E_x\tau_C\le F_1(x)/a_1$. Set $A_0=1$ and $A_1=1/a_1$. For $j\ge2$, the difference-of-powers identity gives $\tau_C^j\le j\sum_{t<\tau_C}(\tau_C-t)^{j-1}$. Apply Tonelli's theorem and condition at time $t$ on $\{t<\tau_C\}$; the remaining hitting time has the law of $\tau_C$ started from $X_t$. Induction therefore yields
\begin{align*}
 E_x\tau_C^j
 &\le jE_x\sum_{t<\tau_C}(\tau_C-t)^{j-1}\\
 &\le jA_{j-1}E_x\sum_{t<\tau_C}F_{j-1}(X_t)
 \le A_jF_j(x),\qquad A_j=jA_{j-1}/a_j.
\end{align*}
The assertion is immediate for starting states in $C$, where $\tau_C=0$. This proves the hitting-moment bound without assuming it in advance.

For a target state $o\in C$, choose a common integer $L\ge1$ and $\epsilon>0$ such that $P_0^L(x,\{o\})\ge\epsilon$ for every $x\in C$. Such $L,\epsilon$ exist by finiteness, irreducibility and aperiodicity; in our applications a self-loop allows explicit padding of each path. Starting in $C$, run $L$ steps and declare success when the endpoint is $o$. Otherwise return to $C$ and repeat. Bounded increments put each endpoint of an $L$-step attempt in a finite set. The preceding hitting bound therefore makes the conditional $r$th moment of each attempt plus its return to $C$ at most a fixed $B_r$. The probability that attempt $j$ is reached is at most $(1-\epsilon)^{j-1}$. If $D_j$ is its duration, Minkowski's inequality gives
\[
 \left\|\sum_{j\ge1}D_j\mathbf1\{j\text{ is reached}\}\right\|_{L^r}
 \le B_r^{1/r}\sum_{j\ge1}(1-\epsilon)^{(j-1)/r}<\infty.
\]
The actual hitting time, or positive return time when starting at $o$, is no larger than this protocol duration. An initial hit of $C$ handles any other fixed starting state. The same argument proves the stated uniform version. This is a specialised proof of the polynomial-drift implications studied by \citet{DoucEtAl2004,JarnerRoberts2002}.
\end{proof}

\begin{theorem}\label{thm:recurrence}
Suppose $F=2$ with arbitrary positive joint probabilities, or more generally $\pi\in\mathcal P_F$. For every fixed $k\ge1$, the joint state chain for $k$ sequences with common profiles is positive recurrent on its reachable class. Its first return time $\tau_k$ to the all-zero state, when started there, satisfies
\begin{equation}\label{eq:all-moments}
 \E_0[\tau_k^a]<\infty\qquad\text{for every finite }a>0.
\end{equation}

\end{theorem}
\begin{proof}
Theorem~\ref{thm:two} or Proposition~\ref{prop:finite-condition} supplies the drift bounds. For integer $r\ge2$, define on the $k$-copy chain
\[
 \mathcal V_{r,k}(x^{(1)},\ldots,x^{(k)})=\sum_{m=1}^k \mathcal V(x^{(m)})^r.
\]
Conditional on the full current state, the marginal $T_r$-step transition of any copy is exactly $P^{T_r}$: future profiles are iid and other copies do not enter its treatment rule. Thus Lemma~\ref{lem:polynomial} yields
\begin{align*}
 P_k^{T_r}\mathcal V_{r,k}
 &\le \mathcal V_{r,k}-\delta_r\sum_m\mathcal V(x^{(m)})^{r-1}+kb_r\\
 &\le \mathcal V_{r,k}-\delta_r\mathcal V_{r,k}^{1-1/r}+kb_r.
\end{align*}
Outside a sufficiently large finite sublevel set $C_{r,k}$, the last positive constant is at most $(\delta_r/2)\mathcal V_{r,k}^{1-1/r}$. Hence
\begin{equation}\label{eq:copy-polynomial}
 P_k^{T_r}\mathcal V_{r,k}
 \le \mathcal V_{r,k}-\tfrac{\delta_r}{2}\mathcal V_{r,k}^{1-1/r}
             +b_{r,k}\one_{C_{r,k}}.
\end{equation}

Apply Lemma~\ref{lem:alloc-return} to the even-step kernel $P_0=P_k^{T_r}$ on its reachable even class and to \eqref{eq:copy-polynomial}. Irreducibility and aperiodicity were established in this subsection; all finite sets are petite there. The all-zero state is accessible. If $\tau_k^{(T_r)}$ is its return time in this skeleton, the original first return satisfies $\tau_k\le T_r\tau_k^{(T_r)}$. Consequently $\E_0\tau_k^r<\infty$. Since $r$ is arbitrary, \eqref{eq:all-moments} follows. Finite mean return and irreducibility also give positive recurrence of the original period-two chain.
\end{proof}

\begin{remark}
The argument proves polynomial drift of arbitrarily high order, using a possibly different sampling interval and finite set at each order. It does not prove geometric ergodicity. Geometric ergodicity is not needed for the following regenerative CLT and moment conclusions. Constants are not claimed uniform as a profile probability approaches zero or as $p$ approaches one.
\end{remark}

We now use returns to simultaneous balance to derive the joint limits.
Let $h(u)$ be the $J$-dimensional profile indicator. For $k$ allocation copies, index treatment paths by $m=1,\ldots,k$. For any nonempty $B\subseteq\{1,\ldots,k\}$ define
\[
 S_{B,n}=\sum_{i=1}^n h(U_i)\prod_{m\in B}Z_i^{(m)},\qquad
 S_{\varnothing,n}=\sum_{i=1}^n\{h(U_i)-\bm\pi\}.
\]
Singleton subsets give within-stratum imbalances; pairs give the unnormalised treatment-sign product processes; triples give the unnormalised triple-product processes. Stack these vectors in $S_n$.

\begin{theorem}\label{thm:clt}\label{lem:supp-label-flip}
Under the conditions of Theorem~\ref{thm:recurrence}, for every fixed $k$,
\begin{equation}\label{eq:joint-clt}
 n^{-1/2}S_n\ \Rightarrow\ \mathcal N(0,\Omega_k),\qquad
 \Omega_k=\frac{\E[WW\trans]}{\E\tau_k},
\end{equation}
where $W$ is the total stacked reward in one all-zero-to-all-zero cycle. Moreover,
\begin{equation}\label{eq:fourth}
 \sup_{n\ge1}\E\norm{n^{-1/2}S_n}^4<\infty,
 \qquad n^{-1}\E[S_nS_n\trans]\longrightarrow\Omega_k.
\end{equation}
Every reward block has mean zero. Cross-covariance blocks for distinct subsets, including the centred profile block, are zero, and covariance blocks for subsets of equal cardinality agree.
\end{theorem}
\begin{proof}
Let $0=T_0<T_1<T_2<\cdots$ be successive visits of the joint state to the all-zero state. The cycle lengths $\tau_j=T_j-T_{j-1}$ and marked cycle rewards $W_j=S_{T_j}-S_{T_{j-1}}$ are iid pairs, by the Markov property and independence of future profile/randomisation variables. Theorem~\ref{thm:recurrence} gives all moments of $\tau_j$. Rewards per observation are bounded, so $\norm{W_j}\le C\tau_j$ and all cycle-reward moments are finite.

The profile reward has zero mean by optional stopping of the bounded-increment martingale $\sum_i\{h(U_i)-\bm\pi\}$ at an integrable cycle length. For a nonempty subset $B$, globally flip all treatment signs of any copy in $B$ within a cycle. The complete marked-cycle law and its endpoints are unchanged, whereas that block of the reward changes sign. Its mean is therefore zero. Flipping a copy in the symmetric difference of two distinct nonempty subsets makes their cross second moment zero. A flip in a nonempty subset similarly kills its cross moment with the centred profile block. Permuting allocation-sequence labels proves equality of same-cardinality covariance blocks.

For completeness, a bound for the incomplete final cycle is obtained before invoking the CLT. Put $N(n)=\max\{j:T_j\le n\}$ and $N_{\rm cyc}(n)=N(n)+1$. For integer renewal times the renewal mass
\[
 u(t)=\sum_{j\ge0}\Pp(T_j=t)
\]
satisfies $u(t)\le1$, because the events in the sum are disjoint. Independence of the next cycle gives, for every $a>0$,
\begin{align}\label{eq:cover-moment}
 \E[\tau_{N_{\rm cyc}(n)}^a]
 &=\sum_{t=0}^n u(t)\E[\tau_k^a\one\{\tau_k>n-t\}]\\
 &\le\sum_{l=0}^{\infty}\E[\tau_k^a\one\{\tau_k>l\}]
 =\E\tau_k^{a+1}<\infty.\nonumber
\end{align}
Thus the contribution of the incomplete final cycle is bounded in probability and is $o_p(\sqrt n)$.

The renewal strong law gives $N(n)/n\to1/\E\tau_k$. The iid finite-variance functional CLT for centred $W_j$, followed by this random time change, yields
\[
 n^{-1/2}\sum_{j=1}^{N(n)}W_j\Rightarrow
 \mathcal N\{0,\E[WW\trans]/\E\tau_k\}.
\]
The time change is valid even though $W_j$ and $\tau_j$ may be dependent: the time-change limit is deterministic, and the iid partial-sum process has the usual tightness. Adding the final incomplete cycle proves \eqref{eq:joint-clt}.

For the moment statement, the discrete-time martingale of full-cycle sums has iid centred increments with finite fourth moments. Doob's inequality and the fourth-moment expansion of an iid centred sum imply
\[
 \E\max_{0\le j\le n+1}\norm{\sum_{l=1}^jW_l}^4\le C_1(n+1)^2.
\]
Since $N_{\rm cyc}(n)\le n+1$ and
\[
 \norm{S_n-\sum_{j=1}^{N_{\rm cyc}(n)}W_j}\le C\tau_{N_{\rm cyc}(n)},
\]
\eqref{eq:cover-moment} with $a=4$ proves $\E\norm{S_n}^4\le C_2(n+1)^2$. This establishes the uniform normalised fourth-moment bound. It implies uniform integrability of each normalised second-product entry. Convergence in distribution then gives the second-moment limit in \eqref{eq:fourth}.

The same flip arguments apply at each deterministic $n$, so all the stated zero cross moments hold before taking limits as well. Finally, the profile block has exactly the iid multinomial covariance $\Pi-\bm\pi\bm\pi\trans$.
\end{proof}

\begin{theorem}\label{prop:supp-allocation-clt}
Under Assumptions~\ref{ass:supp-design} and \ref{ass:supp-profile}, index three allocation sequences by $0,1,2$. There are finite positive-semidefinite matrices $\Sigma_D,\Psi_D$ such that
\begin{align*}
 &\left(R_n,\frac{D_n^{(0)}}{\sqrt n},\frac{D_n^{(1)}}{\sqrt n},
 \frac{D_n^{(2)}}{\sqrt n},G_n^{(01)},G_n^{(02)},G_n^{(12)}\right)\\
 &\hspace{12mm}\Rightarrow
 \mathcal N\!\left(0,
 \diag(\Pi-\bm\pi\bm\pi\trans,
 \Sigma_D,\Sigma_D,\Sigma_D,\Psi_D,\Psi_D,\Psi_D)\right).
\end{align*}
The corresponding second moments converge, and $K_n^{(012)}$ itself has a centred Gaussian limit, possibly degenerate, and bounded fourth moments. In particular $K_n^{(012)}=O_p(1)$. The result for a single allocation sequence follows by marginalisation. Here $\Sigma_S=\Pi-\bm\pi\bm\pi^\top$. All covariance limits are independent of the number of allocation sequences used to represent them, since their finite-$n$ marginal laws are identical. Consequently Assumption~\ref{ass:supp-general-allocation} holds for the Pocock--Simon range method under Assumptions~\ref{ass:supp-design} and \ref{ass:supp-profile}.
\end{theorem}
\begin{proof}
Apply Theorem~\ref{thm:clt} with $k=3$ and retain the profile, singleton, pair and triple blocks. The covariance formula and label symmetries give the displayed blocks; the fourth-moment bound implies the required second-moment convergence. The marginal distribution of a selected set of sequences is unchanged by generating additional sequences, so the covariance limits do not depend on that auxiliary number.
\end{proof}

\Needspace{9\baselineskip}
The same cycle representation determines the nullspace of $\Sigma_D$ and the positive definiteness of $\Psi_D$.
\begin{proposition}\label{prop:rank}
Under Assumptions~\ref{ass:supp-design} and \ref{ass:supp-profile},
\[
 \ker\Sigma_D=\operatorname{range}(\mathsf L\trans),\qquad
 \operatorname{rank}(\Sigma_D)=2^F-(F+1),\qquad \Psi_D\text{ is positive definite}.
\]
\end{proposition}
\begin{proof}
At the end of any one-sequence regeneration cycle, $\mathsf L W_{\{1\}}=0$, so additive profile scores lie in the kernel of the regenerative covariance. Conversely, if $v\notin\operatorname{range}(\mathsf L\trans)$, there is an integer vector $z\in\ker \mathsf L$ with $v\trans z\ne0$. Indeed, the nullspace of the integer matrix $\mathsf L$ has a rational basis; clear denominators in a basis vector having nonzero inner product with $v$.

Realise $z$ by $|z_u|$ arrivals of profile $u$ with treatment sign $\sg(z_u)$. This finite marked path has positive probability and ends at zero marginal state because $\mathsf L z=0$. Decompose it at its intermediate returns to zero. Since its total $v$-reward is nonzero, some first-return segment has nonzero $v$-reward and positive probability. The mean cycle reward is zero by symmetry, so $\E(v\trans W_{\{1\}})^2>0$. This proves the exact kernel and rank.

For the product-process block use any fixed $k\ge2$. A two-step cycle consists of the same profile $u$ twice and opposite signs on the second observation in every copy. Given the first signs, all second signs are preferred by the exact range rule, so the total probability of these two-step cycles is $\pi_u^2p^k$. The reward for a fixed pair is $\pm2h(u)$. Therefore, for every $v\in\R^J$,
\[
 v\trans\Psi_Dv=\frac{\E(v\trans W_{\{1,2\}})^2}{\E\tau_k}
 \ge\frac{4p^k}{\E\tau_k}\sum_u\pi_u^2v_u^2.
\]
This is strictly positive for $v\ne0$.
\end{proof}


\subsection{Covariance estimation from simulated allocation sequences}
For randomisation-design simulation $q$ (profiles and treatment assignments generated without outcomes), define
\[
U_s^{(q)}=
\begin{cases}
D_s^{(q)}/\sqrt{N_s^{(q)}},&N_s^{(q)}>0,\\
0,&N_s^{(q)}=0,
\end{cases}
\]
and let $\widehat{\bm\Gamma}_n$ be the sample covariance of $B_0$ independent vectors $\bm U^{(q)}$. For a realised $\bm N$, define
\begin{equation}\label{eq:supp-Omega}
\widetilde{\bm\Omega}_n(\bm N)
=\bm N^{1/2}\widehat{\bm\Gamma}_n\bm N^{1/2}.
\end{equation}

For a second set of design simulations, generate two allocation sequences independently conditional on a common profile sequence and define
\begin{equation}\label{eq:supp-G-calibration}
\bm G_n^{(q,12)}
=\frac1{\sqrt n}\sum_{i=1}^n
\bm h(S_i^{(q)})Z_i^{(q,1)}Z_i^{(q,2)}.
\end{equation}
Let $\widehat{\bm\Psi}_n$ be the sample covariance of $B_\Psi$ independent vectors $\bm G_n^{(q,12)}$. For the outcome trial, let
\begin{equation}\label{eq:supp-Pi-hat}
\widehat{\bm\Pi}_n=\bm N/n.
\end{equation}
This defines an estimator based on two allocation sequences. The reported implementation using three sequences instead averages the two sample covariance matrices for $\bm G_n^{(01)}$ and $\bm G_n^{(02)}$; it does not take the covariance of their average. Unlike the covariance estimate based on a single allocation sequence in \eqref{eq:supp-Omega}, no rescaling by the realised counts is needed for first-order consistency of the normalised covariance of the treatment-sign product process.


\begin{assumption}\label{ass:supp-calibration}
Randomisation-design simulations use the trial's profile law and the same allocation rule, are independent across simulation replicates, and use random numbers independent of the outcome trial and regenerated reference-test sequences. The replication counts satisfy $B_0(n)\to\infty$ and, when $\Psi_D$ is estimated, $B_\Psi(n)\to\infty$. Alternatively the profile law is consistently estimated and satisfies Proposition~\ref{prop:alloc-plugin}. Covariance regularisation, if used, has operator-norm difference $o_p(1)$ from the raw sample covariance.
\end{assumption}
The final regularisation condition is achieved by symmetrisation and projection onto the positive-semidefinite cone. For numerical removal of small positive eigenvalues, use the relative threshold $t_n=\min(10^{-12},n^{-1})$. Its induced operator-norm error is at most $t_n\max(1,\|\widehat\Sigma\|)$ and is $o_p(1)$, since the raw matrices are $O_p(1)$. This threshold equals $10^{-12}$ at every sample size reported in the numerical studies. The released code implements this threshold.

\begin{corollary}\label{cor:calibration}\label{prop:supp-calibration}
Under Assumption~\ref{ass:supp-general-allocation},
\[
 \sup_n\E\norm{\bm U_n}^4<\infty,\qquad
 \sup_n\E\norm{G_n^{(12)}}^4<\infty.
\]
Under Assumption~\ref{ass:supp-calibration}, the sample covariance estimates are therefore consistent when $B_0(n)\to\infty$ and $B_\Psi(n)\to\infty$:
\[
 \widehat\Gamma_n\xrightarrow{p}\Pi^{-1/2}\Sigma_D\Pi^{-1/2},\qquad
 \widehat\Psi_n\xrightarrow{p}\Psi_D.
\]
The bound for $\bm U_n$ follows below from the allocation fourth-moment bound in Assumption~\ref{ass:supp-general-allocation}; for the range method all these bounds are consequences of Theorem~\ref{prop:supp-allocation-clt}.
\end{corollary}
\begin{proof}
The fourth-moment bound for $G_n^{(12)}$ is part of Assumption~\ref{ass:supp-general-allocation}. For a fixed profile $u$, split according as $N_{n,u}\ge n\pi_u/2$. On this event,
\[
 |U_{n,u}|^4\le\frac{4|D_{n,u}|^4}{n^2\pi_u^2},
\]
whose expectation is uniformly bounded by Assumption~\ref{ass:supp-general-allocation}. On the complementary event, $|D_{n,u}|\le N_{n,u}$ gives $|U_{n,u}|^4\le n^2$ (with the prescribed zero-count convention). Since $N_{n,u}$ is binomial with parameter $\pi_u$, an exponential lower-tail bound gives
\[
 \E[|U_{n,u}|^4\one\{N_{n,u}<n\pi_u/2\}]
 \le n^2\exp(-n\pi_u/8),
\]
which is bounded uniformly in $n$. There are finitely many profiles.

By the count law of large numbers and the allocation CLT,
$\bm U_n\Rightarrow\mathcal N(0,\Pi^{-1/2}\Sigma_D\Pi^{-1/2})$.
The just-proved fourth-moment bound implies convergence of second moments. Means are zero by label symmetry conditional on profiles. For independent design-simulation replicates, each uncentred sample second moment has variance $O(B^{-1})$ uniformly in $n$; the empirical-mean subtraction vanishes as well. This proves both covariance consistency statements. For realised counts, $\widetilde\Omega_n/n=(N/n)^{1/2}\widehat\Gamma_n(N/n)^{1/2}\to\Sigma_D$. Also $\operatorname{tr}(N^+\widetilde\Omega_n)=O_p(1)$, while $J$ is fixed and every stratum is nonempty with probability tending to one. Thus $\widehat\kappa_n^{\rm raw}\to1$, and its positive-part safeguard is asymptotically inactive. The average of the two consistent product-process covariance estimators in the three-sequence implementation is consistent even though the product processes within a simulation replicate are dependent; the replicates themselves are independent.
\end{proof}


We next allow the profile distribution used in the simulations to be estimated from the trial.
\begin{proposition}\label{prop:alloc-plugin}
Fix $F$ and $p_{\rm bc}$. Let $\mathcal P_0$ be a compact full-support set of profile laws on which the drift bounds \eqref{eq:H1}--\eqref{eq:H2} hold with common positive constants. Then the return-time moments and normalised fourth-moment bounds above are uniform on $\mathcal P_0$. The covariance functions $\Sigma_D(\pi)$, $\Psi_D(\pi)$ and $\Gamma(\pi)=\Pi^{-1/2}\Sigma_D(\pi)\Pi^{-1/2}$ are continuous there, and the finite-$n$ covariance matrices of $D_n/\sqrt n$, $G_n^{(12)}$ and $\bm U_n$ converge uniformly to these limits.

Suppose the true law $\pi$ has a neighbourhood of this kind and $\widehat\pi_n\to\pi$ in probability, with all coordinates of $\widehat\pi_n$ positive. Randomisation-design covariance estimation under $\widehat\pi_n$, using random numbers independent of the outcome trial and regenerated test paths conditional on $\widehat\pi_n$, remains consistent when $B_0(n),B_\Psi(n)\to\infty$. This includes smoothed empirical frequencies $(N_s+a)/(n+Ja)$ for any fixed $a>0$. Every positive two-factor law, the five-factor laws in Corollary~\ref{cor:open}, and the prevalence boxes containing the distributions used in the simulation code of Proposition~\ref{prop:alloc-box} satisfy the neighbourhood requirement.
\end{proposition}
\begin{proof}
On $\mathcal P_0$, $K$ has eigenvalues uniformly bounded away from zero and infinity. In the proof of Lemma~\ref{lem:polynomial}, the increment bounds, $C'_r$, $T_r$ and the final finite sublevel set can therefore be chosen uniformly. The same is true for any fixed number of copies. All profile and treatment-sign probabilities along a fixed admissible path have a common positive lower bound. Lemma~\ref{lem:alloc-return} now gives uniform return-time moments of every fixed order.

For a fixed cycle-length cutoff, cycle probabilities are continuous functions of $\pi$: only finitely many marked paths occur, and the treatment probabilities along each path do not depend on $\pi$. A cycle reward has norm at most a constant times its length. Uniform higher return moments permit removal of the cutoff uniformly. Thus $E_\pi\tau_k$ and $E_\pi[WW^\top]$ are continuous, and so is their ratio.

Here is a uniform finite-$n$ argument, which does not assume interchange of limits. Let $N_{\rm cyc}(n)=\min\{j:T_j>n\}$ and $Y_n=\sum_{j=1}^{N_{\rm cyc}(n)}W_j$. Since $\{N_{\rm cyc}(n)\ge j\}=\{T_{j-1}\le n\}$ is measurable before cycle $j$, the stopped martingale second-moment identity gives
\[
 E_\pi[Y_nY_n^\top]=E_\pi N_{\rm cyc}(n)\,E_\pi[WW^\top].
\]
Wald's identity for cycle lengths and \eqref{eq:cover-moment} imply $E_\pi N_{\rm cyc}(n)=n/E_\pi\tau_k+O(1)$ uniformly. Also $\|S_n-Y_n\|\le C\tau_{N_{\rm cyc}(n)}$ has uniformly bounded second moment. Cauchy--Schwarz consequently gives
\[
 \sup_{\pi\in\mathcal P_0}
 \left\|n^{-1}E_\pi[S_nS_n^\top]-\Omega_k(\pi)\right\|=O(n^{-1/2}).
\]
The fourth-moment proof following \eqref{eq:cover-moment} is uniform for the same reason.

For count normalisation, on $\min_s N_s/n\ge\frac12\min_s\pi_s$, the map $t\mapsto t^{-1/2}$ is uniformly Lipschitz. Multinomial fourth moments and the just-established allocation fourth moments give
\[
 \sup_{\pi\in\mathcal P_0}E_\pi
 \left\|\bm U_n-\Pi^{-1/2}D_n/\sqrt n\right\|^2=O(n^{-1}).
\]
To include the complementary event, use $|D_{n,s}|\le N_s$, a common positive lower bound for $\pi_s$, and the exponential binomial lower-tail bound. Cauchy--Schwarz proves the desired uniform covariance convergence for $\bm U_n$.

Conditional on the estimated law lying in a compact neighbourhood $\mathcal P_0$, sample second moments from independent design-simulation replicates have mean-square error $O(B^{-1})$ uniformly. Uniform finite-$n$ convergence, continuity of the covariance functions and $\widehat\pi_n\to\pi$ prove consistency by the triangle inequality. The probability of leaving $\mathcal P_0$ tends to zero. For two factors, the analytic drift constants are common and compactness handles $K$; for the specified five-factor laws, Corollary~\ref{cor:open} and continuity provide the required neighbourhood.
\end{proof}


The covariance reconstruction in the main article can be justified within the same allocation theory. Let $\mathsf L$ have full row rank and satisfy $\Sigma_D\mathsf L^\top=0$, as proved for the overall and marginal contrasts of the range method. Compute $\widehat\Xi_n$ as the sample covariance of $X_n=\mathsf L D_n$ using the first allocation sequence in the same replicates that estimate $\widehat\Gamma_n$. For the realised counts set
\[
 B_N=N^{1/2}\mathsf L^\top,\qquad
 K_N=B_N^\top B_N=\mathsf L N\mathsf L^\top.
\]
On the event that $K_N$ is invertible, set
\[
 Q_N=B_NK_N^{-1},\qquad P_N=B_NK_N^{-1}B_N^\top,
\]
\[
 \widehat\Gamma_n^{\,*}
 =(I-P_N)\widehat\Gamma_n(I-P_N)+Q_N\widehat\Xi_nQ_N^\top,
 \qquad
 \widetilde\Omega_n^{\,*}=N^{1/2}\widehat\Gamma_n^{\,*}N^{1/2}.
\]
If $K_N$ is singular, set $\widehat\Gamma_n^{\,*}=\widehat\Gamma_n$ and retain the original $\widetilde\Omega_n$ for the whole covariance, rather than using an inverse on only some contrasts. Empty joint strata do not by themselves require this fallback: the represented profiles may still give full column rank to $B_N$. Fixed positive profile probabilities imply that the full-rank event has probability tending to one.

Both terms in $\widehat\Gamma_n^{\,*}$ are positive semidefinite. Since $B_N^\top(I-P_N)=0$ and $B_N^\top Q_N=I$, the full-rank event gives the exact identity
\[
 \mathsf L\widetilde\Omega_n^{\,*}\mathsf L^\top
 =B_N^\top\widehat\Gamma_n^{\,*}B_N=\widehat\Xi_n.
\]
The reconstruction is invariant to replacing $\mathsf L$ by $T\mathsf L$ for an invertible matrix $T$, provided $\widehat\Xi_n$ is replaced by $T\widehat\Xi_nT^\top$. Indeed, $B_N$ then becomes $B_NT^\top$, the projector is unchanged, and $Q_N$ becomes $Q_NT^{-1}$. This also leaves $Q_N\widehat\Xi_nQ_N^\top$ unchanged. The cross-covariance between the range of $B_N$ and its orthogonal complement is set to zero. It is not asserted to vanish in the true finite-sample conditional covariance. The matrix $\widehat\Xi_n$ estimates a finite-$n$ covariance averaged over the simulated profile law, not a covariance conditional on the ordered profiles of the outcome trial.

For the first-order limit, label symmetry gives $\E X_n=0$. Assumption~\ref{ass:supp-general-allocation} and its second-moment convergence imply
\[
 \frac1n\E\|X_n\|^2
 =\operatorname{tr}\left\{\mathsf L\frac{\E(D_nD_n^\top)}n\mathsf L^\top\right\}
 \longrightarrow\operatorname{tr}(\mathsf L\Sigma_D\mathsf L^\top)=0.
\]
For $B_0\ge2$, the sample covariance obeys
\[
 0\le\operatorname{tr}(\widehat\Xi_n/n)
 \le\frac1{n(B_0-1)}\sum_{q=1}^{B_0}\|X_n^{(q)}\|^2.
\]
Taking expectations and applying Markov's inequality proves $\widehat\Xi_n/n\to_p0$. Only normalised second-moment convergence is used here. For a consistently estimated profile law, the uniform second-moment convergence in Proposition~\ref{prop:alloc-plugin} gives the same conclusion on the compact neighbourhood used there.

Write $\overline B_n=B_N/\sqrt n$ and $\overline B=\Pi^{1/2}\mathsf L^\top$. Then $\overline B_n\to_p\overline B$, $\overline B^\top\overline B$ is positive definite, and
\[
 P_N\to_p P_0=\overline B(\overline B^\top\overline B)^{-1}\overline B^\top.
\]
By Corollary~\ref{prop:supp-calibration},
$\widehat\Gamma_n\to_p\Gamma=\Pi^{-1/2}\Sigma_D\Pi^{-1/2}$.
The balance condition gives $\Gamma\overline B=0$, so
$(I-P_0)\Gamma(I-P_0)=\Gamma$. Moreover,
\[
 Q_N\widehat\Xi_nQ_N^\top
 =\overline B_n(\overline B_n^\top\overline B_n)^{-1}(\widehat\Xi_n/n)
       (\overline B_n^\top\overline B_n)^{-1}\overline B_n^\top\to_p0.
\]
The singular event has probability tending to zero and uses the original consistent estimator. Consequently,
\[
 \widehat\Gamma_n^{\,*}-\widehat\Gamma_n\to_p0,
 \qquad \widetilde\Omega_n^{\,*}/n\to_p\Sigma_D.
\]
When implementing the rank check in floating-point arithmetic, the code compares the smallest eigenvalue of $K_N$ to
$\min(10^{-12},n^{-1})\max\{1,\lambda_{\max}(K_N)\}$.
Under the preceding conditions this check is eventually inactive in probability. Its fallback is recorded in every simulation.

Recompute $\widehat\kappa_n^{\,*}$ from \eqref{eq:supp-kappa} using $\widetilde\Omega_n^{\,*}$, with the same positive-part convention. If $E_N=\diag\{\mathbf1(N_s>0)\}$, then
\[
 \operatorname{tr}(N^+\widetilde\Omega_n^{\,*})
 =\operatorname{tr}(E_N\widehat\Gamma_n^{\,*})=O_p(1),
 \qquad \widehat\kappa_n^{\,*}\to_p1.
\]
The score-moment result in Lemma~\ref{lem:supp-score-moments} therefore yields, for any fixed tested value $b$,
\[
 \frac{\widehat V_{{\rm S},n}^{\,*}-\widehat V_{{\rm S},n}}n
 =\frac14\bar{\bm r}_n^\top
       \frac{\widetilde\Omega_n^{\,*}-\widetilde\Omega_n}{n}\bar{\bm r}_n
   +\frac14(\widehat\kappa_n^{\,*}-\widehat\kappa_n)
       \frac{\bm e_n^\top\bm e_n}{n}\to_p0.
\]
Using the same estimated effect deviations and the same $\widehat\Psi_n-\widehat\Pi_n$ correction gives the identical difference for the raw reference-variance estimators. When $v_{\rm R}(b)>0$, both variance floors are inactive with probability tending to one. Thus the sampling, conditional-reference and corrected-test results also hold with the reconstructed variances. The local-alternative conclusions follow from the same argument with the row-wise score-moment bounds in Assumption~\ref{ass:supp-local}. This proof uses the balance condition $\Sigma_D\mathsf L^\top=0$; applying the projector to arbitrary covariate directions of an unrestricted design is not justified by it.

The calculation requires only an $(F+1)$-dimensional inverse. On the full-rank event, with $w=N^{1/2}\bar{\bm r}$, the stratum quadratic is
\[
 \bar{\bm r}^\top\widetilde\Omega_n^{\,*}\bar{\bm r}
 =\{(I-P_N)w\}^\top\widehat\Gamma_n\{(I-P_N)w\}
  +(Q_N^\top w)^\top\widehat\Xi_n(Q_N^\top w).
\]
Because $E_NB_N=B_N$, its trace contribution is
\[
 \operatorname{tr}(N^+\widetilde\Omega_n^{\,*})
 =\operatorname{tr}(E_N\widehat\Gamma_n)
  -\operatorname{tr}(K_N^{-1}B_N^\top\widehat\Gamma_nB_N)
  +\operatorname{tr}(\widehat\Xi_nK_N^{-1}).
\]
These identities permit computation without forming a new $J$-dimensional matrix for every score. They are algebraic identities, not claims of a measured speed improvement. The earlier numerical studies use the original covariance construction. Appendix~K reports the separate paired evaluation of both constructions.

\section{Repeated-sampling inference at an average-effect null value}\label{sec:supp-sampling}

\subsection{Participant-level assumptions and score expansion}
\begin{assumption}\label{ass:supp-participants}
The vectors $\{S_i,\bm c_i,Y_i(0),Y_i(1)\}$ are independent and identically distributed across participants. The fixed-dimensional baseline vector $\bm c_i$ contains an intercept and may include baseline functions not used by the allocation rule. The matrix $\bm\Sigma_C=\E[\bm c_i\bm c_i^\top]$ is positive definite. For some $\eta>0$, $\E[\|\bm c_i\|^{4+\eta}+|Y_i(0)|^{4+\eta}+|Y_i(1)|^{4+\eta}]<\infty$. Conditional on the complete profile sequence, the observed and regenerated allocation random numbers are jointly independent of all remaining participant variables $\{\bm c_i,Y_i(0),Y_i(1)\}_{i\ge1}$.
\end{assumption}

Fix a null value $b$ and write
\begin{equation}\label{eq:supp-Qdelta}
Q_i(b)=\frac{Y_i(1)+Y_i(0)-b}{2},
\qquad
\delta_i(b)=Y_i(1)-Y_i(0)-b,
\end{equation}
so that
\begin{equation}\label{eq:supp-Widentity-general}
W_i(b)=Q_i(b)+\frac12Z_i^{(0)}\delta_i(b).
\end{equation}
Let $C$ have rows $\bm c_i^\top$ and put $\bm r(b)=M_C\bm W(b)$, where $M_C=I-C(C^\top C)^+C^\top$. Define the sampling variance estimate by $\widehat V_{{\rm S},n}(\bm r)=\{\bar{\bm r}^\top\widetilde\Omega_n\bar{\bm r}+\widehat\kappa_n\bm e^\top\bm e\}/4$, where $\widehat\kappa_n=\max(\kappa^{\rm raw},0)$ in \eqref{eq:supp-kappa} with $\Omega=\widetilde\Omega_n$, and use the stated zero-residual convention when $n=J_+$. On positive estimated-variance events write $Z_{{\rm S},n}(b)=T_n\{\bm r(b)\}/\sqrt{\widehat V_{{\rm S},n}\{\bm r(b)\}}$. For the adjusted score, define $\bm\beta_b=\bm\Sigma_C^{-1}\E[\bm c_iQ_i(b)]$, $q_i(b)=Q_i(b)-\bm c_i^\top\bm\beta_b$. For the unadjusted score, take $\bm c_i=1$. Let
\[
m_s(b)=\E[q_i(b)\mid S_i=s],
\qquad
\nu_i^Q(b)=q_i(b)-m_{S_i}(b),
\]
\[
d_s(b)=\E[\delta_i(b)\mid S_i=s],
\qquad
\nu_i^\delta(b)=\delta_i(b)-d_{S_i}(b).
\]
Write $\bm m(b)=\{m_1(b),\ldots,m_J(b)\}^\top$ and $\bm d(b)=\{d_1(b),\ldots,d_J(b)\}^\top$. At $\Delta=b$, $\E[\delta_i(b)]=\bm\pi^\top\bm d(b)=0$.

The potential-outcome notation is convenient for the algebra but does
not require the within-participant joint distribution of $Y_i(0)$ and
$Y_i(1)$ to be specified in the binary simulations. Define
$U_i(a;b)=Y_i(a)-\bm c_i^\top\bm\beta_b$. Conditional on $S_i$,
parallelogram identities give
\begin{align}
\E[\nu_i^Q(b)^2]
+\frac14\E[\Var[\delta_i(b)\mid S_i]]
&=
\frac12\E\left[
\Var[U_i(1;b)\mid S_i]
+\Var[U_i(0;b)\mid S_i]
\right],
\label{eq:supp-coupling-R}\\
\E[\nu_i^Q(b)^2]
+\frac14\E[\delta_i(b)^2]
&=
\frac12\E\left[
\Var[U_i(1;b)\mid S_i]
+\Var[U_i(0;b)\mid S_i]
\right]
\nonumber\\
&\quad+\frac14\bm d(b)^\top\bm\Pi\bm d(b).
\label{eq:supp-coupling-S}
\end{align}
The stratum means $\bm m(b)$ and $\bm d(b)$ likewise involve only
conditional marginal means. Hence the variance limits used below are
invariant to the unidentified conditional association between the two
potential outcomes.

Let $\widehat{\bm\beta}_b=(\bm C^\top\bm C)^{-1}\bm C^\top\bm W(b)$ for the adjusted score.

\begin{lemma}\label{lem:supp-beta-general}
Under Assumptions~\ref{ass:supp-general-allocation} and \ref{ass:supp-participants},
\[
\widehat{\bm\beta}_b-\bm\beta_b=O_p(n^{-1/2}),
\qquad
\frac1{\sqrt n}\sum_{i=1}^nZ_i^{(0)}\bm c_i=O_p(1).
\]
If the same allocation law is used for regenerated paths, the
second bound holds jointly for $m=0,1,2$, by Assumption~\ref{ass:supp-general-allocation}.
Consequently, for each fixed allocation sequence,
\begin{equation}\label{eq:supp-beta-rem}
\frac1{\sqrt n}
(\bm Z^{(m)})^\top\bm C
(\widehat{\bm\beta}_b-\bm\beta_b)
=o_p(1).
\end{equation}
\end{lemma}

\begin{proof}
From \eqref{eq:supp-Widentity-general},
\[
\frac1n\bm C^\top\bm W(b)
=\frac1n\sum_i\bm c_iQ_i(b)
+\frac1{2n}\sum_i\bm c_iZ_i^{(0)}\delta_i(b).
\]
For $\bm g_i=\bm c_i\delta_i(b)$, write
$\bm g_i=\E[\bm g_i\mid S_i]+\{\bm g_i-\E[\bm g_i\mid S_i]\}$. The stratum-mean part of the weighted sum is a fixed linear combination of $\bm D_n^{(0)}$ and is $O_p(\sqrt n)$. Conditional on profiles and assignments, the residual part is a sum of independent mean-zero vectors with finite $2+\eta/2$ moments by H\"older's inequality and is also $O_p(\sqrt n)$. Thus the second term is $O_p(n^{-1/2})$. The ordinary law of large numbers and central limit theorem give the result for $\widehat{\bm\beta}_b$. The same decomposition with $\bm c_i$ in place of $\bm g_i$ proves the bound for $n^{-1/2}\sum_iZ_i^{(0)}\bm c_i$. Equation \eqref{eq:supp-beta-rem} follows by multiplication of $O_p(1)$ and $O_p(n^{-1/2})$ terms.
\end{proof}

\begin{proposition}\label{prop:supp-observed-expansion}
Fix $b$, not necessarily equal to $\Delta$. Under Assumptions~\ref{ass:supp-general-allocation} and \ref{ass:supp-participants},
\begin{equation}\label{eq:supp-observed-expansion}
\frac{T_n\{\bm r(b)\}}{\sqrt n}
=
\frac12\bm m(b)^\top\frac{\bm D_n^{(0)}}{\sqrt n}
+\frac1{2\sqrt n}\sum_{i=1}^nZ_i^{(0)}\nu_i^Q(b)
+\frac1{4\sqrt n}\sum_{i=1}^n\delta_i(b)
+o_p(1).
\end{equation}
\end{proposition}

\begin{proof}
Using $(Z_i^{(0)})^2=1$ and \eqref{eq:supp-Widentity-general},
\[
T_n\{\bm r(b)\}
=\frac12\sum_iZ_i^{(0)}q_i(b)
+\frac14\sum_i\delta_i(b)
-\frac12(\bm Z^{(0)})^\top\bm C
(\widehat{\bm\beta}_b-\bm\beta_b).
\]
The last term is $o_p(\sqrt n)$ by Lemma~\ref{lem:supp-beta-general}. Decompose $q_i(b)=m_{S_i}(b)+\nu_i^Q(b)$ and use
\[
\sum_iZ_i^{(0)}m_{S_i}(b)=\bm m(b)^\top\bm D_n^{(0)}.
\]
\end{proof}

\subsection{Joint central limit theorem and consistency of SIGA-S}
\begin{lemma}\label{lem:supp-participant-transform}
Under Assumptions~\ref{ass:supp-general-allocation} and \ref{ass:supp-participants},
\[
\left(
\frac{\bm D_n^{(0)}}{\sqrt n},
\frac1{\sqrt n}\sum_iZ_i^{(0)}\nu_i^Q(b),
\frac1{\sqrt n}\sum_i\delta_i(b)
\right)
\xrightarrow{d}
\mathcal N\left[
\bm0,
\begin{pmatrix}
\bm\Sigma_D&\bm0&\bm0\\
\bm0^\top&\E[\nu_i^Q(b)^2]&0\\
\bm0^\top&0&\E[\delta_i(b)^2]
\end{pmatrix}
\right]
\]
at $\Delta=b$.
\end{lemma}

\begin{proof}
Write $\delta_i(b)=d_{S_i}(b)+\nu_i^\delta(b)$. The stratum-mean contribution to the third component is $\bm d(b)^\top\bm R_n$ because $\bm\pi^\top\bm d(b)=0$ at the null value. Assumption~\ref{ass:supp-general-allocation} gives its joint Gaussian limit and zero covariance with the imbalance block. Conditional on the profiles and observed allocation path, the remaining participant-level terms are independent across $i$, and their conditional variances converge by the law of large numbers. The conditional covariance between the second component and the centred third component is
\[
\frac1n\sum_iZ_i^{(0)}
\E[\nu_i^Q(b)\delta_i(b)\mid S_i]
=O_p(n^{-1/2}),
\]
because it is a fixed linear combination of the components of $\bm D_n^{(0)}/n$. The finite $2+\eta$ moments imply the conditional Lindeberg condition. The variance of the third component is the sum of $\E[\Var[\delta_i(b)\mid S_i]]$ and $\bm d(b)^\top\bm\Sigma_S\bm d(b)=\bm d(b)^\top\bm\Pi\bm d(b)$. Conditional characteristic-function convergence, combined with Assumption~\ref{ass:supp-general-allocation}, gives the joint result. Cross-covariances with the imbalance block vanish because the participant-level transforms have conditional mean zero.
\end{proof}

Let $\bar{\bm r}_n=\bm N^+\bm H^\top\bm r(b)$ and $\bm e_n=\bm M_H\bm r(b)$.

\begin{lemma}\label{lem:supp-score-moments}
Fix $b$, not necessarily equal to the data-generating marginal effect. Under the same assumptions,
\[
\bar{\bm r}_n\xrightarrow{p}\bm m(b),
\]
and
\begin{equation}\label{eq:supp-e2-limit}
\frac{\bm e_n^\top\bm e_n}{n}
\xrightarrow{p}
\E[\nu_i^Q(b)^2]
+\frac14\E[\delta_i(b)^2].
\end{equation}
Moreover,
\[
\frac{\max_i e_{ni}^2}{\bm e_n^\top\bm e_n}
\xrightarrow{p}0
\]
whenever the limiting denominator in \eqref{eq:supp-e2-limit} is positive.
\end{lemma}

\begin{proof}
By Lemma~\ref{lem:supp-beta-general}, replacing $\widehat{\bm\beta}_b$ by $\bm\beta_b$ changes the empirical score norm by $o_p(n)$ and its stratum means by $o_p(1)$. Thus it suffices to consider
$a_i=q_i(b)+Z_i^{(0)}\delta_i(b)/2$. Within each stratum, the average of $q_i(b)$ converges to $m_s(b)$, while
$n_s^{-1}\sum_{S_i=s}Z_i^{(0)}\delta_i(b)=O_p(n_s^{-1/2})$. This bound does not require $\Delta=b$. It proves the stratum-mean limit. After empirical stratum demeaning, expansion of the square gives
\[
 \frac{\bm e_n^\top\bm e_n}{n}
 =\frac1n\sum_i\left\{(\nu_i^Q(b))^2+\frac14\delta_i(b)^2
       +Z_i^{(0)}\nu_i^Q(b)\delta_i(b)\right\}+o_p(1).
\]
The first two averages converge to their expectations. The cross-term average is $o_p(1)$ by the same weighted-sum argument used in Lemma~\ref{lem:supp-participant-transform}; that argument does not use $\Delta=b$. The maximal-term result follows from the finite $4+\eta$ moment and positive profile probabilities.
\end{proof}

For a fixed null value $b$, define the score-quadratic limit
\begin{equation}\label{eq:supp-vS-tilde}
\widetilde v_{\rm S}(b)
=
\frac14\left[
\bm m(b)^\top\bm\Sigma_D\bm m(b)
+\E[\nu_i^Q(b)^2]
+\frac14\E[\delta_i(b)^2]
\right].
\end{equation}

\begin{proposition}\label{prop:supp-fixed-b-quadratic}
Under Assumptions~\ref{ass:supp-general-allocation}, \ref{ass:supp-participants} and \ref{ass:supp-calibration}, for any fixed $b$,
\[
\frac{\widehat V_{{\rm S},n}\{\bm r(b)\}}{n}
\xrightarrow{p}
\widetilde v_{\rm S}(b).
\]
At $\Delta=b$, $\widetilde v_{\rm S}(b)=v_{\rm S}(b)$ and is the repeated-sampling variance limit of the observed statistic. Away from the null value, \eqref{eq:supp-vS-tilde} is the probability limit of the score quadratic used by the test; it is not asserted to equal the centred repeated-sampling variance under the fixed alternative.
\end{proposition}

\begin{proof}
The result follows from Corollary~\ref{prop:supp-calibration}, Lemma~\ref{lem:supp-score-moments} and the quadratic representation \eqref{eq:supp-quadratic}.
\end{proof}

\begin{theorem}\label{thm:supp-sampling}
Under Assumptions~\ref{ass:supp-general-allocation}, \ref{ass:supp-participants} and \ref{ass:supp-calibration}, at $\Delta=b$,
\begin{equation}\label{eq:supp-vS}
\frac{T_n\{\bm r(b)\}}{\sqrt n}
\xrightarrow{d}
\mathcal N\{0,v_{\rm S}(b)\},
\end{equation}
where
\[
v_{\rm S}(b)
=\frac14\left[
\bm m(b)^\top\bm\Sigma_D\bm m(b)
+\E[\nu_i^Q(b)^2]
+\frac14\E[\delta_i(b)^2]
\right].
\]
Furthermore,
\begin{equation}\label{eq:supp-VS-consistency}
\frac{\widehat V_{{\rm S},n}\{\bm r(b)\}}n
\xrightarrow{p} v_{\rm S}(b).
\end{equation}
If $v_{\rm S}(b)>0$, then
\[
Z_{{\rm S},n}(b)
\xrightarrow{d}
\mathcal N(0,1),
\]
and its one- or two-sided normal-approximation $p$-value is asymptotically valid at the average-effect null value.
\end{theorem}

\begin{proof}
Equation \eqref{eq:supp-vS} follows from Proposition~\ref{prop:supp-observed-expansion}, Lemma~\ref{lem:supp-participant-transform} and Slutsky's theorem. Proposition~\ref{prop:supp-fixed-b-quadratic}, evaluated at $\Delta=b$, gives \eqref{eq:supp-VS-consistency}. Studentisation completes the proof.
\end{proof}

For comparison with the unadjusted variance in \citet{BugniCanayShaikh2018}, take $\bm c_i=1$, set $\mu_{as}=\E\{Y_i(a)\mid S_i=s\}$ and $\sigma_{as}^2=\Var\{Y_i(a)\mid S_i=s\}$, and write $\mu_a=\E Y_i(a)$. At $\Delta=b$, $d_s=\mu_{1s}-\mu_{0s}-b$ and $m_s=(\mu_{1s}+\mu_{0s}-\mu_1-\mu_0)/2$. The identities in (\ref{eq:supp-coupling-S}) give
\begin{equation}\label{eq:supp-bugni-comparison}
 16v_{\rm S}
 =4\bm m^\top\Sigma_D\bm m
  +2\sum_s\pi_s(\sigma_{1s}^2+\sigma_{0s}^2)
  +\sum_s\pi_s d_s^2.
\end{equation}
The factor 16 converts the variance of $T_n/\sqrt n$ to that of the treatment-minus-control mean difference: by the exact identity in the main article, $4T_n/\sqrt n$ and $\sqrt n(\bar Y_1-\bar Y_0-b)$ differ by $o_p(1)$ at this null value.

With treatment target $1/2$, our imbalance $D_{n,s}=\sum_{S_i=s}(2A_i-1)$ is twice that in \citet{BugniCanayShaikh2018}. On their allocation domain, therefore, $\Sigma_D=4\operatorname{diag}\{\pi_s\tau(s)\}$. Substitution into (\ref{eq:supp-bugni-comparison}) gives their three terms in equations~(15)--(18): within-stratum outcome variation, stratum-specific effect variation, and allocation imbalance. This establishes agreement on the overlapping unadjusted setting; it does not transfer their diagonal-covariance assumption to minimisation. The treatment-sign product covariance $\Psi_D$, used for the conditional analysis below, is a different quantity.


\begin{corollary}\label{cor:supp-sharp-stratum}
If $\delta_i(b)=0$ almost surely and the sharp null $Y_i(1)-Y_i(0)=b$ is compatible with the outcome support, Theorem~\ref{thm:supp-sampling} reduces to the result for a constant individual treatment effect. If $d_s(b)=0$ for every $s$, individual treatment effects may remain heterogeneous, but the difference between the two variance contributions in Supplementary Appendix C is zero and the conditional randomisation and repeated-sampling variances agree to first order.
\end{corollary}

\begin{assumption}\label{ass:supp-local}
The allocation rule, its fixed full-support profile probabilities, and the number of profiles are unchanged across $n$. In row $n$, the participant vectors $(S_i,\bm c_{ni},Y_{ni}(0),Y_{ni}(1))$ are independent and identically distributed and independent of the allocation random numbers conditional on the profiles. For a fixed $\eta>0$, the $4+\eta$ moments in Assumption~\ref{ass:supp-participants} are bounded uniformly in $n$, and the smallest eigenvalue of $E(\bm c_{ni}\bm c_{ni}^\top)$ is bounded away from zero. Let $\bm\beta_{n,b}$ be the row-specific population projection and define $q_{ni},\nu^Q_{ni},\delta_{ni},\bm m_n,\bm d_n$ as in Appendix~B.1. The conditional means $\bm m_n,\bm d_n$ and the conditional second moments of $(\nu^Q_{ni},\nu^\delta_{ni})$ converge, in every stratum, to those of a law satisfying $\Delta=b$. Finally $\Delta_n=b+h/\sqrt n$ for a fixed real $h$.
\end{assumption}

\begin{corollary}\label{cor:supp-local-S}
Under Assumptions~\ref{ass:supp-general-allocation}, \ref{ass:supp-calibration} and \ref{ass:supp-local}, with $v_{\rm S}(b)>0$, the upper-tail level-$\alpha$ SIGA-S test satisfies
\[
 \Pbb\{Z_{{\rm S},n}(b)>z_{1-\alpha}\}
 \longrightarrow
 1-\Phi\{z_{1-\alpha}-h/(4\sqrt{v_{\rm S}(b)})\}.
\]
\end{corollary}
\begin{proof}
Uniform moment bounds and the uniform eigenvalue bound give the coefficient-estimation expansion in Lemma~\ref{lem:supp-beta-general} row by row. The allocation process is unchanged, so Assumption~\ref{ass:supp-general-allocation} applies without any triangular-array allocation assumption. Conditional on profiles and paths, the uniform $2+\epsilon$ moments give Lindeberg's condition for the participant residual transforms, by the bound $n^{-\epsilon/2}C$. Their conditional covariance matrices converge to the stated limits at the null value, since there are finitely many strata and the relevant conditional moments converge. In the third component of \eqref{eq:supp-observed-expansion}, write
\[
 n^{-1/2}\sum_i\delta_{ni}(b)
 =h+\bm d_n^\top\bm R_n+n^{-1/2}\sum_i\nu^\delta_{ni}(b).
\]
The same conditional characteristic-function argument as in Lemma~\ref{lem:supp-participant-transform} therefore gives the covariance at the null value, with mean shift $h/4$ in the observed statistic. Uniform moments also give the row-wise laws of large numbers for the score quadratic. Covariance consistency gives the denominator limit at the null value. Slutsky's theorem completes the proof. A participant or allocation central limit theorem is not assumed in Assumption~\ref{ass:supp-local}.
\end{proof}

\section{Conditional randomisation theory and covariance correction}\label{sec:supp-randomization}

\subsection{Conditional weak convergence}
Let $\mathcal A_n$ contain the randomisation-design covariance estimation variables. Conditional statements involving an estimated variance denominator condition on $\mathcal A_n$ as well as the stated trial sigma-field. Given the observed trial, the covariance-estimation random numbers and all regenerated reference-test paths are independent. Thus adding $\mathcal A_n$ does not alter the conditional law of a regenerated path. Covariance estimates can nevertheless depend on the observed counts through rescaling or the plug-in profile law.

Generate the observed path $\bm Z^{(0)}$ and two regenerated paths $\bm Z^{(1)}$ and $\bm Z^{(2)}$ independently conditional on the same profile sequence. Let $\mathcal F_n$ be the enlarged sigma-field generated by the profile sequence, baseline variables, both potential outcomes and the observed allocation path. Let $\mathcal O_n$ be the observed-data sigma-field generated by the profile sequence, baseline variables, observed outcomes and observed allocation path; the fixed score is measurable with respect to $\mathcal O_n$. Conditional convergence given $\mathcal F_n$ implies the corresponding conclusion given $\mathcal O_n$ by the tower property and contraction of conditional expectations.

The next lemma relates joint convergence of two conditionally independent replicates to convergence of their conditional law; see \citet{BuecherKojadinovic2019}. We include the argument for the finite-dimensional vectors used here.

\begin{lemma}\label{lem:supp-conditional-convergence}
Let $\bm X_n^{(1)}$ and $\bm X_n^{(2)}$ take values in a fixed finite-dimensional Euclidean space and be conditionally independent and identically distributed given $\mathcal F_n$. If
\[
(\bm X_n^{(1)},\bm X_n^{(2)})
\xrightarrow{d}
(\bm X^{(1)},\bm X^{(2)}),
\]
where the limits are independent with common law $\mathcal Q$, then
\[
d_{\mathrm{BL}}\{\mathcal L(\bm X_n^{(1)}\mid\mathcal F_n),\mathcal Q\}
\xrightarrow{p}0.
\]
\end{lemma}

\begin{proof}
For a bounded Lipschitz function $f$, let $\mathcal Q_nf=\E[f(\bm X_n^{(1)})\mid\mathcal F_n]$. By marginalisation of the joint weak convergence assumed in the lemma, $\bm X_n^{(1)}\xrightarrow{d}\bm X^{(1)}\sim\mathcal Q$. Since $f$ is bounded and continuous,
\[
\E[\mathcal Q_nf]
=
\E[f(\bm X_n^{(1)})]
\longrightarrow
\mathcal Qf.
\]
Because $\bm X_n^{(1)}$ and $\bm X_n^{(2)}$ are conditionally independent and identically distributed given $\mathcal F_n$,
\[
\begin{aligned}
\E[(\mathcal Q_nf)^2]
&=
\E\left[
\E[f(\bm X_n^{(1)})\mid\mathcal F_n]
\E[f(\bm X_n^{(2)})\mid\mathcal F_n]
\right]\\
&=
\E[f(\bm X_n^{(1)})f(\bm X_n^{(2)})].
\end{aligned}
\]
The assumed joint weak convergence and bounded continuity of $(\bm x_1,\bm x_2)\mapsto f(\bm x_1)f(\bm x_2)$ give
\[
\E[f(\bm X_n^{(1)})f(\bm X_n^{(2)})]
\longrightarrow
\E[f(\bm X^{(1)})f(\bm X^{(2)})]
=
(\mathcal Qf)^2,
\]
where the final equality follows from the independence and common law $\mathcal Q$ of the limiting copies. Hence $\mathcal Q_nf\xrightarrow{p}\mathcal Qf$. Let $\{f_k:k\ge1\}$ be a countable convergence-determining class of bounded Lipschitz functions and define
\[
d_0(\mu,\nu)=\sum_{k=1}^\infty 2^{-k}
\frac{|\mu f_k-\nu f_k|}{1+|\mu f_k-\nu f_k|}.
\]
For each fixed $K$, the first $K$ summands converge to zero in probability, while the remaining tail is bounded by $2^{-K}$. Thus $d_0\{\mathcal L(\bm X_n^{(1)}\mid\mathcal F_n),\mathcal Q\}\xrightarrow{p}0$. The metric $d_0$ metrizes weak convergence, as does $d_{\mathrm{BL}}$; because the limit $\mathcal Q$ is deterministic, convergence in probability under one such metric implies convergence in probability under the other.
\end{proof}
\begin{proposition}\label{prop:supp-randomization-residual}
Fix $b$, not necessarily equal to $\Delta$, and let $\bm e_n=\bm M_H\bm r(b)$ be the within-stratum residual of the observed score. Under Assumptions~\ref{ass:supp-general-allocation} and \ref{ass:supp-participants},
\begin{equation}\label{eq:supp-regenerated-expansion}
\frac{(\bm Z^{(1)})^\top\bm e_n}{\sqrt n}
=
\frac1{\sqrt n}\sum_iZ_i^{(1)}
\left\{\nu_i^Q(b)+\frac12Z_i^{(0)}\nu_i^\delta(b)\right\}
+\frac12\bm d(b)^\top\bm G_n^{(01)}
+o_p(1).
\end{equation}
The variance parameter of its limiting conditional normal law is
\begin{equation}\label{eq:supp-sigmaR}
\sigma_{\rm R}^2(b)
=\E[\nu_i^Q(b)^2]
+\frac14\E[\Var[\delta_i(b)\mid S_i]]
+\frac14\bm d(b)^\top\bm\Psi_D\bm d(b).
\end{equation}
By contrast, Lemma~\ref{lem:supp-score-moments} implies
\begin{equation}\label{eq:supp-sigmaS}
\frac{\bm e_n^\top\bm e_n}{n}
\xrightarrow{p}
\sigma_{\rm S}^2(b)
=\E[\nu_i^Q(b)^2]
+\frac14\E[\Var[\delta_i(b)\mid S_i]]
+\frac14\bm d(b)^\top\bm\Pi\bm d(b).
\end{equation}
\end{proposition}

\begin{proof}
Put $f_i=q_i(b)-m_{S_i}(b)+Z_i^{(0)}\delta_i(b)/2$ and $\bm f=(f_1,\ldots,f_n)^\top$. Exactly,
\[
 \bm e_n=M_H\{\bm f-C(\widehat{\bm\beta}_b-\bm\beta_b)\}.
\]
For every fixed stratum, $\sum_{S_i=s}f_i=O_p(\sqrt n)$ by the allocation bound and conditional participant moment bounds. Hence
\[
 n^{-1/2}(\bm Z^{(1)})^\top P_H\bm f
 =n^{-1/2}(\bm D_n^{(1)})^\top N^+H^\top\bm f=o_p(1).
\]
Also $(\bm Z^{(1)})^\top M_H C=O_p(\sqrt n)$: the unprojected term follows from Lemma~\ref{lem:supp-beta-general}, while the projected term is $(\bm D_n^{(1)})^\top N^+H^\top C=O_p(\sqrt n)$. Multiplication by $\widehat{\bm\beta}_b-\bm\beta_b=O_p(n^{-1/2})$ makes its contribution after division by $\sqrt n$ negligible. Decomposing $\delta_i=d_{S_i}+\nu_i^\delta$ now gives \eqref{eq:supp-regenerated-expansion}.

Conditional on profiles and paths, the first sum contains independent centred participant terms. Its variance tends to $E[(\nu_i^Q)^2]+E[(\nu_i^\delta)^2]/4$: the cross term is a fixed linear combination of $\bm D_n^{(0)}/n$. Assumption~\ref{ass:supp-general-allocation} supplies the limit of the treatment-sign product process, independently in the joint Gaussian limit of these participant terms. This identifies the variance parameter of the limiting conditional normal law; that conditional convergence is proved next. Finally $E[\delta_i^2]=E[\operatorname{Var}(\delta_i\mid S_i)]+\bm d^\top\Pi\bm d$ proves \eqref{eq:supp-sigmaS}.
\end{proof}

\begin{theorem}\label{thm:supp-conditional-R}
Fix $b$, not necessarily equal to $\Delta$. Under Assumptions~\ref{ass:supp-general-allocation} and \ref{ass:supp-participants}, conditionally on $\mathcal F_n$,
\begin{equation}\label{eq:supp-conditional-R}
d_{\mathrm{BL}}\left\{
\mathcal L\left[
\begin{pmatrix}
\bm D_n^{(1)}/\sqrt n\\[1mm]
(\bm Z^{(1)})^\top\bm e_n/\sqrt n
\end{pmatrix}
\middle|\mathcal F_n
\right],
\mathcal N\left[
\bm0,
\begin{pmatrix}
\bm\Sigma_D&\bm0\\
\bm0^\top&\sigma_{\rm R}^2(b)
\end{pmatrix}
\right]
\right\}
\xrightarrow{p}0.
\end{equation}
The same convergence holds with $\mathcal F_n$ replaced by the observed-data sigma-field $\mathcal O_n$. Consequently,
\begin{equation}\label{eq:supp-vR}
d_{\mathrm{BL}}\left(
\mathcal L\left(
\left.
\frac{T_n^*\{\bm r(b)\}}{\sqrt n}
\right|
\mathcal O_n
\right),
\mathcal N\{0,v_{\rm R}(b)\}
\right)
\xrightarrow{p}0,
\end{equation}
where
\[
v_{\rm R}(b)
=\frac14\left[
\bm m(b)^\top\bm\Sigma_D\bm m(b)
+\sigma_{\rm R}^2(b)
\right].
\]
\end{theorem}

\begin{proof}
For regenerated copies $r=1,2$, let $\xi_i^{(r)}=Z_i^{(r)}\nu_i^Q+Z_i^{(0)}Z_i^{(r)}\nu_i^\delta/2$. Conditional on profiles and all three paths, these are centred independent vectors across participants. Their diagonal conditional variances converge to $E[(\nu_i^Q)^2]+E[(\nu_i^\delta)^2]/4$. The off-diagonal conditional covariance of their normalised sums is
\begin{align*}
 &\frac1n\sum_iZ_i^{(1)}Z_i^{(2)}
 \{E[(\nu_i^Q)^2\mid S_i]+\tfrac14E[(\nu_i^\delta)^2\mid S_i]\}\\
 &\quad+\frac1n\sum_iZ_i^{(0)}Z_i^{(1)}Z_i^{(2)}E[\nu_i^Q\nu_i^\delta\mid S_i].
\end{align*}
These are fixed linear combinations of $\bm G_n^{(12)}/\sqrt n$ and $\bm K_n^{(012)}/\sqrt n$, and both vanish by Assumption~\ref{ass:supp-general-allocation}. The participant moment condition gives conditional Lindeberg convergence. Multiplying the resulting conditional characteristic functions by the characteristic functions of the allocation blocks, then using bounded convergence in probability, yields joint Gaussian limits in which the participant transforms are independent of the allocation limits. Distinct pair-product and imbalance blocks are uncorrelated by Assumption~\ref{ass:supp-general-allocation}, and hence independent in this joint Gaussian limit.

The two vectors consisting of regenerated imbalance and regenerated residual score therefore converge jointly to independent copies of the displayed Gaussian law. They are conditionally independent and identically distributed given $\mathcal F_n$. Lemma~\ref{lem:supp-conditional-convergence} proves conditional convergence. For bounded Lipschitz functions the discrepancy is bounded, so convergence in probability also gives convergence in $L^1$. Conditional expectation onto $\mathcal O_n\subset\mathcal F_n$ proves the same conclusion there; tightness upgrades the countable determining class to the bounded-Lipschitz metric. The stratum means converge by Lemma~\ref{lem:supp-score-moments}; the exact score projection gives \eqref{eq:supp-vR}. Conditioning additionally on $\mathcal A_n$ leaves the regenerated law unchanged, as explained above.
\end{proof}

\subsection{Estimating the covariance correction}\label{sec:supp-dhat}
For stratum $s$, define
\[
n_{as}=\sum_i\mathbf 1(A_i=a,S_i=s),
\qquad
\bar Y_{as}=n_{as}^{-1}\sum_i\mathbf 1(A_i=a,S_i=s)Y_i,
\]
when $n_{as}>0$. Set $\widehat d_s(b)=0$ if either treatment group is absent from stratum $s$; otherwise let
\begin{equation}\label{eq:supp-dhat}
\widehat d_s(b)=\bar Y_{1s}-\bar Y_{0s}-b.
\end{equation}
At the design stage, the corresponding stratum-specific deviations may instead be calculated directly from the prespecified outcome-generating model; in that case the consistency requirement is immediate when the model values are used exactly.

There is a separate finite-sample issue when an estimated deviation vector is substituted into its quadratic form. Condition on the profiles, assignments and the outcome-independent allocation covariance simulations. Write $M=\widehat\Psi_n-\widehat\Pi_n$. If both arms are represented in every stratum, then $\E(\widehat{\bm d}\mid S,A,M)=\bm d$, the estimation errors in different strata are conditionally independent, and
\[
 \E(\widehat{\bm d}^{\top}M\widehat{\bm d}\mid S,A,M)
 =\bm d^\top M\bm d+\sum_sM_{ss}\tau_s^2,
 \quad
 \tau_s^2=\frac{\sigma_{1s}^2}{n_{1s}}+\frac{\sigma_{0s}^2}{n_{0s}}.
\]
Here $\sigma_{as}^2=\Var\{Y_i(a)\mid S_i=s\}$. If $n_{as}\ge2$ for every arm and stratum, replacing $\sigma_{as}^2$ by the unbiased within-arm sample variance gives an unbiased conditional estimate of each $\tau_s^2$. Subtracting $\sum_sM_{ss}\widehat\tau_s^2$ therefore removes the conditional mean contribution of estimation noise to this quadratic form on that complete-data event. This statement does not extend to an empty arm filled with zero, nor does it imply unbiasedness of the floored variance or of a rejection probability. Under fixed $J$ and the stated positive profile probabilities, $\max_s\widehat\tau_s^2=O_p(n^{-1})$, so the subtraction is $O_p(n^{-1})$ and has no first-order effect after the variance is divided by $n$. The data-based estimator without this subtraction remains the primary implementation; the subtraction is evaluated separately as an exploratory refinement.


\begin{lemma}\label{lem:supp-dhat}
Under Assumptions~\ref{ass:supp-general-allocation} and \ref{ass:supp-participants},
\[
\max_{1\le s\le J}|\widehat d_s(b)-d_s(b)|\xrightarrow{p}0.
\]
The probability that every positive-probability stratum contains participants from both treatment groups tends to one.
\end{lemma}

\begin{proof}
Because $D_{n,s}=O_p(\sqrt n)$ and $n_s/n\xrightarrow{p}\pi_s>0$,
\[
\frac{n_{1s}}{n_s}=\frac12+\frac{D_{n,s}}{2n_s}\xrightarrow{p}\frac12.
\]
The treatment assignments are independent of the potential outcomes conditional on the baseline sequence. Weighted laws of large numbers therefore give
$\bar Y_{1s}\xrightarrow{p}\E[Y_i(1)\mid S_i=s]$ and
$\bar Y_{0s}\xrightarrow{p}\E[Y_i(0)\mid S_i=s]$. Fixed $J$ gives uniformity over strata.
\end{proof}

Define the raw variance estimate for the randomisation distribution
\begin{equation}\label{eq:supp-VRraw}
\widehat V_{{\rm R},n}^{\rm raw}(b)
=\widehat V_{{\rm S},n}\{\bm r(b)\}
+\frac n{16}\widehat{\bm d}(b)^\top
\{\widehat{\bm\Psi}_n-\widehat{\bm\Pi}_n\}
\widehat{\bm d}(b),
\end{equation}
and set
\[
\widehat V_{{\rm R},n}(b)
=\max\{\widehat V_{{\rm R},n}^{\rm raw}(b),
        \epsilon_n\widehat V_{{\rm S},n}\{\bm r(b)\}\},
\qquad \epsilon_n\downarrow0.
\]
On positive estimated-variance events define $Z_{{\rm R},n}(b)=T_n\{\bm r(b)\}/\sqrt{\widehat V_{{\rm R},n}(b)}$. The associated diagnostic ratio is
\begin{equation}\label{eq:supp-rho}
\widehat\rho_n(b)
=
\frac{\widehat V_{{\rm R},n}^{\rm raw}(b)}
{\widehat V_{{\rm S},n}\{\bm r(b)\}}
=
1+
\frac{
\frac n{16}\widehat{\bm d}(b)^\top
\{\widehat{\bm\Psi}_n-\widehat{\bm\Pi}_n\}
\widehat{\bm d}(b)
}{
\widehat V_{{\rm S},n}\{\bm r(b)\}
}.
\end{equation}
The raw variance is used in this ratio so that activation of the floor does not conceal an unstable estimate of the product-process covariance.

\begin{theorem}\label{thm:supp-SIGA-R}
Fix $b$, not necessarily equal to $\Delta$. In addition to the conditions of Theorem~\ref{thm:supp-conditional-R}, suppose Assumption~\ref{ass:supp-calibration} holds, $\widehat{\bm d}(b)\xrightarrow{p}\bm d(b)$ and $v_{\rm R}(b)>0$, where $\widehat{\bm d}$ is measurable with respect to $\mathcal O_n\vee\mathcal A_n$ or is a deterministic model quantity. Then
\[
\frac{\widehat V_{{\rm R},n}(b)}n
\xrightarrow{p} v_{\rm R}(b).
\]
Moreover,
\[
\sup_x\left|
\Pbb_{\Rcal}\left[
\frac{T_n^*\{\bm r(b)\}}
{\widehat V_{{\rm R},n}(b)^{1/2}}
\le x\middle|\mathcal F_n\vee\mathcal A_n
\right]-\Phi(x)
\right|
\xrightarrow{p}0.
\]
The same statement holds conditionally on $\mathcal O_n\vee\mathcal A_n$. Hence the p-values from the ideal version of the reference randomisation test and from SIGA-R are asymptotically equivalent for one- and two-sided tests, without assuming equality of the two limiting variances. This conclusion applies both at a weak-null value and under a fixed alternative used in a design-stage power calculation.
\end{theorem}

\begin{proof}
For any fixed $b$,
\[
\E[\delta_i(b)^2]
=\E[\Var[\delta_i(b)\mid S_i]]
+\bm d(b)^\top\bm\Pi\bm d(b).
\]
Proposition~\ref{prop:supp-fixed-b-quadratic}, Corollary~\ref{prop:supp-calibration} and the assumed consistency of $\widehat{\bm d}(b)$ show that the correction in \eqref{eq:supp-VRraw} replaces the $\bm\Pi$ quadratic form in $\widetilde v_{\rm S}(b)$ by its $\bm\Psi_D$ counterpart. Thus the raw variance converges to $v_{\rm R}(b)$, and the floor is inactive with probability tending to one. The conditional normal-limit conclusion follows from Theorem~\ref{thm:supp-conditional-R} and conditional Slutsky. Continuity of the normal distribution gives the p-value result; the $\mathcal O_n$ version follows as in Theorem~\ref{thm:supp-conditional-R}.


More explicitly, let $F_n$ be the conditional distribution function of $T_n^*/\sqrt{\widehat V_{{\rm R},n}}$ and put $\varepsilon_n=\sup_x|F_n(x)-\Phi(x)|$. Conditional weak convergence to a continuous distribution implies $\varepsilon_n\to_p0$. Since $\Phi$ is continuous, the same bound holds for $|F_n(x-)-\Phi(x)|$. At any observed standardised threshold $t_n$, the inclusive upper tail is $1-F_n(t_n-)$ and the lower tail is $F_n(t_n)$, so their absolute differences from the Gaussian tails are at most $\varepsilon_n$. For $|t_n|>0$, the inclusive absolute-value tail is $F_n(-|t_n|)+1-F_n(|t_n|-)$ and the error is at most $2\varepsilon_n$; at zero both absolute-value tails equal one. These bounds are uniform in the threshold, including thresholds that diverge under a fixed alternative. Thus the asserted absolute $p$-value equivalence does not require a moderate-deviation argument.
\end{proof}

\subsection{Equality of limiting variances, weak-null size and local power}
Let $\Phi$ denote the standard normal distribution function and let $z_\gamma=\Phi^{-1}(\gamma)$.

\begin{corollary}\label{cor:supp-variance-equality}
At $\Delta=b$,
\begin{equation}\label{eq:supp-gap}
v_{\rm R}(b)-v_{\rm S}(b)
=\frac1{16}\bm d(b)^\top
(\bm\Psi_D-\bm\Pi)\bm d(b).
\end{equation}
Thus SIGA-S and the conditional randomisation distribution of the reference test agree to first order if and only if their limiting variances agree, equivalently when
\[
\bm d(b)^\top\bm\Psi_D\bm d(b)
=\bm d(b)^\top\bm\Pi\bm d(b).
\]
This condition is automatic when $\bm d(b)=\bm0$, including the sharp null $Y_i(1)-Y_i(0)=b$ whenever that individual-level hypothesis is compatible with the outcome support, and null values that hold on average within every stratum.
\end{corollary}

Under label symmetry, each sign has conditional mean zero and variance one given the profiles. The conditional sign covariance in the following identity is therefore also the correlation matrix of the treatment indicators. Its entries relate assignments of different participants within the same sequence. Independent repetitions of that sequence give the entrywise square in the identity.

\begin{proposition}\label{prop:supp-sign}
Let $\mathcal C_n=\E(\bm Z\bm Z^\top\mid S_1,\ldots,S_n)$ under the label-symmetric allocation rule, and let $\bm H$ be its profile-indicator matrix. Then
\begin{align}
 \Cov(\bm G_n^{(01)}\mid S_1,\ldots,S_n)
   &=n^{-1}\bm H^\top(\mathcal C_n\circ\mathcal C_n)\bm H,\label{eq:supp-square-covariance}\\
 v^\top\{\Cov(\bm G_n^{(01)}\mid S_1,\ldots,S_n)-\bm N/n\}v
   &=\frac2n\sum_{i<j}v_{S_i}v_{S_j}(\mathcal C_n)_{ij}^2.\label{eq:supp-offdiagonal-square}
\end{align}
Under Assumption~\ref{ass:supp-general-allocation}, the expectations of these covariance matrices converge to $\bm\Psi_D$. For $J\ge2$, let $B_\pi$ be an orthonormal basis matrix for $(\sqrt{\bm\pi})^\perp$ and define
\[
 \mathcal R_\pi=B_\pi^\top\bm\Pi^{-1/2}\bm\Psi_D\bm\Pi^{-1/2}B_\pi.
\]
On $\bm\pi^\top v=0$, the extrema of $v^\top\bm\Psi_Dv/(v^\top\bm\Pi v)$ are the smallest and largest eigenvalues of $\mathcal R_\pi$. In particular, $v^\top(\bm\Psi_D-\bm\Pi)v\ge0$ for every such $v$ if and only if $\lambda_{\min}(\mathcal R_\pi)\ge1$, and the quadratic form vanishes for every such $v$ if and only if $\mathcal R_\pi=I_{J-1}$.
\end{proposition}
\begin{proof}
Label symmetry gives $\E(Z_i^{(a)}\mid S_1,\ldots,S_n)=0$. The two allocation sequences are conditionally independent, so the product has conditional mean zero and
\[
 \E[Z_i^{(0)}Z_i^{(1)}Z_j^{(0)}Z_j^{(1)}\mid S_1,\ldots,S_n]
       =(\mathcal C_n)_{ij}^2.
\]
Summing with the profile indicators proves \eqref{eq:supp-square-covariance}. Each diagonal entry of $\mathcal C_n$ equals one; removing those terms proves \eqref{eq:supp-offdiagonal-square}. The unconditional covariance is the expectation of the conditional covariance because the conditional mean is zero. The second-moment convergence in Assumption~\ref{ass:supp-general-allocation} identifies its limit as $\bm\Psi_D$.

If $\bm\pi^\top v=0$, write $u=\bm\Pi^{1/2}v=B_\pi t$. Then $v^\top\bm\Pi v=t^\top t$ and $v^\top\bm\Psi_Dv=t^\top\mathcal R_\pi t$. The Rayleigh quotient gives the extrema and the two equivalences.

To verify that the negative directions are relevant to the outcome class, choose any nonzero mean-zero $v$ with $v^\top(\bm\Psi_D-\bm\Pi)v<0$. Let $\epsilon$ be independent of $S$, centred with variance $\sigma^2>0$ and all moments finite, and use an intercept-only adjustment with
\[
 Y(0)=\epsilon-\tfrac{t}{2}v_S,\qquad
 Y(1)=\epsilon+b+\tfrac{t}{2}v_S,\qquad t\ne0.
\]
This law satisfies the participant assumptions and $\Delta=b$, with $\delta=t v_S$ and $Q=\epsilon$. Hence $\bm d=t v$, both limiting variances are positive, and their difference is negative. Corollary~\ref{cor:supp-size} gives anti-conservativeness. The same construction proves necessity of the criterion over the admissible continuous-outcome class. For $J=1$, the only mean-zero effect vector is zero, so the criterion is vacuous.
\end{proof}

\begin{corollary}\label{cor:supp-stratified-sign}
Let $\mathcal S_n=\sigma(S_1,\ldots,S_n)$ and retain the label symmetry and conditional independence of regenerated sequences used in Proposition~\ref{prop:supp-sign}. If $(\mathcal C_n)_{ij}=0$ whenever $S_i\ne S_j$, then, for every $n$,
\begin{equation}\label{eq:supp-stratified-sign}
 \Cov(\bm G_n^{(01)}\mid\mathcal S_n)-\bm N/n
 =\diag\left\{\frac2n\sum_{\substack{i<j\,:\ S_i=S_j=s}}
                  (\mathcal C_n)_{ij}^{\,2}:s=1,\ldots,J\right\}\succeq0.
\end{equation}
In particular, the condition holds when the assignment vectors in different strata are independent conditional on $\mathcal S_n$. Under Assumption~\ref{ass:supp-general-allocation}, $\bm\Psi_D-\bm\Pi\succeq0$. Under the participant assumptions, this yields $v_{\rm R}(b)\ge v_{\rm S}(b)$ at $\Delta=b$ and the limiting null rejection probabilities stated in the main-text corollary.
\end{corollary}
\begin{proof}
The off-diagonal profile entries in \eqref{eq:supp-square-covariance} contain only products with $S_i\ne S_j$ and vanish. In a diagonal entry the $i=j$ terms sum to $N_s/n$, since $Z_i^2=1$. Pairing the remaining $i,j$ and $j,i$ terms gives \eqref{eq:supp-stratified-sign}. Conditional independence across strata implies the required zero covariance because label symmetry gives $\E(Z_i\mid\mathcal S_n)=0$. Taking expectations, using $\E(\bm N/n)=\bm\Pi$, and then applying the second-moment convergence in Assumption~\ref{ass:supp-general-allocation} proves the limiting matrix inequality. Equation~\eqref{eq:supp-gap} and Corollary~\ref{cor:supp-size} give the result for the test. Equality is allowed.
\end{proof}
This conclusion uses zero between-stratum assignment covariances, a weaker property than full conditional independence. From \eqref{eq:supp-offdiagonal-square}, a negative quadratic form requires a contribution from different strata with effect weights of opposite signs. Dependence between strata is therefore necessary for a negative limiting difference in this setting, but it is not sufficient. Appendix~I shows that the finite-sample inequality fails for the range method with three or more factors.

For known $\bm\pi$, let $\widehat{\mathcal R}_\pi$ replace $\bm\Psi_D$ by an estimate. The variational characterisation gives
\[
 |\lambda_{\min}(\widehat{\mathcal R}_\pi)-\lambda_{\min}(\mathcal R_\pi)|
 \le\|\bm\Pi^{-1/2}(\widehat{\bm\Psi}-\bm\Psi_D)\bm\Pi^{-1/2}\|_{\rm op}.
\]
Thus a simulated negative direction requires an error assessment before it is interpreted as a negative limiting direction. There are two separate errors when the simulation uses a finite trial size: estimation of that finite-size covariance and its difference from the limiting covariance. Numerical nonnegativity at selected parameters likewise gives no universal sign theorem. With estimated $\bm\pi$, consistent covariance estimation and a locally continuous choice of basis give the same limiting eigenvalues. An equivalent computation uses any basis of $\bm\pi^\perp$ and a generalised eigenvalue problem with denominator $\bm\Pi$; the released code uses this form.

\begin{corollary}\label{cor:supp-size}
Under Assumptions~\ref{ass:supp-general-allocation} and \ref{ass:supp-participants}, suppose $v_{\rm S}(b),v_{\rm R}(b)>0$ and $0<\alpha<1/2$. For the ideal upper-tail reference test, the limiting rejection probability at $\Delta=b$ is
\begin{equation}\label{eq:supp-size}
1-\Phi\left\{z_{1-\alpha}\sqrt{v_{\rm R}(b)/v_{\rm S}(b)}\right\}.
\end{equation}
For the absolute-value two-sided test with nominal level $\alpha$, the corresponding limit is
\begin{equation}\label{eq:supp-size-two-sided}
2\left[1-\Phi\left\{z_{1-\alpha/2}\sqrt{v_{\rm R}(b)/v_{\rm S}(b)}\right\}\right].
\end{equation}
Both limits equal their nominal levels when the two limiting variances are equal, are conservative when $v_{\rm R}(b)>v_{\rm S}(b)$ and are anti-conservative when $v_{\rm R}(b)<v_{\rm S}(b)$.
\end{corollary}

\begin{proof}
Theorem~\ref{thm:supp-conditional-R}, continuity of the positive-variance normal distribution and conditional quantile convergence imply that the conditional upper-tail and absolute-value critical values are asymptotically $z_{1-\alpha}\sqrt{n v_{\rm R}(b)}$ and $z_{1-\alpha/2}\sqrt{n v_{\rm R}(b)}$, respectively. Proposition~\ref{prop:supp-observed-expansion} and Lemma~\ref{lem:supp-participant-transform} give the observed central limit theorem with variance $v_{\rm S}(b)$, without requiring a covariance estimate. Standardisation and symmetry of the limit at the null value give \eqref{eq:supp-size} and \eqref{eq:supp-size-two-sided}.
\end{proof}

\begin{corollary}\label{cor:supp-local-R}
Under Assumptions~\ref{ass:supp-general-allocation} and \ref{ass:supp-local}, suppose $v_{\rm S}(b),v_{\rm R}(b)>0$. For SIGA-R, also impose Assumption~\ref{ass:supp-calibration} and $\widehat{\bm d}_n(b)\to_p\bm d(b)$ for the limiting law at the null value. The limiting rejection probability of the ideal upper-tail reference test, and of SIGA-R, is
\begin{equation}\label{eq:supp-local-R}
1-\Phi\left[
\frac{z_{1-\alpha}\sqrt{v_{\rm R}(b)}-h/4}
{\sqrt{v_{\rm S}(b)}}
\right].
\end{equation}
For the absolute-value two-sided test, the corresponding limit is
\begin{align}
&1-\Phi\left[
\frac{z_{1-\alpha/2}\sqrt{v_{\rm R}(b)}-h/4}
{\sqrt{v_{\rm S}(b)}}
\right]
+\Phi\left[
\frac{-z_{1-\alpha/2}\sqrt{v_{\rm R}(b)}-h/4}
{\sqrt{v_{\rm S}(b)}}
\right].
\label{eq:supp-local-R-two-sided}
\end{align}
When the two limiting variances are equal, these expressions equal the corresponding local-power limits of SIGA-S.
\end{corollary}
\begin{proof}
The allocation process is the same in every row. Repeating the conditional characteristic-function argument of Theorem~\ref{thm:supp-conditional-R}, with row-specific participant residuals, gives the centred conditional normal limit with variance $v_{\rm R}(b)$. The uniform moment conditions imply conditional Lindeberg convergence and the row-wise laws of large numbers; the convergent stratum moments identify that variance. Thus the ideal conditional critical values, divided by $\sqrt n$, converge to $z_{1-\alpha}\sqrt{v_{\rm R}(b)}$ or $z_{1-\alpha/2}\sqrt{v_{\rm R}(b)}$. The observed statistic has limit $N\{h/4,v_{\rm S}(b)\}$ by the proof of Corollary~\ref{cor:supp-local-S}. Applying the critical values to that limit gives the displayed formulas. For SIGA-R, row-wise score-quadratic convergence, covariance consistency and the stated consistency of $\widehat{\bm d}_n$ give the same critical-value limits. The model-implied row vector satisfies this consistency by Assumption~\ref{ass:supp-local}; the sample-mean estimator in \eqref{eq:supp-dhat} does so by the same uniform moment and treatment-fraction arguments as Lemma~\ref{lem:supp-dhat}. Monte Carlo reference tests have these same limits when $B_n\to\infty$.
\end{proof}

\begin{corollary}\label{cor:supp-NI-EQ}
Under Assumptions~\ref{ass:supp-general-allocation}, \ref{ass:supp-participants} and \ref{ass:supp-calibration}, all statements are pointwise in the participant law and require positive relevant variance limits. At a fixed lower non-inferiority null value $L$, SIGA-S has asymptotic upper-tail size $\alpha$. At any fixed law with $\Delta<L$ and a positive score-quadratic limit, its rejection probability tends to zero. Thus it controls the one-sided null pointwise. The SIGA-S intersection--union test for fixed $L<U$ has pointwise asymptotic level at most $\alpha$. For the reference test and SIGA-R, substitute their component size limits; equality of the limiting variances gives exact calibration at the null value, and $v_R\ge v_S>0$ gives conservative level control.
\end{corollary}
\begin{proof}
Away from a null value $b$, \eqref{eq:supp-observed-expansion} implies $T_n\{\bm r(b)\}/n\to(\Delta-b)/4$, and the score quadratic divided by $n$ converges to a finite positive value under the stated condition. The test statistic therefore diverges in the appropriate direction. At the null value apply Theorem~\ref{thm:supp-sampling} or Corollary~\ref{cor:supp-size}. Under either part of the equivalence null, the intersection--union rejection event is contained in that component rejection event.
\end{proof}

\begin{corollary}\label{cor:supp-MC}
Under the conditions of Theorem~\ref{thm:supp-SIGA-R}, let $p_{\rm RT}$ be the ideal conditional randomisation $p$-value and let $\widehat p_{{\rm RT},B}$ be its inclusive plus-one estimate from $B$ independent regenerated paths. Then, conditionally on the observed data, $\widehat p_{{\rm RT},B}-p_{\rm RT}=O_p(B^{-1/2})$, and $|\widehat p_{{\rm RT},B}-p_{{\rm SIGA-R}}|=o_p(1)+O_p(B^{-1/2})$.
\end{corollary}

\begin{proof}
Conditional on the observed data, the exceedance indicators are independent Bernoulli variables with success probability $p_{\rm RT}$. The conditional bias of the plus-one estimate is $(1-p_{\rm RT})/(B+1)$ and its variance is $Bp_{\rm RT}(1-p_{\rm RT})/(B+1)^2$. These bounds give the first order statement; Theorem~\ref{thm:supp-SIGA-R} and the triangle inequality give the second. Taking $B=B_n\to\infty$ gives the stated equivalence generally; with a fixed $B$, the Monte Carlo error need not vanish.
\end{proof}

\subsection{The variance-corrected reference test}\label{sec:supp-crt}
Let $\widehat\lambda_n(b)=[\widehat V_{{\rm R},n}(b)/\widehat V_{{\rm S},n}\{\bm r(b)\}]^{1/2}$ and define the corrected two-sided Monte Carlo $p$-value
\begin{equation}\label{eq:supp-crt-p}
 \widehat p_{\rm CRT}=\frac{1+\sum_{j=1}^B
 \mathbf 1\{|T^*_{n,j}|\ge\widehat\lambda_n(b)\,|T_n|\ \text{or}\ |T^*_{n,j}|=|T_n|\}}{B+1},
\end{equation}
with one-sided versions using $T^*_{n,j}\ge\widehat\lambda_n(b)T_n$ or $T^*_{n,j}\le\widehat\lambda_n(b)T_n$, in each case also counting the ties $T^*_{n,j}=T_n$; a regenerated statistic tied with the observed statistic thus always counts as at least as extreme, as in the reference test. The following theorem is the corrected-test theorem of Section~3.4 of the main text.

\begin{theorem}\label{thm:supp-crt}
Under Assumptions~\ref{ass:supp-general-allocation}, \ref{ass:supp-participants} and \ref{ass:supp-calibration}, at $\Delta=b$ with $v_{\rm S}(b),v_{\rm R}(b)>0$, $\widehat{\bm d}(b)\to_p\bm d(b)$ and $B=B_n\to\infty$, the one- and two-sided versions of $\widehat p_{\rm CRT}$ converge in distribution to the uniform law on $(0,1)$. Under Assumption~\ref{ass:supp-local}, with $\widehat{\bm d}_n(b)\to_p\bm d(b)$ for the limiting law at the null value, the upper-tail rejection probability converges to $1-\Phi\{z_{1-\alpha}-h/(4\sqrt{v_{\rm S}(b)})\}$. Under the sharp null $Y_i(1)-Y_i(0)=b$ with the estimator \eqref{eq:supp-dhat}, $\widehat\lambda_n(b)^2=1+O_p(n^{-1})$.
\end{theorem}

\begin{proof}
Let $Z_{{\rm S},n}=T_n/\widehat V_{{\rm S},n}^{1/2}$ and let $\varepsilon_n$ be the Kolmogorov distance between the conditional distribution of $T_n^*/\widehat V_{{\rm R},n}^{1/2}$ given $\mathcal O_n\vee\mathcal A_n$ and the standard normal law. Theorem~\ref{thm:supp-SIGA-R} gives $\varepsilon_n\to_p0$. Without the original-scale tie clause, the ideal two-sided $p$-value differs from $2\{1-\Phi(|Z_{{\rm S},n}|)\}$ by at most $2\varepsilon_n$, using cdf left limits for inclusive tails. A cdf within $\varepsilon_n$ of a continuous cdf has every atom bounded by $2\varepsilon_n$. The event $|T_n^*|=|T_n|$ contains at most two atoms and therefore has conditional probability at most $4\varepsilon_n$. The tie-inclusive $p$-value thus differs from the normal tail by at most $6\varepsilon_n$. This uniform bound holds even though the threshold is observed-data dependent. The sampling theorem and the continuous mapping theorem give convergence to the uniform law. The Monte Carlo estimate adds $O_p(B^{-1/2})$ error, conditionally on the observed data and covariance estimates. The same argument applies to one-sided versions, whose tie event contains one atom.

Under Assumption~\ref{ass:supp-local}, the conditional normal approximation and the vanishing atom bound continue to hold by Corollary~\ref{cor:supp-local-R}. The upper-tail corrected $p$-value is $1-\Phi(Z_{{\rm S},n})+o_p(1)$. Since $Z_{{\rm S},n}\Rightarrow\mathcal N\{h/(4\sqrt{v_{\rm S}(b)}),1\}$ by Corollary~\ref{cor:supp-local-S}, continuity at the limiting critical value gives the stated rejection probability. No exact finite-sample quantile equivalence is used for the extra tie clause.

Under the sharp null, $\delta_i(b)=0$, so $\widehat d_s(b)$ is the within-stratum difference between the two arms of the sample means of $Y_i(0)$ and is $O_p(n^{-1/2})$ for each $s$. Indeed, conditional on profiles and assignments, the two arm means have the same stratum mean and their difference has variance $\operatorname{Var}\{Y_i(0)\mid S_i=s\}(n_{1s}^{-1}+n_{0s}^{-1})$. Both arm counts divided by $n$ converge to $\pi_s/2>0$, so conditional Chebyshev's inequality gives the rate; the event of an empty arm has probability tending to zero. Since $\widehat{\bm\Psi}_n-\widehat{\bm\Pi}_n=O_p(1)$ and $\widehat V_{{\rm S},n}/n\to_pv_{\rm S}(b)>0$, the correction term in \eqref{eq:supp-VRraw} divided by $\widehat V_{{\rm S},n}$ is $O_p(n^{-1})$, and the floor is inactive with probability tending to one.
\end{proof}


\section{Continuity correction for an unadjusted binary score}\label{sec:supp-lattice}

This appendix defines the continuity-corrected normal approximation, averaged over the distribution of the treated count, that is used for the reported SIGA-S and SIGA-R analyses of unadjusted binary superiority at the null value $b=0$. The restriction $b=0$ refers to the tested superiority null value, not to the treatment effect under which an outer trial is generated; the same tail calculation is therefore used when evaluating either type I error or power. The refinement uses the same observed statistic and the procedure-specific variance estimate---$\widehat V_{{\rm S},n}$ for SIGA-S or $\widehat V_{{\rm R},n}$ for SIGA-R---as the corresponding ordinary Gaussian calculation and changes only the conversion of the statistic into a $p$-value. It is a numerical finite-sample refinement rather than an exact enumeration of the randomisation distribution and is separate from the covariance correction.

For a given outer trial, let
\[
Y_+=\sum_{i=1}^nY_i,
\qquad
K=\sum_{i=1}^nA_i,
\qquad
Y_T=\sum_{i=1}^nA_iY_i.
\]
For the centred unadjusted score $\bm r=\bm Y-\overline Y\bm 1_n$, where $\overline Y=Y_+/n$,
\begin{align}
T_n(\bm r)&=\sum_{i=1}^n(A_i-\tfrac12)(Y_i-Y_+/n)\nonumber\\
 &=Y_T-\frac12Y_+-\frac{Y_+}{n}(K-\tfrac12n)
 =Y_T-\frac{KY_+}{n}.
\label{eq:supp-binarylattice}
\end{align}
Conditional on $K=k$ and the observed total number of successes $Y_+$, the feasible values of the treatment-group success count are
\[
\mathcal X_k=\mathbb Z\cap[\max\{0,k-(n-Y_+)\},\min\{k,Y_+\}].
\]
For $x\in\mathcal X_k$, the corresponding feasible value of the statistic is
\[
t_{k,x}=x-\frac{kY_+}{n}.
\]
Thus, conditional on $K=k$, the statistic lies on a finite lattice with unit spacing, although the location and support of the lattice depend on $k$.

Let $V$ denote the variance estimate used by the procedure, with $V=\widehat V_{{\rm S},n}(\bm r)$ for SIGA-S and $V=\widehat V_{{\rm R},n}(0)$ for SIGA-R, and let $\Phi$ denote the standard normal cumulative distribution function. Assign normal probability mass to the unit-width cell centred at the lattice point $t_{k,x}$ using
\begin{equation}\label{eq:supp-cellmass}
g_{k,x}(V)
=
\Phi\left(\frac{t_{k,x}+1/2}{\sqrt V}\right)
-
\Phi\left(\frac{t_{k,x}-1/2}{\sqrt V}\right).
\end{equation}
Because the continuous normal approximation also assigns probability outside the feasible support for a given $k$, normalise the cell probabilities over $\mathcal X_k$ using
\[
G_k(V)=\sum_{x\in\mathcal X_k}g_{k,x}(V).
\]

In simulation of profiles and assignments $q$, let
\[
K^{(q)}=\sum_{i=1}^nA_i^{(q)}.
\]
The design-stage probability that exactly $k$ participants are assigned to treatment is estimated by
\[
\widehat w_{n,k}
=
\frac1{B_0}
\sum_{q=1}^{B_0}
\mathbf 1\{K^{(q)}=k\},
\qquad
k=0,\ldots,n.
\]
These weights are retained from the relevant randomisation-design covariance estimation. The covariance-estimation simulations for SIGA-S, using one allocation sequence, and SIGA-R, using three sequences, both record the treated-count distribution, and $\sum_{k=0}^n\widehat w_{n,k}=1$.

The inclusive upper-tail probability is
\begin{equation}\label{eq:supp-lattice-upper}
p_{\mathrm{mix}}^{+}
=
\sum_{k=0}^n\widehat w_{n,k}
\frac{\displaystyle\sum_{\substack{x\in\mathcal X_k:\ t_{k,x}\ge T_n(\bm r)}}g_{k,x}(V)}{G_k(V)}.
\end{equation}
The inclusive lower-tail probability is
\begin{equation}\label{eq:supp-lattice-lower}
p_{\mathrm{mix}}^{-}
=
\sum_{k=0}^n\widehat w_{n,k}
\frac{\displaystyle\sum_{\substack{x\in\mathcal X_k:\ t_{k,x}\le T_n(\bm r)}}g_{k,x}(V)}{G_k(V)}.
\end{equation}
The inclusive absolute-value two-sided probability, corresponding to the two-sided tail convention used for the reference randomisation test, is
\begin{equation}\label{eq:supp-lattice-pvalue}
p_{\mathrm{mix}}^{\pm}
=
\sum_{k=0}^n\widehat w_{n,k}
\frac{\displaystyle\sum_{\substack{x\in\mathcal X_k:\ |t_{k,x}|\ge |T_n(\bm r)|}}g_{k,x}(V)}{G_k(V)}.
\end{equation}

The mixture over $k$ accounts for variation in the realised treatment-group size under biased-coin minimisation, while the cell probabilities retain the unit-spaced support of the unadjusted binary statistic conditional on each $k$. The resulting distribution remains a working numerical approximation: the allocation design enters through $V$ and $\widehat w_{n,k}$, but the exact design-specific probabilities of the individual lattice points are not enumerated.


For the tie-inclusive corrected test at $b=0$ with the unadjusted binary score, use $V=\widehat V_{{\rm R},n}(0)$ in these cell probabilities and sum over its actual rejection region. With $\widehat\lambda_n=(\widehat V_{{\rm R},n}/\widehat V_{{\rm S},n})^{1/2}$, the corresponding approximation is
\[
 \sum_{k=0}^n\widehat w_{n,k}
 \frac{\displaystyle\sum_{x\in\mathcal X_k}g_{k,x}(\widehat V_{{\rm R},n})
 \mathbf1\{|t_{k,x}|\ge\widehat\lambda_n|T_n|\ \text{or}\ |t_{k,x}|=|T_n|\}}
 {G_k(\widehat V_{{\rm R},n})}.
\]
The union counts each lattice point once. The uncorrected reference approximation instead uses the region $|t_{k,x}|\ge|T_n|$ with that same reference variance. Merely replacing $\widehat V_{{\rm R},n}$ by $\widehat V_{{\rm S},n}$ in the uncorrected lattice tail is not the same finite-sample calculation. Either the original or reconstructed variance estimates can be used, consistently in both the masses and the scale factor.

This continuity correction is not used for baseline-adjusted binary superiority or for non-inferiority and equivalence because the corresponding scores do not generally reduce to the same one-dimensional, unit-spaced lattice indexed only by $K$ and the treatment-group success count. These analyses use the ordinary Gaussian tail probabilities based on $Z_{{\rm S},n}(b)$ or, when the conditional randomisation distribution is the target, $Z_{{\rm R},n}(b)$. For $\star\in\{{\rm S},{\rm R}\}$, the upper-tail, lower-tail and two-sided Gaussian probabilities are
\[
p_{\star}^{+}(b)=1-\Phi\{Z_{\star,n}(b)\},
\qquad
p_{\star}^{-}(b)=\Phi\{Z_{\star,n}(b)\},
\qquad
p_{\star}^{\pm}(b)=2\left[1-\Phi\{|Z_{\star,n}(b)|\}\right].
\]
Non-inferiority at a lower null value $L$ uses $p_{\star}^{+}(L)$. Equivalence with limits $L<U$ uses the intersection--union $p$-value
\[
p_{\star,\mathrm{EQ}}
=
\max\left\{
p_{\star}^{+}(L),
p_{\star}^{-}(U)
\right\}.
\]

\begin{proposition}\label{prop:lnm}
At the null value $b=0$, suppose $Y_+/n\to p_Y\in(0,1)$ in probability and the procedure-specific variance $V_n/n\to v>0$ in probability. Let the treated-count weights estimated from the randomisation-design simulations be those in this appendix, computed under a fixed known profile law satisfying Assumption~\ref{ass:supp-general-allocation}, or under an estimated range-method profile law satisfying Proposition~\ref{prop:alloc-plugin}. The mixture distribution of the score has Kolmogorov distance $o_p(1)$ from $\mathcal N(0,V_n)$. Hence its inclusive upper-, lower-, and absolute-value two-sided $p$-values differ from the corresponding Gaussian $p$-values by $o_p(1)$, uniformly over the observed score threshold.
\end{proposition}
\begin{proof}
For a candidate treated count $k$, the score lattice is
\[
 t_{k,x}=x-kY_+/n,\qquad
 \max\{0,k-(n-Y_+)\}\le x\le\min\{k,Y_+\},\quad x\in\Z.
\]
Write $y=Y_+/n$. If both $k/n$ and $y$ belong to $[\epsilon,1-\epsilon]$, its lower endpoint is at most $-\epsilon^2n$ and its upper endpoint is at least $\epsilon^2n$. On an event with probability tending to one, $V_n$ is between two fixed positive multiples of $n$. The total normal mass outside the feasible lattice cells is then at most $2\exp(-c n)$, for a constant $c>0$.

Normal cell masses use intervals of width one centred at each lattice point. Before truncation, summing these masses up to a threshold rounds that threshold by at most one unit. The maximum normal density is $(2\pi V_n)^{-1/2}$. After truncation to the feasible cells and normalisation, the conditional lattice cdf therefore differs uniformly from the normal cdf by at most
\[
 C V_n^{-1/2}+C\exp(-cn),
\]
with constants uniform over the stated central range of $k$.

The treated count from the randomisation-design simulations satisfies $K/n=(1+D_{n,0}/n)/2\to1/2$ by Assumption~\ref{ass:supp-general-allocation}. The expected empirical mixture weight outside $k/n\in[\epsilon,1-\epsilon]$ is exactly the corresponding randomisation-design probability and tends to zero. Markov's inequality therefore makes that empirical outside weight $o_p(1)$, for any positive number of independent design-simulation replicates. For the estimated-law case, condition on $\widehat\pi_n$ lying in the compact neighbourhood of Proposition~\ref{prop:alloc-plugin}; its uniform fourth-moment bound gives $\sup_{\rho\in\mathcal P_0}\Pbb_\rho(|D_{n,0}|>c n)=O(n^{-2})$ for fixed $c>0$. The same argument applies uniformly on this event, and its complement has probability tending to zero. Since all component cdfs are bounded by one, mixing the preceding central-range bound proves the Kolmogorov claim. The assumptions on $Y_+/n$ and $V_n/n$ make the central-range event have probability tending to one for a sufficiently small fixed $\epsilon>0$.

Inclusive one-sided tails inherit the same uniform error by taking cdf left limits where appropriate. An absolute-value tail is obtained from two cdf values, with at most twice the error; at threshold zero both absolute-value tails equal one. This proves the $p$-value statement. It is a first-order equivalence result, not a finite-sample bound between the lattice mixture and the true randomisation distribution.
\end{proof}



Under the conditions of Proposition~\ref{prop:lnm}, applied with $V=\widehat V_{{\rm R},n}$, the corrected lattice tail just defined differs by $o_p(1)$ from $2[1-\Phi\{|T_n|/\sqrt{\widehat V_{{\rm S},n}}\}]$. To see this, the cdf bound in that proposition is uniform in the threshold, including the random threshold $\widehat\lambda_n|T_n|$. On the central treated-count event used in its proof, the mass of a single normal cell is at most $C/\sqrt{\widehat V_{{\rm R},n}}$ after normalisation. Mixing preserves this bound; the remaining treated-count weight is $o_p(1)$. The extra raw-scale tie event contains at most two points and consequently has probability $o_p(1)$. The normal tail at the rescaled threshold is exactly the displayed sampling-standardised normal tail. The result is first-order equivalence, not a finite-sample exactness guarantee or a proof that the cell approximation is uniformly more accurate than an ordinary Gaussian tail.

\section{Additional simulation settings}\label{sec:supp-simulation-tables}

The numerical study contains two independent Monte Carlo experiments over the same 56 scenario configurations. Supplementary Tables~1 and~2 give the fixed-alternative settings used in both experiments. Complete null-value and power results are in Supplementary Tables 5 and 6; the main article reports their summaries. Within each experiment, the SIGA procedure and its corresponding reference randomisation test used common outer trials; the two experiments used independent outer simulations.

\subsection{Continuous-outcome study}\label{sec:supp-continuous-model}

The continuous-outcome study provided a regular setting with a homogeneous additive treatment effect across joint strata. It was not designed to make the difference between the two limiting variances artificially large. Four design--sample-size combinations were considered. With two independent binary minimisation factors, the factor prevalences were $\Pbb(X_{i1}=1)=0.50$ and $\Pbb(X_{i2}=1)=0.40$, and the nominal sample sizes were 100 and 500 participants per group. With five independent binary minimisation factors, the corresponding prevalences were
$(0.50,0.40,0.30,0.20,0.10)$, and the nominal sample sizes were 200 and 1000 participants per group.

The control potential outcome was generated as
\[
Y_i(0)
=
\sum_{j=1}^F
\beta_j\{X_{ij}-\Pbb(X_{ij}=1)\}
+
\xi\{X_{i1}X_{i2}-\E[X_{i1}X_{i2}]\}
+
\varepsilon_i,
\]
where $F$ is the number of minimisation factors, $\varepsilon_i\sim\mathcal N(0,1)$ and $\xi=0.25$. For simplicity, the main-effect coefficient vector was set equal to the corresponding vector of factor prevalences. Thus, $\bm\beta=(0.50,0.40)^\top$ with two factors and
$\bm\beta=(0.50,0.40,0.30,0.20,0.10)^\top$ with five factors. The treatment potential outcome was generated as $Y_i(1)=Y_i(0)+\tau$, where $\tau$ is the true additive treatment effect. Consequently, for the SIGA-R calculation at null value $b$, the model-implied stratum-specific deviations were $d_s(b)=\tau-b$, $s=1,\ldots,J$. The adjusted score was obtained from a baseline-only linear working model containing the minimisation-factor main effects but not the interaction term. The unadjusted and adjusted reference randomisation tests used the same fixed score vectors for each null value as the corresponding SIGA procedures.

Superiority was assessed using a two-sided level-$0.05$ test. Non-inferiority was assessed at the lower additive null value $-0.20$ using a one-sided level-$0.025$ test. Equivalence was assessed over the interval $(-0.45,0.45)$ using two one-sided level-$0.05$ tests. Type I error was evaluated at the relevant null value with a constant individual treatment effect and, for equivalence, separately at the lower and upper limits. Fixed alternatives were selected in separate pilot simulations conducted before the final study so that the reference randomisation test had approximately 85\% power, with values between 80\% and 90\% regarded as acceptable. The resulting prespecified alternatives are listed in Supplementary Table~1.

\subsection{Binary-outcome study}\label{sec:supp-binary-model}

The binary-outcome study used the same four minimisation design--sample-size combinations and allocation settings as the continuous-outcome study. Unlike the homogeneous additive continuous model, the common-log-odds model below can produce different risk differences across joint strata because the baseline risks vary across strata. It therefore provides a nonsharp setting without deliberately maximising the difference between the two variance targets.

The conditional marginal potential-outcome probabilities were specified by
\[
\operatorname{logit}\Pbb\{Y_i(a)=1\mid\bm X_i\}
=
\alpha_0
+
\sum_{j=1}^F\gamma_jX_{ij}
+
\zeta X_{i1}X_{i2}
+
\theta a.
\]
For simulation, the observed outcome was drawn independently across participants as a Bernoulli variable with the probability corresponding to the realised assignment. This fully defines the observed-data law without requiring an otherwise unidentified coupling of the two binary potential outcomes. Supplementary Appendix B.1 shows that the variance combinations entering SIGA-S and SIGA-R depend only on the two conditional marginal outcome laws and not on that coupling.

The intercept $\alpha_0$ was calibrated so that the marginal control risk was $0.60$. For each scenario, $\theta$ was determined numerically by summing over the finite covariate distribution so that the resulting marginal risk difference agreed with its prespecified value to numerical tolerance. The first $F$ entries of $(0.35,-0.25,0.20,-0.15,0.10)$ were used as the main-effect coefficients, and $\zeta=0.20$. The adjusted score was obtained by baseline-only linear residualisation of $\bm W(b)$ on the minimisation-factor main effects, with the interaction term omitted.

For the SIGA-R calculation, the model-implied deviation at the null value in stratum $s$ was $d_s(b)=\Pbb\{Y_i(1)=1\mid S_i=s\}-\Pbb\{Y_i(0)=1\mid S_i=s\}
-b$. These values were calculated directly from the prespecified outcome-generating model rather than estimated separately in each outer trial.

Superiority was assessed using a two-sided level-$0.05$ test, and non-inferiority was assessed at the lower null value $L=-0.10$ using a one-sided level-$0.025$ test. Equivalence was assessed using two one-sided level-$0.05$ tests. The symmetric equivalence limits were $(-0.21,0.21)$ for the two-factor design with 100 participants per group and $(-0.15,0.15)$ for the five-factor design with 200 participants per group; both larger designs used $(-0.10,0.10)$. For equivalence, type I error was evaluated separately at the lower and upper limits.

As in the continuous-outcome study, the design-specific treatment effects were selected in separate pilot calculations conducted before the final study to place the power of the reference randomisation test in an informative range of approximately 80--90\%. For the two smaller binary-outcome designs, the equivalence limits were also selected in these pilot calculations for the same purpose. These simulation-specific limits were used only to compare the methods and should not be interpreted as clinically recommended equivalence margins. The prespecified treatment effects and equivalence limits are listed in Supplementary Table~2.

For unadjusted binary superiority at the null value $b=0$, both SIGA-S and SIGA-R used the continuity-corrected normal approximation described in Supplementary Appendix D, with their respective sampling and conditional randomisation variance estimates. Baseline-adjusted binary superiority and all binary non-inferiority and equivalence analyses used ordinary Gaussian tail probabilities.


\subsection{Trial-based reconstruction}\label{sec:supp-trial-model}
We illustrate the use of SIGA for trial planning using the published planning characteristics of SWIFT DIRECT, a multicentre randomised non-inferiority trial comparing direct mechanical thrombectomy with intravenous alteplase followed by thrombectomy in patients with acute ischaemic stroke \citep{FischerEtAl2022Protocol,FischerEtAl2022Trial}. The binary primary endpoint was functional independence, defined as a modified Rankin Scale score of 0--2 at 90 days. The protocol planned a total of 404 participants and reported 80\% power under an assumed response probability of 0.622 in each group, an absolute non-inferiority margin of 0.12 on the treatment-minus-control risk-difference scale, and a one-sided significance level of 0.05. The treatment coefficient in the outcome model is set to zero, so that both marginal response probabilities are 0.622 and the true risk difference is zero.

The illustration retains the five dichotomised factors reported for the trial's minimisation and covariate-adjusted analyses: baseline National Institutes of Health Stroke Scale score ($\le17$ versus $>17$), age ($<70$ versus $\ge70$ years), occlusion location (M1 only versus internal-carotid-artery-related occlusion), tandem-lesion status, and Alberta Stroke Program Early CT Score (4--7 versus 8--10). The working prevalences of the adverse categories are 0.466, 0.592, 0.287, 0.154 and 0.300, respectively. The NIHSS and age prevalences are approximated from the pooled medians and interquartile ranges using truncated-normal working distributions. The occlusion and tandem-lesion prevalences are based on pooled counts, whereas the ASPECTS prevalence is based on a discrete working distribution chosen to reproduce the reported pooled median of 8 and interquartile range of 7--9. Because the joint distribution of the five factors was not reported, they are generated independently. No correlation structure is imposed in this trial-based reconstruction without participant-level information. The illustration is therefore neither a participant-level reanalysis nor an exact reconstruction of the completed trial.

To make the minimisation factors prognostic, binary outcomes are generated from a logistic model containing their five main effects. The coefficients for the adverse categories are set to $(-0.55,-0.25,-0.40,-0.25,-0.60)$ for NIHSS, age, occlusion location, tandem lesion and ASPECTS, respectively. These values represent moderate working prognostic associations and are treated as design assumptions because factor-specific outcome effects were not available from the aggregate reports. No interaction terms are included. The intercept is calibrated so that the marginal response probability is 0.622, and the treatment coefficient is set exactly to zero. The treatment and control groups therefore have identical conditional and marginal response probabilities.

The published trial used deterministic minimisation, but the aggregate publications do not provide the participant-level information required to reconstruct the realised allocation process. To illustrate prospective power evaluation under a fully specified stochastic randomisation design, treatment is assigned using the stochastic Pocock--Simon range method described in Assumption~\ref{ass:supp-design} with $p_{\rm bc}=0.80$. This rule is a prespecified randomised analogue used for the illustration and is not intended to reproduce the completed trial's allocation process exactly. Non-inferiority at the null value $-0.12$ is evaluated using both the unadjusted score and baseline-only linear residualisation on the five factor main effects. SIGA-S, SIGA-R and the reference randomisation test are evaluated on common outer trials. At this nonsharp binary risk-difference null value, the reference randomisation test with the observed score held fixed cannot generally be represented by the individual-level sharp null $Y_i(1)-Y_i(0)=b$ for binary potential outcomes. SIGA-S targets the marginal risk-difference test, which is valid at the null value, whereas SIGA-R targets the conditional distribution of the reference randomisation test with the observed score held fixed; their agreement is therefore assessed empirically rather than assumed.


\renewcommand{\tablename}{Supplementary Table}
\setcounter{table}{0}

\begin{table}[!htbp]
\centering
\caption{Treatment effects used in the continuous-outcome power scenarios.}
\label{tab:power-effects}
\normalsize
\setlength{\tabcolsep}{8pt}
\renewcommand{\arraystretch}{1.10}
\begin{tabular}{ccccc}
\toprule
Factors & Size/group & Superiority $\tau$ & Non-inferiority $\tau$ & Equivalence $\tau$ \\
\midrule
2 & 100 & 0.430 & 0.230 & 0.000 \\
2 & 500 & 0.190 & -0.010 & 0.280 \\
5 & 200 & 0.300 & 0.100 & 0.180 \\
5 & 1000 & 0.135 & -0.065 & 0.330 \\
\bottomrule
\end{tabular}
\end{table}

\clearpage
\begin{table}[!htbp]
\centering
\caption{Marginal risk differences and equivalence limits used in the binary-outcome power scenarios.}
\label{tab:binary-effects}
\small
\setlength{\tabcolsep}{6pt}
\renewcommand{\arraystretch}{1.12}
\begin{tabular}{cccccc}
\toprule
Factors & Size/group & \shortstack{Superiority\\$\Delta$} & \shortstack{Non-inferiority\\$\Delta$} & \shortstack{Equivalence\\limits} & \shortstack{Equivalence\\$\Delta$} \\
\midrule
2 & 100 & 0.200 & 0.100 & $(-0.210,0.210)$ & 0.000 \\
2 & 500 & 0.100 & 0.000 & $(-0.100,0.100)$ & 0.000 \\
5 & 200 & 0.140 & 0.050 & $(-0.150,0.150)$ & 0.000 \\
5 & 1000 & 0.070 & -0.030 & $(-0.100,0.100)$ & 0.040 \\
\bottomrule
\end{tabular}
\end{table}

\section{Computation-time benchmark}\label{sec:supp-timing}

The full operating-characteristic calculations were distributed across multiple Apple computers because simulations consisting of 100,000 outer trials and 4,999 regenerated allocation paths per randomisation test could not be completed on the M3 Ultra workstation alone within the available computing window. To obtain hardware-standardised computation-time estimates for all three procedures, the timing assessment was conducted separately and exclusively on a single Apple M3 Ultra workstation.

The workstation had 28 CPU cores available, but each benchmark process was restricted to one core. For each outcome, analysis, objective, factor configuration and sample size, 30 representative outer trials were used in each of three independent repeats with system sleep disabled, local storage and $B=4{,}999$ regenerated allocation paths per randomisation test. Randomisation-test time was measured directly across the 30 trials in each repeat. Because individual SIGA-S and SIGA-R calls were too short for reliable direct timing, each procedure was first warmed up and then repeated in blocks until both at least one second of cumulative method-specific time and at least 1,000 timed calls had been recorded in each repeat. The resulting mean time per trial was multiplied by 100,000 and divided by 60 to obtain the projected cumulative minutes reported in the main article. These values are single-core projections for 100,000 analyses and are not the observed wall-clock duration of the multi-machine operating-characteristic calculations; the ratio between procedures is per core and is unaffected by parallelisation.

Common data generation, score construction and the one-time randomisation-design covariance calculations were excluded from the analysis-time measurements. Covariance estimation for SIGA-S with one allocation sequence and for SIGA-R with three sequences was timed separately in three benchmark repetitions, with 100,000 randomisation-design simulation replicates in each repetition; only one covariance estimation is required in practice for each design and total sample size. Unadjusted and adjusted analyses were timed separately for both outcome types and all four sample-size settings. For unadjusted binary superiority at $b=0$, both SIGA-S and SIGA-R timings include the continuity correction of Appendix~D; the other analyses use the ordinary Gaussian tail calculation.


\clearpage

\section{Targeted variance-difference sensitivity analysis}\label{sec:supp-variance-validation}

The targeted variance-difference sensitivity analysis examined settings in which the repeated-sampling and conditional randomisation variances could differ more clearly than in the practically motivated benchmark scenarios. It was intended to verify the variance-gap mechanism and the behaviour of SIGA-R as an approximation to the reference randomisation test with the observed score held fixed, rather than to represent typical clinical-trial conditions.

Six independent-factor allocation designs were used. Designs D1 and D2 had two binary minimisation factors with marginal probabilities $(0.50,0.40)$, 100 or 500 participants per group and $p_{\rm bc}=0.80$. Designs D3 and D4 had five binary factors with marginal probabilities $(0.50,0.40,0.30,0.20,0.10)$, 200 or 1000 participants per group and $p_{\rm bc}=0.80$. Designs D5 and D6 used the D1 and D2 factor distributions and sample sizes with the stronger biased coin $p_{\rm bc}=0.95$. All designs used the equal-weight range method specified in the main article.

For every design, 100,000 three-sequence randomisation-design simulations estimated $\widehat{\bm\Gamma}_n$ and $\widehat{\bm\Psi}_n$. For D5 and D6, the maximum-ratio treatment-effect direction was selected using an independent 200,000-replicate allocation covariance estimation and was then evaluated using the separate 100,000-replicate analysis-stage covariance estimate. The absolute probability-weighted cosines between the independently selected direction and the corresponding direction from the analysis-stage covariance estimate were 1.000 for D5 and 0.999 for D6. Each outcome scenario used 100,000 independently generated outer trials. The reference randomisation test regenerated 4,999 paths per outer trial and used the inclusive plus-one upper-tail $p$-value at level 0.025. The model-known vector $\bm d(b)$ was used in SIGA-R so that the comparison isolated the covariance of the treatment-sign product process rather than error from estimating stratum-specific effects.

Scenarios S01, S03, S05 and S07 used the realistic continuous model with deviations proportional to the centred first factor, $\max_s|d_s|=0.50$, outcome standard deviation one and individual-effect standard deviation 0.25. Scenarios S02, S04, S06 and S08 used the common-log-odds binary model from the main simulation study. These eight scenarios represented moderate heterogeneity under the four $p_{\rm bc}=0.80$ designs. Scenarios S09 and S11 used the pure-symmetric continuous model with shared-noise standard deviation 0.05 and $\max_s|d_s|=1.00$ along the independently selected maximum-ratio direction. Scenarios S10 and S12 used the corresponding symmetric binary stress model, scaling the direction toward $\max_s|d_s|=0.80$ subject to valid treatment and control probabilities. The continuous and binary null values were $-0.20$ and $-0.10$, respectively. Supplementary Table~\ref{tab:supp-product-process-scenarios} gives the complete scenario definitions.


The model labels in Supplementary Tables~\ref{tab:supp-product-process-scenarios} and \ref{tab:supp-product-process-sensitivity-results} refer to the following distributions. Write $x_f(s)\in\{0,1\}$ for factor $f$ in profile $s$, $p_f=E\{x_f(S)\}$, and
\[
 \mu_s=\sum_{f=1}^F\beta_f\{x_f(s)-p_f\}
       +.25\{x_1(s)x_2(s)-E[x_1(S)x_2(S)]\},
\]
where $(\beta_1,\ldots,\beta_F)$ is the first $F$ entries of $(.50,.40,.30,.20,.10)$. All scenarios below have $\Delta=b$ and $\sum_s\pi_s d_s=0$.

In S01, S03, S05 and S07 (label \texttt{realistic}), the potential outcomes can be generated as
\[
 Y_i(0)=\mu_{S_i}+\varepsilon_i,\qquad
 Y_i(1)=Y_i(0)+b+d_{S_i}+\zeta_i,
\]
with independent $\varepsilon_i\sim N(0,1)$ and $\zeta_i\sim N(0,.25^2)$, independent also of profiles and allocation. Let $v_s=x_1(s)-p_1$ and divide $v$ by $\max_s|v_s|$; set $d_s=.50v_s$. The common-log-odds binary scenarios use the logistic model in Appendix~E.2, with the treatment coefficient chosen for the stated null value.

For S09 and S11 (label \texttt{pure-symmetric}), set
\[
 Y_i(0)=\varepsilon_i-(b+d_{S_i})/2,\qquad
 Y_i(1)=\varepsilon_i+(b+d_{S_i})/2,
 \qquad \varepsilon_i\sim N(0,.05^2).
\]
Here $d=v$, where $v$ has weighted mean zero and maximum absolute entry one and is selected to maximise
\[
 \frac{v^\top\widehat\Psi_{\rm dir}v}{v^\top\Pi v}
 \quad\text{subject to }\bm\pi^\top v=0,\quad v\ne0.
\]
The matrix $\widehat\Psi_{\rm dir}$ is estimated from the separate 200,000-replicate direction-selection simulations, not from the simulations used to evaluate the procedures. The table label \texttt{robust-max-ratio} identifies this generalised-eigenvector option in the released code.

S10 and S12 (label \texttt{symmetric-stress}) use the same independently chosen direction, with
\[
 d=t v,\quad t=\min(.80,.98t_{\max}),\qquad
 t_{\max}=\sup\{a\ge0:\max_s|b+a v_s|\le.90\},
\]
\[
 \Pbb\{Y_i(0)=1\mid S_i=s\}=\{1-b-d_s\}/2,\qquad
 \Pbb\{Y_i(1)=1\mid S_i=s\}=\{1+b+d_s\}/2.
\]
Observed binary outcomes are drawn independently using the probability for the assigned treatment. These definitions reproduce the model and direction options in the released simulation code \citep{KojimaSIGA2026}; they do not change the numerical results reported below.

In scenarios S01--S08, the mean diagnostic ratio $\widehat\rho_n$ ranged from 1.000 to 1.004 and the largest absolute difference between either SIGA procedure and the reference randomisation test was 0.14 percentage points. In the strongly heterogeneous continuous scenarios S09 and S11, the mean ratio was 1.193--1.194; reference-test rejection probabilities were 1.62--1.64\%, compared with 1.60--1.65\% for SIGA-R and 2.43--2.52\% for SIGA-S. In the strongly heterogeneous binary scenarios S10 and S12, the mean ratio was 1.113--1.115; SIGA-R rejection probabilities were 2.01--2.02\%, compared with 1.77--1.91\% for the reference test and 2.19--2.40\% for SIGA-S. Across all 12 scenarios and both analyses, the largest absolute SIGA-S--reference-test and SIGA-R--reference-test differences were 0.90 and 0.24 percentage points, respectively. The floor in the SIGA-R variance was not active in any scenario--analysis combination. Complete rejection probabilities are reported in Supplementary Table~\ref{tab:supp-product-process-sensitivity-results}.

\clearpage
\begingroup
\small
\setlength{\tabcolsep}{3.95pt}
\renewcommand{\arraystretch}{1.06}
\setlength{\LTleft}{\fill}
\setlength{\LTright}{\fill}

\begin{longtable}{@{}ccccclllc@{}}
\caption{Definitions of scenarios in the targeted variance-difference sensitivity analysis.}
\label{tab:supp-product-process-scenarios}\\
\toprule
Code & Design & $F$ & \shortstack{$n$/\\group} & $p_{\rm bc}$ & Outcome & Model & Direction & Null value \\
\midrule
\endfirsthead

\multicolumn{9}{c}{\tablename\ \thetable{} -- continued}\\
\toprule
Code & Design & $F$ & \shortstack{$n$/\\group} & $p_{\rm bc}$ & Outcome & Model & Direction & Null value \\
\midrule
\endhead

\midrule
\multicolumn{9}{r}{Continued on next page}\\
\endfoot

\bottomrule
\endlastfoot

S01 & D1 & 2 & 100 & 0.80 & Continuous & realistic & first-factor-contrast & -0.20 \\
S02 & D1 & 2 & 100 & 0.80 & Binary & common-log-odds & common-log-odds & -0.10 \\
S03 & D2 & 2 & 500 & 0.80 & Continuous & realistic & first-factor-contrast & -0.20 \\
S04 & D2 & 2 & 500 & 0.80 & Binary & common-log-odds & common-log-odds & -0.10 \\
S05 & D3 & 5 & 200 & 0.80 & Continuous & realistic & first-factor-contrast & -0.20 \\
S06 & D3 & 5 & 200 & 0.80 & Binary & common-log-odds & common-log-odds & -0.10 \\
S07 & D4 & 5 & 1000 & 0.80 & Continuous & realistic & first-factor-contrast & -0.20 \\
S08 & D4 & 5 & 1000 & 0.80 & Binary & common-log-odds & common-log-odds & -0.10 \\
S09 & D5 & 2 & 100 & 0.95 & Continuous & pure-symmetric & robust-max-ratio & -0.20 \\
S10 & D5 & 2 & 100 & 0.95 & Binary & symmetric-stress & robust-max-ratio & -0.10 \\
S11 & D6 & 2 & 500 & 0.95 & Continuous & pure-symmetric & robust-max-ratio & -0.20 \\
S12 & D6 & 2 & 500 & 0.95 & Binary & symmetric-stress & robust-max-ratio & -0.10 \\

\end{longtable}

\noindent{\small The continuous and binary scenarios use null values $-0.20$ and $-0.10$, respectively, and all tests are one-sided at level 0.025. Designs D1--D4 use $p_{\rm bc}=0.80$, whereas D5 and D6 use $p_{\rm bc}=0.95$.}
\endgroup

\clearpage
\begingroup
\small
\setlength{\tabcolsep}{3.2pt}
\renewcommand{\arraystretch}{0.95}
\setlength{\LTleft}{\fill}
\setlength{\LTright}{\fill}

\begin{longtable}{@{}ccclllcccc@{}}
\caption{Rejection probabilities in the targeted variance-difference sensitivity analysis.}
\label{tab:supp-product-process-sensitivity-results}\\
\toprule
Code & $F$ & \shortstack{$n$/\\group} & Outcome & Model & Analysis &
\shortstack{Mean\\$\widehat\rho_n$} & SIGA-S & SIGA-R & RT \\
\midrule
\endfirsthead

\multicolumn{10}{c}{\tablename\ \thetable{} -- continued}\\
\toprule
Code & $F$ & \shortstack{$n$/\\group} & Outcome & Model & Analysis &
\shortstack{Mean\\$\widehat\rho_n$} & SIGA-S & SIGA-R & RT \\
\midrule
\endhead

\midrule
\multicolumn{10}{r}{Continued on next page}\\
\endfoot

\bottomrule
\endlastfoot

S01 & 2 & 100 & Continuous & realistic & Unadjusted & 1.003 & 2.44 & 2.42 & 2.46 \\
S01 & 2 & 100 & Continuous & realistic & Adjusted   & 1.004 & 2.44 & 2.42 & 2.46 \\
S02 & 2 & 100 & Binary & common-log-odds & Unadjusted & 1.000 & 2.46 & 2.46 & 2.33 \\
S02 & 2 & 100 & Binary & common-log-odds & Adjusted   & 1.000 & 2.39 & 2.39 & 2.33 \\
S03 & 2 & 500 & Continuous & realistic & Unadjusted & 1.003 & 2.48 & 2.46 & 2.48 \\
S03 & 2 & 500 & Continuous & realistic & Adjusted   & 1.003 & 2.49 & 2.47 & 2.48 \\
S04 & 2 & 500 & Binary & common-log-odds & Unadjusted & 1.000 & 2.48 & 2.48 & 2.41 \\
S04 & 2 & 500 & Binary & common-log-odds & Adjusted   & 1.000 & 2.46 & 2.46 & 2.42 \\
S05 & 5 & 200 & Continuous & realistic & Unadjusted & 1.002 & 2.38 & 2.37 & 2.42 \\
S05 & 5 & 200 & Continuous & realistic & Adjusted   & 1.002 & 2.44 & 2.42 & 2.43 \\
S06 & 5 & 200 & Binary & common-log-odds & Unadjusted & 1.000 & 2.36 & 2.36 & 2.37 \\
S06 & 5 & 200 & Binary & common-log-odds & Adjusted   & 1.000 & 2.38 & 2.38 & 2.39 \\
S07 & 5 & 1000 & Continuous & realistic & Unadjusted & 1.002 & 2.47 & 2.46 & 2.47 \\
S07 & 5 & 1000 & Continuous & realistic & Adjusted   & 1.002 & 2.49 & 2.48 & 2.47 \\
S08 & 5 & 1000 & Binary & common-log-odds & Unadjusted & 1.000 & 2.35 & 2.35 & 2.38 \\
S08 & 5 & 1000 & Binary & common-log-odds & Adjusted   & 1.000 & 2.38 & 2.38 & 2.39 \\
S09 & 2 & 100 & Continuous & pure-symmetric & Unadjusted & 1.193 & 2.43 & 1.60 & 1.63 \\
S09 & 2 & 100 & Continuous & pure-symmetric & Adjusted   & 1.193 & 2.43 & 1.60 & 1.64 \\
S10 & 2 & 100 & Binary & symmetric-stress & Unadjusted & 1.115 & 2.19 & 2.02 & 1.77 \\
S10 & 2 & 100 & Binary & symmetric-stress & Adjusted   & 1.115 & 2.23 & 2.02 & 1.80 \\
S11 & 2 & 500 & Continuous & pure-symmetric & Unadjusted & 1.194 & 2.52 & 1.65 & 1.62 \\
S11 & 2 & 500 & Continuous & pure-symmetric & Adjusted   & 1.194 & 2.52 & 1.65 & 1.62 \\
S12 & 2 & 500 & Binary & symmetric-stress & Unadjusted & 1.113 & 2.39 & 2.01 & 1.91 \\
S12 & 2 & 500 & Binary & symmetric-stress & Adjusted   & 1.113 & 2.40 & 2.01 & 1.91 \\

\end{longtable}

\begin{singlespace}
\noindent{\small Rejection probabilities are percentages at a nominal one-sided level of 2.5\%. Codes S01--S08 denote moderately heterogeneous scenarios under $p_{\rm bc}=0.80$, whereas S09--S12 denote deliberately strong treatment-effect-heterogeneity scenarios under $p_{\rm bc}=0.95$.}
\end{singlespace}
\endgroup


\clearpage
\begin{landscape}
\begin{singlespace}
\begingroup
\small
\setlength{\tabcolsep}{1.2pt}
\renewcommand{\arraystretch}{1.06}
\setlength{\LTleft}{\fill}
\setlength{\LTright}{\fill}

\begin{longtable}{@{}
>{\raggedright\arraybackslash}m{62pt}
>{\centering\arraybackslash}m{16pt}
>{\centering\arraybackslash}m{34pt}
>{\raggedright\arraybackslash}m{77pt}
>{\raggedright\arraybackslash}m{64pt}
>{\centering\arraybackslash}m{42pt}
>{\centering\arraybackslash}m{42pt}
>{\centering\arraybackslash}m{42pt}
>{\centering\arraybackslash}m{42pt}
>{\centering\arraybackslash}m{32pt}
>{\centering\arraybackslash}m{68pt}
@{}}

\caption{Complete operating characteristics for the unadjusted analysis in the main simulation studies.}
\label{tab:supp-simulation-unadjusted}\\

\toprule
Outcome & $F$ & \shortstack{$n$/\\group} & Objective & Measure & Effect &
\multicolumn{2}{c}{Repeated-sampling study} &
\multicolumn{3}{c}{Conditional-randomisation study} \\
\cmidrule(lr){7-8}\cmidrule(lr){9-11}
& & & & & &
SIGA-S & RT & SIGA-R & RT & Mean $\widehat\rho_n$ \\
\midrule
\endfirsthead

\multicolumn{11}{c}{\tablename\ \thetable{} -- continued}\\
\toprule
Outcome & $F$ & \shortstack{$n$/\\group} & Objective & Measure & Effect &
\multicolumn{2}{c}{Repeated-sampling study} &
\multicolumn{3}{c}{Conditional-randomisation study} \\
\cmidrule(lr){7-8}\cmidrule(lr){9-11}
& & & & & &
SIGA-S & RT & SIGA-R & RT & Mean $\widehat\rho_n$ \\
\midrule
\endhead

\midrule
\multicolumn{11}{r}{Continued on next page}\\
\endfoot

\bottomrule
\endlastfoot

Continuous & 2 & 100 & Superiority & Type I error & +0.000 & 4.88 & 4.90 & 4.97 & 5.01 & 1.000 \\
Continuous & 2 & 100 & Superiority & Power & +0.430 & 85.20 & 85.17 & 84.92 & 85.17 & 1.010 \\
Continuous & 2 & 100 & Non-inferiority & Type I error & -0.200 & 2.44 & 2.46 & 2.50 & 2.54 & 1.000 \\
Continuous & 2 & 100 & Non-inferiority & Power & +0.230 & 85.14 & 85.06 & 84.85 & 85.06 & 1.010 \\
Continuous & 2 & 100 & Equivalence & Error at $L$ & -0.450 & 4.89 & 4.90 & 5.02 & 5.02 & 1.000/1.038 \\
Continuous & 2 & 100 & Equivalence & Error at $U$ & +0.450 & 4.99 & 4.95 & 4.96 & 4.92 & 1.038/1.000 \\
Continuous & 2 & 100 & Equivalence & Power & +0.000 & 86.96 & 86.91 & 86.37 & 86.48 & 1.011/1.011 \\
\cmidrule(lr){1-11}

Continuous & 2 & 500 & Superiority & Type I error & +0.000 & 4.89 & 4.88 & 5.08 & 5.09 & 1.000 \\
Continuous & 2 & 500 & Superiority & Power & +0.190 & 84.78 & 84.75 & 84.76 & 84.77 & 1.002 \\
Continuous & 2 & 500 & Non-inferiority & Type I error & -0.200 & 2.52 & 2.53 & 2.51 & 2.50 & 1.000 \\
Continuous & 2 & 500 & Non-inferiority & Power & -0.010 & 84.95 & 84.93 & 84.95 & 84.92 & 1.002 \\
Continuous & 2 & 500 & Equivalence & Error at $L$ & -0.450 & 5.02 & 5.01 & 4.98 & 4.98 & 1.000/1.035 \\
Continuous & 2 & 500 & Equivalence & Error at $U$ & +0.450 & 4.93 & 4.92 & 4.98 & 5.02 & 1.035/1.000 \\
Continuous & 2 & 500 & Equivalence & Power & +0.280 & 84.83 & 84.81 & 85.04 & 85.02 & 1.024/1.001 \\
\cmidrule(lr){1-11}

Binary & 2 & 100 & Superiority & Type I error & +0.000 & 4.16 & 4.29 & 4.08 & 4.25 & 1.000 \\
Binary & 2 & 100 & Superiority & Power & +0.200 & 86.11 & 86.21 & 86.00 & 86.32 & 1.011 \\
Binary & 2 & 100 & Non-inferiority & Type I error & -0.100 & 2.44 & 2.28 & 2.44 & 2.25 & 1.000 \\
Binary & 2 & 100 & Non-inferiority & Power & +0.100 & 84.08 & 83.85 & 83.90 & 83.85 & 1.010 \\
Binary & 2 & 100 & Equivalence & Error at $L$ & -0.210 & 4.77 & 4.68 & 4.74 & 4.65 & 1.000/1.035 \\
Binary & 2 & 100 & Equivalence & Error at $U$ & +0.210 & 4.90 & 4.89 & 4.97 & 5.00 & 1.042/1.000 \\
Binary & 2 & 100 & Equivalence & Power & +0.000 & 83.48 & 83.54 & 83.37 & 83.61 & 1.010/1.010 \\
\cmidrule(lr){1-11}

Binary & 2 & 500 & Superiority & Type I error & +0.000 & 4.55 & 4.54 & 4.68 & 4.67 & 1.000 \\
Binary & 2 & 500 & Superiority & Power & +0.100 & 91.07 & 91.06 & 91.09 & 91.19 & 1.002 \\
Binary & 2 & 500 & Non-inferiority & Type I error & -0.100 & 2.41 & 2.33 & 2.54 & 2.46 & 1.000 \\
Binary & 2 & 500 & Non-inferiority & Power & +0.000 & 89.68 & 89.69 & 89.79 & 89.87 & 1.002 \\
Binary & 2 & 500 & Equivalence & Error at $L$ & -0.100 & 5.04 & 4.89 & 4.96 & 4.79 & 1.000/1.008 \\
Binary & 2 & 500 & Equivalence & Error at $U$ & +0.100 & 5.05 & 4.95 & 5.04 & 4.97 & 1.009/1.000 \\
Binary & 2 & 500 & Equivalence & Power & +0.000 & 89.04 & 89.06 & 88.94 & 88.95 & 1.002/1.002 \\
\cmidrule(lr){1-11}

Continuous & 5 & 200 & Superiority & Type I error & +0.000 & 4.89 & 4.95 & 4.94 & 5.04 & 1.000 \\
Continuous & 5 & 200 & Superiority & Power & +0.300 & 84.18 & 84.32 & 84.13 & 84.37 & 1.004 \\
Continuous & 5 & 200 & Non-inferiority & Type I error & -0.200 & 2.48 & 2.52 & 2.45 & 2.51 & 1.000 \\
Continuous & 5 & 200 & Non-inferiority & Power & +0.100 & 84.28 & 84.36 & 84.26 & 84.46 & 1.004 \\
Continuous & 5 & 200 & Equivalence & Error at $L$ & -0.450 & 4.90 & 4.96 & 4.82 & 4.89 & 1.000/1.030 \\
Continuous & 5 & 200 & Equivalence & Error at $U$ & +0.450 & 4.89 & 4.95 & 4.94 & 5.02 & 1.030/1.000 \\
Continuous & 5 & 200 & Equivalence & Power & +0.180 & 84.67 & 84.75 & 84.74 & 84.90 & 1.016/1.003 \\
\cmidrule(lr){1-11}

Continuous & 5 & 1000 & Superiority & Type I error & +0.000 & 5.02 & 5.01 & 5.12 & 5.11 & 1.000 \\
Continuous & 5 & 1000 & Superiority & Power & +0.135 & 85.11 & 85.13 & 85.22 & 85.23 & 1.001 \\
Continuous & 5 & 1000 & Non-inferiority & Type I error & -0.200 & 2.50 & 2.50 & 2.56 & 2.57 & 1.000 \\
Continuous & 5 & 1000 & Non-inferiority & Power & -0.065 & 85.43 & 85.38 & 85.15 & 85.16 & 1.001 \\
Continuous & 5 & 1000 & Equivalence & Error at $L$ & -0.450 & 4.95 & 4.99 & 4.95 & 4.95 & 1.000/1.026 \\
Continuous & 5 & 1000 & Equivalence & Error at $U$ & +0.450 & 5.09 & 5.12 & 4.87 & 4.92 & 1.026/1.000 \\
Continuous & 5 & 1000 & Equivalence & Power & +0.330 & 84.75 & 84.78 & 84.75 & 84.80 & 1.020/1.001 \\
\cmidrule(lr){1-11}

Binary & 5 & 200 & Superiority & Type I error & +0.000 & 4.58 & 4.58 & 4.47 & 4.45 & 1.000 \\
Binary & 5 & 200 & Superiority & Power & +0.140 & 83.99 & 84.03 & 83.93 & 84.06 & 1.004 \\
Binary & 5 & 200 & Non-inferiority & Type I error & -0.100 & 2.38 & 2.38 & 2.31 & 2.30 & 1.000 \\
Binary & 5 & 200 & Non-inferiority & Power & +0.050 & 87.43 & 87.33 & 87.20 & 87.19 & 1.004 \\
Binary & 5 & 200 & Equivalence & Error at $L$ & -0.150 & 4.69 & 4.79 & 4.65 & 4.76 & 1.000/1.015 \\
Binary & 5 & 200 & Equivalence & Error at $U$ & +0.150 & 4.88 & 4.88 & 5.03 & 4.99 & 1.017/1.000 \\
Binary & 5 & 200 & Equivalence & Power & +0.000 & 84.89 & 84.85 & 84.64 & 84.74 & 1.004/1.004 \\
\cmidrule(lr){1-11}

Binary & 5 & 1000 & Superiority & Type I error & +0.000 & 4.67 & 4.70 & 4.64 & 4.66 & 1.000 \\
Binary & 5 & 1000 & Superiority & Power & +0.070 & 90.05 & 90.03 & 90.23 & 90.19 & 1.001 \\
Binary & 5 & 1000 & Non-inferiority & Type I error & -0.100 & 2.45 & 2.47 & 2.40 & 2.42 & 1.000 \\
Binary & 5 & 1000 & Non-inferiority & Power & -0.030 & 89.27 & 89.15 & 89.35 & 89.24 & 1.001 \\
Binary & 5 & 1000 & Equivalence & Error at $L$ & -0.100 & 4.64 & 4.71 & 4.79 & 4.84 & 1.000/1.006 \\
Binary & 5 & 1000 & Equivalence & Error at $U$ & +0.100 & 4.98 & 4.95 & 5.00 & 4.97 & 1.007/1.000 \\
Binary & 5 & 1000 & Equivalence & Power & +0.040 & 87.22 & 87.22 & 87.38 & 87.35 & 1.003/1.001 \\

\end{longtable}

\noindent{\small Rejection probabilities are percentages. The repeated-sampling and conditional-randomisation studies used independent outer simulations, so the two RT columns need not be numerically identical. With 100,000 outer trials the Monte Carlo standard error of a rejection probability is 0.07 points at 5\% and 0.11 points at 85\%. For equivalence, $L$ and $U$ denote the lower and upper limits, and two variance ratios are reported in lower-limit/upper-limit order.}

\endgroup
\end{singlespace}
\end{landscape}


\clearpage
\begin{landscape}
\begin{singlespace}
\begingroup
\small
\setlength{\tabcolsep}{1.2pt}
\renewcommand{\arraystretch}{1.06}
\setlength{\LTleft}{\fill}
\setlength{\LTright}{\fill}

\begin{longtable}{@{}
>{\raggedright\arraybackslash}m{62pt}
>{\centering\arraybackslash}m{16pt}
>{\centering\arraybackslash}m{34pt}
>{\raggedright\arraybackslash}m{77pt}
>{\raggedright\arraybackslash}m{64pt}
>{\centering\arraybackslash}m{42pt}
>{\centering\arraybackslash}m{42pt}
>{\centering\arraybackslash}m{42pt}
>{\centering\arraybackslash}m{42pt}
>{\centering\arraybackslash}m{32pt}
>{\centering\arraybackslash}m{68pt}
@{}}

\caption{Complete operating characteristics for the baseline-adjusted analysis in the main simulation studies.}
\label{tab:supp-simulation-adjusted}\\

\toprule
Outcome & $F$ & \shortstack{$n$/\\group} & Objective & Measure & Effect &
\multicolumn{2}{c}{Repeated-sampling study} &
\multicolumn{3}{c}{Conditional-randomisation study} \\
\cmidrule(lr){7-8}\cmidrule(lr){9-11}
& & & & & &
SIGA-S & RT & SIGA-R & RT & Mean $\widehat\rho_n$ \\
\midrule
\endfirsthead

\multicolumn{11}{c}{\tablename\ \thetable{} -- continued}\\
\toprule
Outcome & $F$ & \shortstack{$n$/\\group} & Objective & Measure & Effect &
\multicolumn{2}{c}{Repeated-sampling study} &
\multicolumn{3}{c}{Conditional-randomisation study} \\
\cmidrule(lr){7-8}\cmidrule(lr){9-11}
& & & & & &
SIGA-S & RT & SIGA-R & RT & Mean $\widehat\rho_n$ \\
\midrule
\endhead

\midrule
\multicolumn{11}{r}{Continued on next page}\\
\endfoot

\bottomrule
\endlastfoot

Continuous & 2 & 100 & Superiority & Type I error & +0.000 & 4.87 & 4.88 & 4.93 & 4.98 & 1.000 \\
Continuous & 2 & 100 & Superiority & Power & +0.430 & 85.50 & 85.44 & 85.20 & 85.40 & 1.010 \\
Continuous & 2 & 100 & Non-inferiority & Type I error & -0.200 & 2.44 & 2.44 & 2.51 & 2.51 & 1.000 \\
Continuous & 2 & 100 & Non-inferiority & Power & +0.230 & 85.41 & 85.37 & 85.08 & 85.29 & 1.010 \\
Continuous & 2 & 100 & Equivalence & Error at $L$ & -0.450 & 4.92 & 4.94 & 5.01 & 5.00 & 1.000/1.038 \\
Continuous & 2 & 100 & Equivalence & Error at $U$ & +0.450 & 4.96 & 4.94 & 4.98 & 4.94 & 1.038/1.000 \\
Continuous & 2 & 100 & Equivalence & Power & +0.000 & 87.28 & 87.17 & 86.68 & 86.83 & 1.011/1.011 \\
\cmidrule(lr){1-11}

Continuous & 2 & 500 & Superiority & Type I error & +0.000 & 4.87 & 4.88 & 5.08 & 5.07 & 1.000 \\
Continuous & 2 & 500 & Superiority & Power & +0.190 & 84.80 & 84.78 & 84.87 & 84.85 & 1.002 \\
Continuous & 2 & 500 & Non-inferiority & Type I error & -0.200 & 2.52 & 2.54 & 2.50 & 2.51 & 1.000 \\
Continuous & 2 & 500 & Non-inferiority & Power & -0.010 & 84.99 & 84.95 & 84.97 & 84.96 & 1.002 \\
Continuous & 2 & 500 & Equivalence & Error at $L$ & -0.450 & 5.04 & 5.00 & 4.96 & 4.97 & 1.000/1.035 \\
Continuous & 2 & 500 & Equivalence & Error at $U$ & +0.450 & 4.90 & 4.92 & 4.98 & 5.00 & 1.035/1.000 \\
Continuous & 2 & 500 & Equivalence & Power & +0.280 & 84.87 & 84.85 & 85.09 & 85.04 & 1.025/1.001 \\
\cmidrule(lr){1-11}

Binary & 2 & 100 & Superiority & Type I error & +0.000 & 5.10 & 4.92 & 5.02 & 4.81 & 1.000 \\
Binary & 2 & 100 & Superiority & Power & +0.200 & 87.92 & 87.50 & 87.83 & 87.62 & 1.011 \\
Binary & 2 & 100 & Non-inferiority & Type I error & -0.100 & 2.37 & 2.32 & 2.36 & 2.28 & 1.000 \\
Binary & 2 & 100 & Non-inferiority & Power & +0.100 & 84.13 & 84.06 & 83.90 & 84.00 & 1.010 \\
Binary & 2 & 100 & Equivalence & Error at $L$ & -0.210 & 4.76 & 4.66 & 4.73 & 4.62 & 1.000/1.035 \\
Binary & 2 & 100 & Equivalence & Error at $U$ & +0.210 & 4.88 & 4.86 & 4.96 & 4.99 & 1.042/1.000 \\
Binary & 2 & 100 & Equivalence & Power & +0.000 & 83.64 & 83.58 & 83.48 & 83.65 & 1.010/1.010 \\
\cmidrule(lr){1-11}

Binary & 2 & 500 & Superiority & Type I error & +0.000 & 4.85 & 4.83 & 5.00 & 4.97 & 1.000 \\
Binary & 2 & 500 & Superiority & Power & +0.100 & 91.62 & 91.51 & 91.66 & 91.61 & 1.002 \\
Binary & 2 & 500 & Non-inferiority & Type I error & -0.100 & 2.39 & 2.35 & 2.51 & 2.47 & 1.000 \\
Binary & 2 & 500 & Non-inferiority & Power & +0.000 & 89.72 & 89.73 & 89.83 & 89.93 & 1.002 \\
Binary & 2 & 500 & Equivalence & Error at $L$ & -0.100 & 5.00 & 4.92 & 4.91 & 4.82 & 1.000/1.008 \\
Binary & 2 & 500 & Equivalence & Error at $U$ & +0.100 & 5.06 & 5.00 & 5.03 & 5.00 & 1.009/1.000 \\
Binary & 2 & 500 & Equivalence & Power & +0.000 & 89.07 & 89.11 & 88.98 & 89.00 & 1.002/1.002 \\
\cmidrule(lr){1-11}

Continuous & 5 & 200 & Superiority & Type I error & +0.000 & 4.99 & 4.97 & 5.00 & 4.99 & 1.000 \\
Continuous & 5 & 200 & Superiority & Power & +0.300 & 84.71 & 84.65 & 84.73 & 84.79 & 1.004 \\
Continuous & 5 & 200 & Non-inferiority & Type I error & -0.200 & 2.53 & 2.52 & 2.51 & 2.54 & 1.000 \\
Continuous & 5 & 200 & Non-inferiority & Power & +0.100 & 84.79 & 84.75 & 84.79 & 84.81 & 1.004 \\
Continuous & 5 & 200 & Equivalence & Error at $L$ & -0.450 & 4.98 & 4.99 & 4.92 & 4.91 & 1.000/1.030 \\
Continuous & 5 & 200 & Equivalence & Error at $U$ & +0.450 & 4.95 & 4.97 & 5.07 & 5.08 & 1.030/1.000 \\
Continuous & 5 & 200 & Equivalence & Power & +0.180 & 85.14 & 85.06 & 85.11 & 85.09 & 1.016/1.003 \\
\cmidrule(lr){1-11}

Continuous & 5 & 1000 & Superiority & Type I error & +0.000 & 5.01 & 5.00 & 5.11 & 5.11 & 1.000 \\
Continuous & 5 & 1000 & Superiority & Power & +0.135 & 85.20 & 85.17 & 85.37 & 85.33 & 1.001 \\
Continuous & 5 & 1000 & Non-inferiority & Type I error & -0.200 & 2.54 & 2.52 & 2.57 & 2.59 & 1.000 \\
Continuous & 5 & 1000 & Non-inferiority & Power & -0.065 & 85.52 & 85.44 & 85.28 & 85.27 & 1.001 \\
Continuous & 5 & 1000 & Equivalence & Error at $L$ & -0.450 & 4.98 & 4.98 & 4.97 & 4.94 & 1.000/1.026 \\
Continuous & 5 & 1000 & Equivalence & Error at $U$ & +0.450 & 5.12 & 5.11 & 4.92 & 4.93 & 1.026/1.000 \\
Continuous & 5 & 1000 & Equivalence & Power & +0.330 & 84.84 & 84.82 & 84.87 & 84.83 & 1.020/1.001 \\
\cmidrule(lr){1-11}

Binary & 5 & 200 & Superiority & Type I error & +0.000 & 5.08 & 5.10 & 4.97 & 5.00 & 1.000 \\
Binary & 5 & 200 & Superiority & Power & +0.140 & 85.32 & 85.24 & 85.27 & 85.24 & 1.004 \\
Binary & 5 & 200 & Non-inferiority & Type I error & -0.100 & 2.41 & 2.42 & 2.33 & 2.33 & 1.000 \\
Binary & 5 & 200 & Non-inferiority & Power & +0.050 & 87.50 & 87.43 & 87.30 & 87.32 & 1.004 \\
Binary & 5 & 200 & Equivalence & Error at $L$ & -0.150 & 4.84 & 4.85 & 4.82 & 4.83 & 1.000/1.015 \\
Binary & 5 & 200 & Equivalence & Error at $U$ & +0.150 & 4.90 & 4.91 & 5.06 & 5.03 & 1.017/1.000 \\
Binary & 5 & 200 & Equivalence & Power & +0.000 & 85.06 & 84.97 & 84.90 & 84.88 & 1.004/1.004 \\
\cmidrule(lr){1-11}

Binary & 5 & 1000 & Superiority & Type I error & +0.000 & 4.92 & 4.92 & 4.91 & 4.88 & 1.000 \\
Binary & 5 & 1000 & Superiority & Power & +0.070 & 90.45 & 90.42 & 90.61 & 90.56 & 1.001 \\
Binary & 5 & 1000 & Non-inferiority & Type I error & -0.100 & 2.48 & 2.50 & 2.43 & 2.44 & 1.000 \\
Binary & 5 & 1000 & Non-inferiority & Power & -0.030 & 89.22 & 89.19 & 89.30 & 89.28 & 1.001 \\
Binary & 5 & 1000 & Equivalence & Error at $L$ & -0.100 & 4.72 & 4.75 & 4.86 & 4.86 & 1.000/1.006 \\
Binary & 5 & 1000 & Equivalence & Error at $U$ & +0.100 & 4.97 & 4.97 & 5.00 & 4.99 & 1.007/1.000 \\
Binary & 5 & 1000 & Equivalence & Power & +0.040 & 87.25 & 87.26 & 87.41 & 87.43 & 1.003/1.001 \\

\end{longtable}

\noindent{\small Rejection probabilities are percentages. The repeated-sampling and conditional-randomisation studies used independent outer simulations, so the two RT columns need not be numerically identical. With 100,000 outer trials the Monte Carlo standard error of a rejection probability is 0.07 points at 5\% and 0.11 points at 85\%. For equivalence, $L$ and $U$ denote the lower and upper limits, and two variance ratios are reported in lower-limit/upper-limit order.}

\endgroup
\end{singlespace}
\end{landscape}

\clearpage
\begin{singlespace}
\begingroup
\small
\setlength{\tabcolsep}{3.2pt}
\renewcommand{\arraystretch}{1.06}
\setlength{\LTleft}{\fill}
\setlength{\LTright}{\fill}

\begin{longtable}{@{}lcrllrrr@{}}
\caption{Projected single-core computation times for 100,000 analyses in the simulation settings.}
\label{tab:supp-computation-times}\\
\toprule
Outcome & $F$ & $n$/group & Analysis & Objective & SIGA-S (min) & SIGA-R (min) & RT (min) \\
\midrule
\endfirsthead

\multicolumn{8}{c}{\tablename\ \thetable{} -- continued}\\
\toprule
Outcome & $F$ & $n$/group & Analysis & Objective & SIGA-S (min) & SIGA-R (min) & RT (min) \\
\midrule
\endhead

\midrule
\multicolumn{8}{r}{Continued on next page}\\
\endfoot

\bottomrule
\endlastfoot

Binary & 2 & 100 & Adjusted & Equivalence & 0.129 & 0.152 & 143.44 \\
Binary & 2 & 100 & Unadjusted & Equivalence & 0.134 & 0.152 & 142.98 \\
Binary & 2 & 100 & Adjusted & Non-inferiority & 0.067 & 0.080 & 140.35 \\
Binary & 2 & 100 & Unadjusted & Non-inferiority & 0.069 & 0.078 & 139.87 \\
Binary & 2 & 100 & Adjusted & Superiority & 0.066 & 0.079 & 135.65 \\
Binary & 2 & 100 & Unadjusted & Superiority & 0.160 & 0.170 & 134.63 \\

Binary & 2 & 500 & Adjusted & Equivalence & 0.261 & 0.286 & 690.28 \\
Binary & 2 & 500 & Unadjusted & Equivalence & 0.259 & 0.290 & 678.54 \\
Binary & 2 & 500 & Adjusted & Non-inferiority & 0.134 & 0.145 & 672.37 \\
Binary & 2 & 500 & Unadjusted & Non-inferiority & 0.133 & 0.148 & 666.31 \\
Binary & 2 & 500 & Adjusted & Superiority & 0.132 & 0.145 & 659.76 \\
Binary & 2 & 500 & Unadjusted & Superiority & 0.423 & 0.466 & 649.35 \\

Binary & 5 & 200 & Adjusted & Equivalence & 0.243 & 0.270 & 395.85 \\
Binary & 5 & 200 & Unadjusted & Equivalence & 0.242 & 0.269 & 400.57 \\
Binary & 5 & 200 & Adjusted & Non-inferiority & 0.124 & 0.138 & 385.80 \\
Binary & 5 & 200 & Unadjusted & Non-inferiority & 0.123 & 0.138 & 387.52 \\
Binary & 5 & 200 & Adjusted & Superiority & 0.121 & 0.138 & 385.39 \\
Binary & 5 & 200 & Unadjusted & Superiority & 0.284 & 0.288 & 388.74 \\

Binary & 5 & 1000 & Adjusted & Equivalence & 0.770 & 0.815 & 2021.50 \\
Binary & 5 & 1000 & Unadjusted & Equivalence & 0.761 & 0.788 & 2032.31 \\
Binary & 5 & 1000 & Adjusted & Non-inferiority & 0.376 & 0.410 & 1884.43 \\
Binary & 5 & 1000 & Unadjusted & Non-inferiority & 0.372 & 0.404 & 1899.41 \\
Binary & 5 & 1000 & Adjusted & Superiority & 0.389 & 0.401 & 1937.07 \\
Binary & 5 & 1000 & Unadjusted & Superiority & 0.945 & 0.981 & 1967.33 \\

Continuous & 2 & 100 & Adjusted & Equivalence & 0.127 & 0.151 & 140.20 \\
Continuous & 2 & 100 & Unadjusted & Equivalence & 0.130 & 0.153 & 138.30 \\
Continuous & 2 & 100 & Adjusted & Non-inferiority & 0.066 & 0.077 & 133.26 \\
Continuous & 2 & 100 & Unadjusted & Non-inferiority & 0.064 & 0.080 & 130.52 \\
Continuous & 2 & 100 & Adjusted & Superiority & 0.065 & 0.078 & 130.83 \\
Continuous & 2 & 100 & Unadjusted & Superiority & 0.066 & 0.079 & 132.30 \\

Continuous & 2 & 500 & Adjusted & Equivalence & 0.264 & 0.286 & 682.17 \\
Continuous & 2 & 500 & Unadjusted & Equivalence & 0.260 & 0.283 & 677.35 \\
Continuous & 2 & 500 & Adjusted & Non-inferiority & 0.130 & 0.145 & 642.52 \\
Continuous & 2 & 500 & Unadjusted & Non-inferiority & 0.131 & 0.143 & 641.98 \\
Continuous & 2 & 500 & Adjusted & Superiority & 0.131 & 0.144 & 644.17 \\
Continuous & 2 & 500 & Unadjusted & Superiority & 0.131 & 0.146 & 648.56 \\

Continuous & 5 & 200 & Adjusted & Equivalence & 0.241 & 0.272 & 385.46 \\
Continuous & 5 & 200 & Unadjusted & Equivalence & 0.242 & 0.271 & 384.39 \\
Continuous & 5 & 200 & Adjusted & Non-inferiority & 0.123 & 0.140 & 373.74 \\
Continuous & 5 & 200 & Unadjusted & Non-inferiority & 0.121 & 0.137 & 375.70 \\
Continuous & 5 & 200 & Adjusted & Superiority & 0.123 & 0.138 & 383.56 \\
Continuous & 5 & 200 & Unadjusted & Superiority & 0.124 & 0.136 & 388.22 \\

Continuous & 5 & 1000 & Adjusted & Equivalence & 0.776 & 0.796 & 2005.37 \\
Continuous & 5 & 1000 & Unadjusted & Equivalence & 0.764 & 0.765 & 2018.74 \\
Continuous & 5 & 1000 & Adjusted & Non-inferiority & 0.393 & 0.402 & 1909.06 \\
Continuous & 5 & 1000 & Unadjusted & Non-inferiority & 0.379 & 0.394 & 1938.11 \\
Continuous & 5 & 1000 & Adjusted & Superiority & 0.384 & 0.402 & 1917.67 \\
Continuous & 5 & 1000 & Unadjusted & Superiority & 0.370 & 0.412 & 1882.48 \\

\end{longtable}

\endgroup
\end{singlespace}

\begin{table}[!htbp]
\centering
\caption{Single-core computation times for one reusable randomisation-design covariance estimation.}
\label{tab:supp-calibration-times}
\begin{threeparttable}
\begin{singlespace}
\begin{tabular}{rrrrrr}
\toprule
Factors & $n$/group & Total $n$ & Replicates & SIGA-S mean (s) & SIGA-R mean (s) \\
\midrule
2 & 100 & 200 & 100,000 & 3.989 & 8.933 \\
2 & 500 & 1000 & 100,000 & 18.627 & 44.052 \\
5 & 200 & 400 & 100,000 & 12.146 & 29.103 \\
5 & 1000 & 2000 & 100,000 & 59.791 & 140.779 \\
\bottomrule
\end{tabular}
\begin{tablenotes}[flushleft]
\small
\item SIGA-S estimates covariance using one allocation sequence per replicate. SIGA-R uses three sequences per replicate to estimate both the imbalance covariance and the covariance of treatment-sign products. Each mean is based on 3 independent benchmark repetitions.
\end{tablenotes}
\end{singlespace}
\end{threeparttable}
\end{table}


\clearpage
\subsection{Covariance ratios and independent evaluation}\label{sec:supp-eigen-diagnostics}\label{sec:supp-independent-directions}
The supplemental sensitivity study also recorded the minimum and maximum generalised covariance ratios. The calculation restricts the effect direction to $\bm\pi^\top v=0$; equivalently it uses the matrix in Proposition~\ref{prop:supp-sign}. Table~\ref{tab:supp-eigen-diagnostics} reports these estimates based on 100,000 covariance-estimation replicates per design. Appendix~I evaluates the finite-sample covariance for one to three factors by a deterministic recursion, without Monte Carlo error.
\begin{table}[!htbp]
\centering
\caption{Estimated smallest and largest covariance ratios on $\bm\pi^\top v=0$.}
\label{tab:supp-eigen-diagnostics}
\begin{threeparttable}
\begin{singlespace}
\begin{tabular}{rrrrr}
\toprule
Factors & Total $n$ & $p_{\rm bc}$ & Minimum & Maximum\\
\midrule
2 & 200 & .80 & 1.046187 & 1.061922\\
2 & 1000 & .80 & 1.043654 & 1.056350\\
5 & 400 & .80 & .981641 & 1.049601\\
5 & 2000 & .80 & .981502 & 1.043298\\
2 & 200 & .95 & 1.063163 & 1.197081\\
2 & 1000 & .95 & 1.063146 & 1.196307\\
\bottomrule
\end{tabular}
\begin{tablenotes}[flushleft]\small
\item All six designs use independent binary factors with the prevalences specified for the main numerical study. Values are taken from the released code \citep{KojimaSIGA2026}; covariance-estimation randomness and finite-trial-size approximation remain present.
\end{tablenotes}
\end{singlespace}
\end{threeparttable}
\end{table}
The five-factor entries are same-simulation minima. For a fixed known $\bm\pi$, the sample covariance is unbiased for its finite-$n$ target. Since the smallest eigenvalue is a concave function of a symmetric matrix, selecting the minimum from that estimated matrix can lower its expected value. We therefore evaluated selected directions with independent simulations as described below. The constraint is $\bm\pi^\top v=0$; for unequal profile probabilities, an interaction vector such as $(1,-1,-1,1)^\top$ generally needs centring by its $\bm\pi$-weighted mean.


Additional assignment-only simulations evaluated the two five-factor designs with $p_{\rm bc}=.80$ and total sizes 400 and 2000. The factor probabilities were $(.50,.40,.30,.20,.10)$ and the allocation rule was unchanged. Each setting used 200,000 independent triples to estimate the covariance and select the smallest-ratio direction, followed by 100,000 new independent triples to evaluate that direction. A third setting used independent fair assignments, implemented by $p_{\rm bc}=.50$, and total size 400. In this control setting the true product covariance is exactly $\bm\Pi$ for every $n$. These calculations generated profiles and assignments only; no trial outcomes or randomisation-test rejection indicators were simulated.

Let $\overline\Psi_n=\Cov(\bm G_n^{(01)})$ for a fixed trial size $n$. The selection matrix is the average of the centred sample covariance matrices of $\bm G_n^{(01)}$ and $\bm G_n^{(02)}$, using divisor $B_{\rm sel}-1$ in each. We select a minimising direction $\widehat v$ under $\bm\pi^\top\widehat v=0$ and normalise it so that $\widehat v^\top\bm\Pi\widehat v=1$. In evaluation replicate $r$, define
\[
 A_r=\tfrac12\left[\{\widehat v^\top\bm G_{n,r}^{(01)}\}^2+
                        \{\widehat v^\top\bm G_{n,r}^{(02)}\}^2\right].
\]
Conditional on the selection simulations, $\E(A_r)=\widehat v^\top\overline\Psi_n\widehat v$ because the product vectors have exactly zero means. The estimate is $\bar A$ and its Monte Carlo standard error is the sample standard deviation of the $A_r$ divided by $\sqrt{B_{\rm eval}}$. Both products from one triple enter the same $A_r$, so their dependence is included in this standard error. The reported intervals are normal Monte Carlo intervals, $\bar A\pm1.96\,\widehat{\operatorname{se}}(\bar A)$, conditional on the selected direction. They quantify simulation uncertainty in a fixed finite-$n$ ratio, not uncertainty in the smallest eigenvalue of $\bm\Psi_D$.

\begin{table}[!htbp]
\centering
\caption{Independent evaluation of selected smallest-ratio directions for five factors.}
\label{tab:supp-independent-directions}
\begin{threeparttable}
\begin{singlespace}
\small
\setlength{\tabcolsep}{4pt}
\begin{tabular}{@{}lrrrrl@{}}
\toprule
Allocation & Total $n$ & \shortstack{Selection\\minimum} & \shortstack{Independent\\ratio} & \shortstack{Monte Carlo\\standard error} & 95\% interval\\
\midrule
Minimisation & 400 & 0.986416 & 1.004618 & 0.003259 & [0.998231, 1.011005]\\
Minimisation & 2000 & 0.987215 & 1.004979 & 0.003192 & [0.998722, 1.011235]\\
\shortstack[l]{Independent\\assignment} & 400 & 0.984303 & 0.998805 & 0.003183 & [0.992567, 1.005043]\\
\bottomrule
\end{tabular}
\begin{tablenotes}[flushleft]\small
\item Each row uses 200,000 selection triples and a separate 100,000 evaluation triples. Both minimisation rows have $p_{\rm bc}=.80$. The independent-assignment row has exactly unit true ratio for every admissible direction. Seeds, the code generating the range-method assignments, the selected directions and covariance summaries accompany the manuscript.
\end{tablenotes}
\end{singlespace}
\end{threeparttable}
\end{table}

The independent estimates in both minimisation settings exceeded one and their intervals included one. Thus these calculations provided no evidence of a ratio below one along the selected directions. The control also produced a selection minimum below one, while its independent estimate was consistent with its known value of one. This comparison shows why a minimum computed and evaluated using the same covariance estimate cannot establish a negative direction. The results concern the selected finite-size directions and do not establish nonnegativity for all directions, other profile laws, or the limiting minimisation covariance; Appendix~I shows by exact enumeration at $n=3$, and by deterministic recursion for larger $n$, that the finite-sample ratio can be below one for three factors.


\section{Examples of other allocation rules}\label{sec:supp-other-designs}
This appendix verifies the common inference condition for several allocation rules. The examples use their own original assignment laws when regenerating treatment sequences. Their proofs are independent of the range-method recurrence proof. Corollary~\ref{cor:supp-stratified-sign} explains the common nonnegative sign for the designs that allocate independently within strata; the separate proofs below also verify the required joint limits and moments.

\subsection{Independent assignment and stratified permuted blocks}\label{sec:supp-independent}\label{sec:supp-block-extension}
\begin{proposition}\label{prop:supp-independent}
If each assignment is an independent fair treatment sign, independent also of the profiles, Assumption~\ref{ass:supp-general-allocation} holds with $\Sigma_D=\Psi_D=\Pi$.
\end{proposition}
\begin{proof}
Stack the centred profile indicator and $h(S_i)\prod_{a\in B}Z_i^{(a)}$ for every nonempty subset $B$ of the three sequences. The stacked vectors are independent and identically distributed and bounded. Their means are zero. Distinct nonempty subsets have zero cross-covariance because some independent sign occurs to an odd power; the same argument gives zero covariance with the centred profile vector. Each nonempty block has covariance $\Pi$. The multivariate central limit theorem gives the joint Gaussian limit, including the triple product. The fourth-moment expansion of sums of centred independent bounded vectors is $O(n^2)$, and second moments converge (in this case they equal their limits at every $n$).
\end{proof}

Next consider balanced permuted blocks formed separately within each stratum.
\begin{proposition}\label{prop:supp-blocks}
Within stratum $s$, allocate successive arrivals in blocks of fixed even size $\ell_s\ge2$, choosing uniformly a balanced sign vector in each block. The vectors are independent across blocks, strata and regenerated allocation sequences. The block boundaries are the same in all sequences, as determined by the ordered profiles; the final block may be incomplete. For independent and identically distributed profiles with fixed $J$ and positive $\pi_s$, Assumption~\ref{ass:supp-general-allocation} holds with
\begin{equation}\label{eq:supp-block-covariance}
 \Sigma_D=0,\qquad
 \Psi_D=\diag\{\pi_s\ell_s/(\ell_s-1):s=1,\ldots,J\}.
\end{equation}
\end{proposition}
\begin{proof}
For a complete block of size $\ell$, let $Z^{(a)}=(Z_1^{(a)},\ldots,Z_\ell^{(a)})$ be the balanced treatment vector in sequence $a$. Label symmetry gives $E Z_j^{(a)}=0$. Exchangeability within the block and $\sum_jZ_j^{(a)}=0$ give
\[
 E[Z_j^{(a)}Z_t^{(a)}]=-\frac1{\ell-1}\quad(j\ne t).
\]
For distinct sequences $a,b$, set $B_{ab}=\sum_{j=1}^\ell Z_j^{(a)}Z_j^{(b)}$. Independence of the two block vectors gives $E B_{ab}=0$ and
\begin{equation}\label{eq:supp-block-pair-variance}
 E B_{ab}^2=\ell+\ell(\ell-1)\left(\frac1{\ell-1}\right)^2
             =\frac{\ell^2}{\ell-1}.
\end{equation}
For the triple product $B_{012}=\sum_jZ_j^{(0)}Z_j^{(1)}Z_j^{(2)}$, the same calculation yields
\[
 E B_{012}=0,\qquad E B_{012}^2=\ell-\frac{\ell}{(\ell-1)^2}.
\]
Any cross moment between different pairs is zero: globally flip the signs of a sequence appearing in exactly one of the two pairs. A pair and the triple product are likewise uncorrelated. This argument preserves each balanced block law.

Conditional on the complete profile sequence, there are $q_{n,s}=\lfloor N_{n,s}/\ell_s\rfloor$ complete blocks in stratum $s$. All block vectors consisting of the three pair rewards and the triple reward are conditionally independent across blocks and strata, centred and uniformly bounded. Since $q_{n,s}/n\to\pi_s/\ell_s$, their conditional covariance matrices after normalisation by $n$ converge to deterministic matrices. The Lindeberg condition is automatic because the block sizes are fixed. The multivariate conditional central limit theorem therefore gives pair-block covariance $\pi_s\ell_s/(\ell_s-1)$ and triple-block covariance $\pi_s\{1-(\ell_s-1)^{-2}\}$. Each incomplete block contributes a bounded remainder, and $|D_{n,s}^{(a)}|\le\ell_s$ because completed blocks are balanced. Thus all normalised imbalance vectors tend to zero.

To include $R_n$ jointly, the conditional characteristic function of the block sums given the profiles converges in probability to that of the stated deterministic Gaussian law. It is bounded by one, so the convergence is also in $L^1$. Multiplication by the characteristic function factor for $R_n$, followed by expectation, and the multinomial central limit theorem show that the limiting profile vector is independent of all product blocks. This proves the joint limit in Assumption~\ref{ass:supp-general-allocation} and also proves tightness of the normalised triple product.

Conditional on the profiles, the fourth-moment expansion of any sum of these independent centred bounded block vectors is bounded by $C(n+1)^2$, with a constant depending only on the fixed block sizes and $J$. The incomplete-block remainders are bounded. The multinomial profile vector has the same fourth-moment order. Consequently the normalised fourth moments are uniformly bounded. They imply uniform integrability of the normalised quadratic products and hence convergence of their expectations. This proves all the moment requirements in the assumption.
\end{proof}

The distinction from within-stratum permutation is important. Here the original consecutive blocks are preserved on every regenerated sequence. The group of permutations of all treatment labels within a stratum, used for example in Section~4.3 of \citet{BugniCanayShaikh2018}, generally permits vectors that violate these original block boundaries. Formula~(\ref{eq:supp-block-covariance}) pertains to regeneration under the fixed blocks, not to that larger permutation group. The companion script exhaustively checks the finite-block covariance identities for sizes two, four and six using exact rational arithmetic; the limit argument is the proof above, not a simulation claim.

At the average-effect null value $\Delta=b$, Corollary~\ref{cor:supp-variance-equality} and (\ref{eq:supp-block-covariance}) give
\begin{equation}\label{eq:supp-block-gap}
 v_{\rm R}(b)-v_{\rm S}(b)
   =\frac1{16}\sum_{s=1}^J\frac{\pi_s d_s(b)^2}{\ell_s-1}.
\end{equation}
For common block sizes two and four, $\Sigma_D=0$ in both cases but $\Psi_D=2\Pi$ and $\Psi_D=4\Pi/3$, respectively. Thus the usual limiting imbalance covariance alone does not determine the distribution obtained by regenerating the original assignments.

The design in Example~3.4 of \citet{BugniCanayShaikh2018} instead assigns a fixed number of treated participants uniformly across the entire stratum. It should not be identified with consecutive blocks of a fixed short size. Their permutation tests in Section~4.3 also recalculate studentised statistics, whereas the reference statistic here retains its observed score. Formula~(\ref{eq:supp-block-gap}) therefore does not contradict their weak-null results. The present calculation assumes known, fixed block boundaries and sizes; it does not assert the same formula when blocks are random or their boundaries are re-generated.


For block size two, these calculations concern random pairing within a fixed finite set of strata. Variance conservativeness due to treatment-effect heterogeneity is also studied in finely stratified experiments \citep{Fogarty2018} and matched-pair sampling-inference theory \citep{BaiRomanoShaikh2022}. Those results address their own estimators and pairing assumptions. In particular, pairing that makes continuous covariates increasingly close is different from the fixed-stratum arrival-order blocks considered here; no equivalence of the resulting variance formulas is presumed.

The distinction between tests of a sharp null and tests of an average-effect null also appears in the finite-population analysis of stratified and matched-pair experiments by \citet{Ding2017}. That comparison concerns fixed potential outcomes. The result here uses a fixed number of baseline strata and repeated sampling of participants, so its variance formula and estimand should be distinguished from finite-population paired-design formulas.

\subsection{Stratified Efron biased-coin allocation}\label{sec:supp-efron}
Consider any fixed finite set of strata, with iid profile probabilities $\pi_s>0$. Within stratum $s$, use Efron's biased coin with fixed $p_s\in(1/2,1)$: select the underrepresented treatment with probability $p_s$, and use a fair coin at balance. Each stratum has its own random numbers. All regenerated sequences use the same ordered profiles and independent random-number streams. Single-sequence balance and assignment covariances for this rule are established in the literature \citep{BugniCanayShaikh2018,MarkaryanRosenberger2010}. Here we verify the joint condition involving regenerated sequences.

\begin{proposition}\label{prop:supp-efron}
The specified design satisfies Assumption~\ref{ass:supp-general-allocation}, with
\begin{equation}\label{eq:supp-efron-cov}
 \Sigma_D=0,\qquad \Psi_D=\diag\{\pi_s\psi(p_s):s=1,\ldots,J\},\qquad \psi(p_s)\ge1.
\end{equation}
For a single stratum started from balance, let $C_m(p)=\E[\bm Z_{1:m}\bm Z_{1:m}^\top]$. Then the finite limit
\begin{equation}\label{eq:supp-efron-psi}
 \psi(p)=\lim_{m\to\infty}\frac{\|C_m(p)\|_{\rm F}^2}{m}
\end{equation}
exists. Consequently, at an average-effect null value,
\[
 v_{\rm R}(b)-v_{\rm S}(b)
      =\frac1{16}\sum_s\pi_s\{\psi(p_s)-1\}d_s(b)^2\ge0.
\]
\end{proposition}
\begin{proof}
Write $q_s=1-p_s$ and choose $0<\theta_s<\log(p_s/q_s)$. Then
$r_s=q_s e^{\theta_s}+p_s e^{-\theta_s}<1$. If $D$ is the current stratum imbalance, direct conditioning gives
\[
 \E(e^{\theta_s|D+Z|}\mid D=d)=
 \begin{cases}r_s e^{\theta_s|d|},&d\ne0,\\ e^{\theta_s},&d=0.\end{cases}
\]
For $k$ repetitions let $x_s^{(a)}$ be their stratum imbalances and set
$\mathcal W(x)=1+\sum_{a=1}^k\sum_{s=1}^J e^{\theta_s|x_s^{(a)}|}$.
A global step updates stratum $s$ with probability $\pi_s$. Linearity of conditional expectation, without independence of the updates across repetitions, gives
\[
 P_k\mathcal W\le\rho \mathcal W+c,\qquad
 \rho=\max_s\{1-\pi_s(1-r_s)\}<1,
\]
for a finite $c$. This also implies $\sup_n E_0 \mathcal W(X_n)<\infty$.

All admissible moves have admissible reverses. The class reachable from zero is irreducible and has period two: the parity of the total imbalance changes at each step, and a profile followed by the same profile with reversed signs permits a two-step return. Its even-step kernel is irreducible and aperiodic. From $P_k^2\mathcal W\le\rho^2 \mathcal W+c(1+\rho)$, for every integer $r\ge2$ we obtain, outside a finite sublevel set,
\[
 P_k^2\mathcal W\le \mathcal W-\tfrac12(1-\rho^2)\mathcal W^{1-1/r}.
\]
Lemma~\ref{lem:alloc-return} therefore gives every finite moment of the return time to simultaneous zero.

Regenerate the joint chain at those returns. For the vector consisting of centred profile counts and all nonempty sign products, the cycle rewards have mean zero and all moments. The proof of Theorem~\ref{thm:clt} applies: bounded rewards, the optional-stopping argument for the profile block, and label flips give the joint normal limit, vanishing cross-covariances between different subsets of repetitions, bounded normalised fourth moments and second-moment convergence. Exponential-moment boundedness of each state coordinate gives $D_{n,s}^{(a)}/\sqrt n\to0$ in $L^4$, so $\Sigma_D=0$.

For one stratum and two repetitions, the same argument shows that
$f_p(m):=\Var\{\sum_{i=1}^m Z_i^{(0)}Z_i^{(1)}\}$ satisfies $f_p(m)/m\to\psi(p)<\infty$. Conditional independence and label symmetry give $f_p(m)=\|C_m(p)\|_{\rm F}^2$. Its diagonal terms total $m$, which proves $\psi(p)\ge1$. Given all profiles, the random streams in distinct strata are independent and all product sums have mean zero; hence cross-stratum covariances vanish. In stratum $s$ the conditional variance is $f_{p_s}(N_{n,s})$. The convergence of $f_p(m)/m$ makes this ratio bounded, while $N_{n,s}/n\to\pi_s$ and $N_{n,s}\to\infty$ almost surely. Dominated convergence gives $n^{-1}E f_{p_s}(N_{n,s})\to\pi_s\psi(p_s)$. This identifies \eqref{eq:supp-efron-cov}. Substitution into \eqref{eq:supp-gap} completes the proof.
\end{proof}
The coefficient in \eqref{eq:supp-efron-psi} can be evaluated using the single-sequence assignment covariance, for which \citet[Corollary~2.1]{MarkaryanRosenberger2010} give exact finite-sample formulas. The result covers arbitrary finite stratum definitions; it does not extend the binary-factor range-method drift proof to multilevel marginal minimisation.

\subsection{A counterexample to a universal sign claim}\label{sec:supp-sign-counterexample}
The next construction isolates the sign issue under the general allocation condition. It is an illustrative sequential rule, rather than a proposed clinical-trial design. Take two iid profiles with probabilities $(1/2,1/2)$ and group consecutive arrivals into pairs. In each pair assign the first participant by a fair coin. If the profiles differ, assign the second participant to the opposite treatment with probability $r\in(1/2,1)$; if they agree, use an independent fair coin. Different pairs use independent random numbers, and regenerated sequences retain these consecutive pairs and use fresh random numbers. Put $a=2r-1$.

\Needspace{9\baselineskip}
\begin{proposition}\label{prop:supp-negative-example}
This sequential rule satisfies Assumption~\ref{ass:supp-general-allocation}, with
\[
 \Sigma_D=\begin{pmatrix}1/2&-a/4\\-a/4&1/2\end{pmatrix},\qquad
 \Psi_D=\begin{pmatrix}1/2&a^2/4\\a^2/4&1/2\end{pmatrix}.
\]
For $d=(t,-t)^\top$ with $t\ne0$, the variance difference at $\Delta=b$ is $-a^2t^2/32<0$.
\end{proposition}
\begin{proof}
For the two signs within a pair, the conditional correlation is $-a$ when the profiles differ and zero otherwise. Its contribution to the off-diagonal entry of the covariance of the two profile imbalances is therefore $-a/2$ per pair, or $-a/4$ per participant. Each diagonal entry is $1/2$: the expected count of either profile is one per pair and, when both profiles agree, the two signs are independent. For two independent repetitions, the conditional correlation of the sign products is the square of the original correlation. This gives diagonal entries $1/2$ and off-diagonal entries $a^2/4$ for $\Psi_D$.

Stack the centred profile counts and every nonempty sign product of three repetitions, summing within consecutive pairs. The block vectors are iid and bounded. Label flips give zero means for the sign blocks and zero cross-covariances between distinct subsets and with the centred profile block. Thus the multivariate iid central limit theorem gives the required joint limit; the bounded incomplete pair is negligible. Fourth moments of centred block sums are $O(n^2)$, second moments converge, and the triple product is tight. Finally,
$d^\top(\Psi_D-\Pi)d=-a^2t^2/2$ gives the displayed difference. The continuous-outcome construction in Proposition~\ref{prop:supp-sign} supplies a nondegenerate weak-null example.
\end{proof}
This example shows why positive semidefiniteness of $\Psi_D$ alone cannot establish conservativeness. A sign conclusion for minimisation requires the restricted quadratic form of that design, rather than the label-symmetry and joint-limit conditions alone.

\section{Sign of the variance difference under minimisation}\label{sec:supp-sign-min}
This appendix proves the allocation examples summarised in Section~4.3 of the main text: the proposition giving a sufficient condition for equality of the variances, the corollary on deterministic one-factor range minimisation, the Walsh diagonalisation under the uniform profile law, and the three-participant computation. It then describes the exact recursion used for the restricted ratios in the main text and reports the complete values. Throughout, $\bm G_n=\sum_{i\le n}\bm h(S_i)\xi_i$ with $\xi_i=Z_i^{(1)}Z_i^{(2)}$ for two independent repetitions driven by the same profiles, so that $\Var(\bm G_n)/n=n^{-1}\E\{\bm H^\top(\mathcal C_n\circ\mathcal C_n)\bm H\}$ by Proposition~\ref{prop:supp-sign}, and $\mathcal H_{i-1}$ is the $\sigma$-field generated by $S_1,\ldots,S_{i-1}$ and the allocation random numbers of both repetitions for the first $i-1$ participants.

\subsection{A sufficient condition for equality of the two variances}
\begin{proposition}\label{prop:supp-blind}
Suppose that, for every $i$, $\E(\xi_i\mid\mathcal H_{i-1},S_i)$ does not depend on $S_i$. Then for every $n$ and every $v$ with $\bm\pi^\top v=0$, $\Var(v^\top\bm G_n)=n\,v^\top\bm\Pi v$. If Assumption~\ref{ass:supp-general-allocation} also holds, then $v^\top(\bm\Psi_D-\bm\Pi)v=0$ and $v_{\rm R}(b)=v_{\rm S}(b)$ at every null value $b$.
\end{proposition}
\begin{proof}
Let $\zeta_i=v_{S_i}\xi_i$ and let $\mathcal H_i$ be generated by $\mathcal H_{i-1}$, $S_i$ and the random numbers of step $i$. Since $S_i$ is independent of $\mathcal H_{i-1}$ and $\psi_i=\E(\xi_i\mid\mathcal H_{i-1},S_i)$ is $\mathcal H_{i-1}$-measurable by assumption,
\[
 \E(\zeta_i\mid\mathcal H_{i-1})=\E\{v_{S_i}\E(\xi_i\mid\mathcal H_{i-1},S_i)\mid\mathcal H_{i-1}\}
 =\psi_i\E(v_{S_i})=\psi_i\bm\pi^\top v=0.
\]
Thus $(\zeta_i)$ is a martingale difference sequence with respect to $(\mathcal H_i)$, and $\E\zeta_i^2=\E v_{S_i}^2=v^\top\bm\Pi v$ because $\xi_i^2=1$. Hence $\E\{(\sum_{i\le n}\zeta_i)^2\}=\sum_{i\le n}\E\zeta_i^2=n\,v^\top\bm\Pi v$. Under Assumption~\ref{ass:supp-general-allocation} the normalised second moments converge, so $v^\top\bm\Psi_Dv=v^\top\bm\Pi v$, and Corollary~\ref{cor:supp-variance-equality} gives $v_{\rm R}(b)=v_{\rm S}(b)$.
\end{proof}
For the stochastic range method, \eqref{eq:mean} gives $\E(\xi_i\mid\mathcal H_{i-1},S_i)=\eta^2f_{x^{(1)}}(U_i)f_{x^{(2)}}(U_i)$, which depends on the incoming profile unless the two repetitions occupy states for which $f_{x^{(1)}}(u)f_{x^{(2)}}(u)$ is constant in $u$; the condition therefore fails in general.

\subsection{Deterministic one-factor range minimisation}
\begin{corollary}\label{cor:supp-det1}
Let $F=1$, $p_{\rm bc}=1$ with ties resolved by a fair coin, and $\pi_+\in(0,1)$. The contrast state takes values in $\{(0,0),(\pm1,\pm1),(0,\pm2)\}$, ties occur only at $(0,0)$, and two independent repetitions driven by the same profiles satisfy $x^{(2)}_i=\epsilon\,x^{(1)}_i$ with a sign $\epsilon$ that is constant between visits to $(0,0)$ and is refreshed by independent fair coins at those visits. Proposition~\ref{prop:supp-blind} applies, Assumption~\ref{ass:supp-general-allocation} holds with $\Sigma_D=0$ and
\begin{equation}\label{eq:supp-det1-psi}
 \Psi_D=\bm\Pi+\frac{1+\pi_+-\pi_+^2}{\pi_+(1-\pi_+)}\,\bm\pi\bm\pi^\top .
\end{equation}
\end{corollary}
\begin{proof}
Write the state as $(D_+,D_-)$, the imbalances at the two levels, so that $D_0=D_++D_-$ and the contrast state is $(D_++D_-,D_+-D_-)$. With $p_{\rm bc}=1$ the sign of participant $i$ with level $u$ is $-\sg\{\sg(D_0)+\sg(D_u)\}$, with a fair coin when this is zero. From $(0,0)$ every arrival leads to $(\pm1,0)$ or $(0,\pm1)$; from $(1,0)$ an arrival at level $+$ returns to $(0,0)$ and an arrival at level $-$ leads to $(1,-1)$; from $(1,-1)$ the two levels lead to $(0,-1)$ and $(1,0)$; from $(0,-1)$ they lead to $(1,-1)$ and $(0,0)$; the remaining states follow by symmetry. The reachable class is therefore $\{(0,0),(\pm1,0),(0,\pm1),(1,-1),(-1,1)\}$, which is the stated set of contrast states. A tie requires $\sg(D_0)=-\sg(D_u)\ne0$ or $D_0=D_u=0$; the first needs $|D_{-u}|\ge2$, which is impossible in the class, and the second forces $(0,0)$.

Two repetitions start at $(0,0)$ and draw independent fair coins, so after the first participant $x^{(2)}_1=\epsilon x^{(1)}_1$ with $\epsilon=Z^{(1)}_1Z^{(2)}_1$. If $x^{(2)}=\epsilon x^{(1)}\ne0$, then $f_{x^{(2)}}(u)=\epsilon f_{x^{(1)}}(u)\ne0$ by \eqref{eq:fx}, both decisions are deterministic and $Z^{(2)}=-f_{x^{(2)}}(u)=\epsilon Z^{(1)}$, so the relation is preserved and $\xi_i=\epsilon$ until both repetitions are at $(0,0)$, where fresh fair coins are drawn. Hence $\E(\xi_i\mid\mathcal H_{i-1},S_i)$ equals $\epsilon$ off balance and $0$ at balance, independently of $S_i$, and Proposition~\ref{prop:supp-blind} applies. The joint chain of any fixed number of repetitions lives on a finite class and is irreducible with period two, so the regenerative argument of Theorem~\ref{thm:clt} applies verbatim and yields Assumption~\ref{ass:supp-general-allocation}; the within-stratum imbalances are bounded, so $\Sigma_D=0$. Formula~\eqref{eq:supp-det1-psi} was obtained by solving the first-step equations for the mean cycle length and the second moment of the cycle reward of the $13$-state pair chain in exact rational arithmetic in $\pi_+$, as in the proof of Theorem~\ref{thm:clt} (file \texttt{exact\_p1\_F1.py}); the recursion of Appendix~I.4 reproduces it to $O(n^{-1})$ for $n$ up to $4000$. Within an excursion from balance all signs are deterministic functions of the profiles, so the squared correlation of any two participants of the same excursion equals one.
\end{proof}

\subsection{Uniform profile law and three participants}
\begin{proposition}\label{prop:supp-walsh}
Let the factors be independent with prevalence $1/2$ each and let the rule be the equal-weight range method. For $A\subseteq\{1,\ldots,F\}$ put $\chi_A(u)=\prod_{f\in A}u_f$ for sign-coded profiles $u$. Then $\Var(\bm G_n)/n$ is diagonalised by the characters $\chi_A$ for every $n$. If Assumption~\ref{ass:supp-general-allocation} holds, so is $\Psi_D$, with $\Psi_D\chi_A=2^{-F}\{1+\widehat g(A)\}\chi_A$, where $\widehat g(A)=2^{-F}\sum_{w\in\{-1,1\}^F}\widetilde g(w)\chi_A(w)$ and $\widetilde g(w)=\lim_nn^{-1}\sum_{i\ne j}\E(\mathcal C_{n,ij}^2\mid U_i\circ U_j=w)$, $\circ$ denoting coordinatewise multiplication. The eigenvalues of $\mathcal R_\pi$ are $1+\widehat g(A)$ for $A\ne\varnothing$, they depend on $A$ only through $|A|$, and $\lambda_{\min}(\mathcal R_\pi)\ge1$ if and only if $\widehat g(A)\ge0$ for every nonempty $A$.
\end{proposition}
\begin{proof}
Relabelling the two levels of factor $f$ for all participants leaves the uniform profile law unchanged and leaves the allocation law given the profiles unchanged, because the imbalance at each participant's own level is the same before and after the relabelling. Hence $\E(\mathcal C_{n,ij}^2\mid U_i=s,U_j=t)$ is invariant under $(s,t)\mapsto(s^f,t^f)$ for every $f$, where $s^f$ reverses the $f$th coordinate, and depends on $(s,t)$ only through $s\circ t$. Since $\Pp(U_i=s,U_j=t)=4^{-F}$ for $i\ne j$,
\[
 \Var(\bm G_n)/n=2^{-F}I+4^{-F}\{\widetilde g_n(s\circ t)\}_{s,t},\qquad
 \widetilde g_n(w)=n^{-1}\sum_{i\ne j}\E(\mathcal C_{n,ij}^2\mid U_i\circ U_j=w).
\]
The second matrix is a convolution operator on the group $\{-1,1\}^F$, and
$\sum_t\widetilde g_n(s\circ t)\chi_A(t)=\chi_A(s)\sum_w\widetilde g_n(w)\chi_A(w)$, so the characters are eigenvectors with the stated eigenvalues; the limit follows from the second-moment convergence in Assumption~\ref{ass:supp-general-allocation}. Because $\bm\Pi=2^{-F}I$ and $\sqrt{\bm\pi}\propto\chi_\varnothing$, the eigenvalues of $\mathcal R_\pi$ are $1+\widehat g(A)$ for $A\ne\varnothing$. Permuting factors preserves the law under equal weights, so $\widehat g(A)$ depends only on $|A|$.
\end{proof}

\begin{proposition}\label{prop:supp-n3}
For the range method with $n=3$ participants and any profile law, $\mathcal C_{12}=-\eta$, $\mathcal C_{13}=-\eta(q+p\sigma)$ and $\mathcal C_{23}=-\eta(q-p\sigma)$, where $p=p_{\rm bc}$, $q=1-p$, $\sigma=\sg(m_1-m_2)$, and $m_1$ and $m_2$ count the factors on which participants 1 and 2 differ and participant 3 agrees with participant 1 and with participant 2, respectively. Consequently, for $\bm\pi^\top v=0$,
\begin{equation}\label{eq:supp-n3}
 v^\top\{\Var(\bm G_3)/3-\bm\Pi\}v
 =\tfrac23\eta^2\E\{v_{S_1}v_{S_3}(q+p\sigma)^2+v_{S_2}v_{S_3}(q-p\sigma)^2\}.
\end{equation}
Under the uniform law, along $v=\chi_{\{1,\ldots,F\}}$,
\begin{align}\label{eq:supp-n3-uniform}
 \frac{v^\top\Var(\bm G_3)v}{3v^\top\bm\Pi v}-1
 &=\frac{2\eta^2}{3\cdot2^{F}}
 \Big[4pq\,\E\{(-1)^{F-m}\sg(2m-F)\}\nonumber\\
 &\qquad\quad+2p^2\E\{(-1)^{F-m}\mathbf 1(2m\ne F)\}\Big],
 \qquad m\sim{\rm Bin}(F,1/2),
\end{align}
which equals $\tfrac43pq\eta^2$ ($F=1$), $\tfrac16p^2\eta^2$ ($F=2$), $-\tfrac16pq\eta^2$ ($F=3$) and $-\tfrac1{32}p^2\eta^2$ ($F=4$).
\end{proposition}
\begin{proof}
$Z_1$ is a fair sign. With $X_1=Z_1a(U_1)$ every term in \eqref{eq:fx} equals $\sg\{Z_1(1+u_fU_{1f})\}\in\{0,Z_1\}$, so $f_{X_1}(u)=Z_1$ for every $u$ and $Z_2=-Z_1$ with probability $p$; thus $\mathcal C_{12}=-\eta$. If $Z_2=-Z_1$, then $X_2=Z_1(0,U_1-U_2)$ and $f_{X_2}(U_3)=Z_1\sg(m_1-m_2)=Z_1\sigma$, because a factor on which participants 1 and 2 differ contributes $\sg(2Z_1U_{1f}U_{3f})=Z_1U_{1f}U_{3f}$ and the other factors contribute zero. If $Z_2=Z_1$, then $X_2=Z_1(2,U_1+U_2)$, every term in \eqref{eq:fx} is $0$ or $Z_1$, and $f_{X_2}(U_3)=Z_1$. Hence $\E(Z_3\mid S,Z_1,Z_2)=-\eta Z_1\sigma$ or $-\eta Z_1$ in the two cases, which gives $\mathcal C_{13}=\E(Z_1Z_3\mid S)=-\eta(q+p\sigma)$ and $\mathcal C_{23}=-\eta(q-p\sigma)$. The identity \eqref{eq:supp-n3} follows from \eqref{eq:supp-offdiagonal-square} and $\E(v_{S_1}v_{S_2})=0$.

Under the uniform law with $v=\chi_{\{1,\ldots,F\}}$, $\E(v_{S_1}v_{S_3}\sigma^k)$ for $k\in\{1,2\}$ vanishes unless participants 1 and 2 differ on every factor, because otherwise $v_{S_1}v_{S_3}$ contains a factor $U_{1f}U_{3f}$ that is a fair sign independent of $\sigma$; the event has probability $2^{-F}$, and conditionally $m_1\sim{\rm Bin}(F,1/2)$, $m_2=F-m_1$, $v_{S_1}v_{S_3}=(-1)^{m_2}$ and $\sigma=\sg(2m_1-F)$. Interchanging participants 1 and 2 maps $\sigma$ to $-\sigma$ and $v_{S_1}v_{S_3}$ to $v_{S_2}v_{S_3}$. Expanding the squares in \eqref{eq:supp-n3} and collecting terms gives \eqref{eq:supp-n3-uniform}. For odd $F$ the second expectation vanishes and the first equals $(-1)^{(F-1)/2}\binom{F-1}{(F-1)/2}2^{-(F-1)}$; for even $F$ the first expectation vanishes by symmetry and the second equals $-(-1)^{F/2}\binom{F}{F/2}2^{-F}$. The four stated values follow, and they were also checked by exact enumeration of all profile and assignment sequences (file \texttt{exact\_bruteforce.py}).
\end{proof}

\subsection{Forward recursion for the finite-sample covariance}\label{sec:supp-dp}
Let $V_i=(x^{(1)}_i,x^{(2)}_i)$ be the pair of contrast states after $i$ participants, and for a pair state $\varsigma$ define $P_i(\varsigma)=\Pp(V_i=\varsigma)$, $\bm M_i(\varsigma)=\E\{\bm G_i\mathbf 1(V_i=\varsigma)\}\in\R^J$ and $\bm S_i=\E(\bm G_i\bm G_i^\top)$. Given $V_i=\varsigma=(x^{(1)},x^{(2)})$, the next participant has profile $u$ with probability $\pi_u$ and signs $(z^{(1)},z^{(2)})$ with probability $\prod_{m=1}^2\{1-\eta z^{(m)}f_{x^{(m)}}(u)\}/2$ by \eqref{eq:mean}, leading to $\varsigma'=(x^{(1)}+z^{(1)}a(u),x^{(2)}+z^{(2)}a(u))$ with reward $\bm r=h(u)z^{(1)}z^{(2)}$. Writing $w$ for the product of these probabilities,
\begin{align*}
 P_{i+1}(\varsigma')&=\sum wP_i(\varsigma),\qquad
 \bm M_{i+1}(\varsigma')=\sum w\{\bm M_i(\varsigma)+P_i(\varsigma)\bm r\},\\
 \bm S_{i+1}&=\bm S_i+\sum w\{\bm M_i(\varsigma)\bm r^\top
                 +\bm r\bm M_i(\varsigma)^\top+P_i(\varsigma)\bm r\bm r^\top\},
\end{align*}
where the sums run over all $(\varsigma,u,z^{(1)},z^{(2)})$ leading to $\varsigma'$. Since $\E\bm G_n=0$, $\Var(\bm G_n)=\bm S_n$. The recursion is exact before numerical truncation. Floating-point evaluation introduces rounding error in addition to state truncation; the reported symmetry checks are numerical diagnostics, not certified rounding-error bounds. To keep the number of pair states manageable, states whose probability falls below a tolerance are discarded together with their descendants. If the discarded paths have total probability $\varepsilon_n$, then, since $|v^\top\bm G_n|\le n\max_s|v_s|$ on every path and $v^\top\bm\Pi v\ge\min_s\pi_s\max_sv_s^2$, the resulting error in $v^\top\Var(\bm G_n)v/(nv^\top\bm\Pi v)$ is at most $\varepsilon_nn/\min_s\pi_s$ in absolute value. The tolerance was $10^{-16}$ for one factor ($\varepsilon_{80}\le1.4\times10^{-10}$, error below $10^{-7}$), $10^{-13}$ or $10^{-14}$ for two factors ($\varepsilon_{40}\le4.5\times10^{-7}$, error below $10^{-4}$), zero for three factors with $n\le9$ (no state discarded), $10^{-11}$ for three factors with $n=10$ ($\varepsilon_{10}=7.3\times10^{-6}$, error below $6\times10^{-4}$) and $10^{-12}$ for the three-factor law with unequal prevalences ($\varepsilon_8=6.4\times10^{-8}$, error below $10^{-5}$). Entries are therefore reported to four decimals for one and two factors and to three or four decimals for three factors. The restricted eigenvalues are those of $B_\pi^\top\bm\Pi^{-1/2}\{\Var(\bm G_n)/n\}\bm\Pi^{-1/2}B_\pi$ as in Proposition~\ref{prop:supp-sign}. For the uniform law, the computed matrices are diagonal in the Walsh basis to within $10^{-15}$, as Proposition~\ref{prop:supp-walsh} requires, and for $p_{\rm bc}=1$ with one factor the ratio equals one to machine precision for every $n\le60$, as Corollary~\ref{cor:supp-det1} requires. The values for $n=3$ and $n=4$ with three factors agree with exact rational enumeration ($619/625$ and $398389/400000$ along the three-way interaction). The code is supplied with the computational files.

\begin{table}[!htbp]
\centering
\caption{Restricted ratios $v^\top\Var(\bm G_n)v/(nv^\top\bm\Pi v)$ on $\bm\pi^\top v=0$ for the stochastic range method with one factor ($J=2$, $p=p_{\rm bc}$), from the deterministic recursion with tolerance $10^{-16}$.}
\label{tab:supp-exact-sign}
\begin{threeparttable}
\begin{singlespace}
\footnotesize
\setlength{\tabcolsep}{3.5pt}
\begin{tabular}{@{}llllllll@{}}
\toprule
$(\pi_+,p)$ & $n=10$ & $n=20$ & $n=40$ & $n=60$ & $n=80$ & $1/n$ extrapolation & Discarded mass\\
\midrule
$(.5,.6)$ & 1.0208 & 1.0208 & 1.0188 & 1.0168 & 1.0158 & 1.013 & $1.4\times10^{-10}$\\
$(.5,.7)$ & 1.0634 & 1.0585 & 1.0501 & 1.0455 & 1.0430 & 1.036 & $7.2\times10^{-11}$\\
$(.5,.8)$ & 1.0977 & 1.0899 & 1.0815 & 1.0781 & 1.0764 & 1.071 & $3.0\times10^{-11}$\\
$(.3,.8)$ & 1.0996 & 1.0956 & 1.0855 & 1.0803 & 1.0778 & 1.070 & $2.9\times10^{-11}$\\
$(.1,.8)$ & 1.0762 & 1.1018 & 1.1058 & 1.0994 & 1.0936 & 1.081 & $2.0\times10^{-11}$\\
$(.5,.95)$ & 1.0597 & 1.0595 & 1.0593 & 1.0592 & 1.0592 & 1.059 & $2.8\times10^{-12}$\\
$(.3,.95)$ & 1.0604 & 1.0583 & 1.0558 & 1.0549 & 1.0544 & 1.053 & $2.8\times10^{-12}$\\
$(.1,.95)$ & 1.0505 & 1.0598 & 1.0541 & 1.0485 & 1.0450 & 1.036 & $2.8\times10^{-12}$\\
$(.5,.99)$ & 1.0146 & 1.0148 & 1.0150 & 1.0150 & 1.0150 & 1.015 & $3.3\times10^{-13}$\\
$(.5,.999)$ & 1.0015 & 1.0016 & 1.0016 & 1.0016 & 1.0016 & 1.0016 & $4.7\times10^{-13}$\\
any $\pi_+$, $p=1$ & 1 & 1 & 1 & 1 & 1 & 1 & 0\\
\bottomrule
\end{tabular}
\begin{tablenotes}[flushleft]\footnotesize
\item The $1/n$ extrapolation uses $n=40$ and $n=80$. The discarded mass at $n=80$ bounds the error of every entry by $\varepsilon_{80}\cdot80/\min_s\pi_s<10^{-7}$. The smallest ratio over $1\le n\le80$ equals one in every row, attained at $n\le2$.
\end{tablenotes}
\end{singlespace}
\end{threeparttable}
\end{table}

\begin{table}[!htbp]
\centering
\caption{Restricted eigenvalues of $B_\pi^\top\bm\Pi^{-1/2}\{\Var(\bm G_n)/n\}\bm\Pi^{-1/2}B_\pi$ for the stochastic range method with two and three factors ($p=p_{\rm bc}$), from the deterministic recursion; $\varepsilon_n$ is the discarded mass at the largest $n$ shown.}
\label{tab:supp-exact-sign-23}
\begin{threeparttable}
\begin{singlespace}
\footnotesize
\begin{tabular}{@{}p{56mm}p{98mm}@{}}
\toprule
Design & Eigenvalues by $n$\\
\midrule
$F=2$, $(.5,.4)$, $p=.8$, three eigenvalues & $n=10$: 1.0586, 1.0836, 1.0837; $n=20$: 1.0604, 1.0790, 1.0796; $n=30$: 1.0573, 1.0743, 1.0746; $n=40$: 1.0544, 1.0710, 1.0710 ($\varepsilon_{40}=4.5\times10^{-7}$)\\
$F=2$, $(.5,.4)$, $p=.95$, three eigenvalues & $n=10$: 1.0594, 1.0622, 1.1894; $n=20$: 1.0640, 1.0666, 1.2050; $n=30$: 1.0653, 1.0674, 1.2054 ($\varepsilon_{30}=3.9\times10^{-9}$)\\
$F=2$, uniform, $p=.8$, Walsh characters & $n=30$: $|A|=1$: 1.0734 (twice); $|A|=2$: 1.0594 ($\varepsilon_{30}=2.3\times10^{-7}$)\\
$F=3$, uniform, $p=.8$, $|A|=3$ & $n=3,\ldots,10$: .9904, .9960, .9949, .9964, .9962, .9969, .9969, .9974 (exact for $n\le9$; $\varepsilon_{10}=7.3\times10^{-6}$)\\
$F=3$, uniform, $p=.8$, $|A|=2$ & $n=3,4,6,8,10$: 1.0192, 1.0160, 1.0285, 1.0356, 1.0398\\
$F=3$, uniform, $p=.8$, $|A|=1$ & $n=3,4,6,8,10$: 1.0480, 1.0514, 1.0643, 1.0683, 1.0694\\
$F=3$, $(.5,.4,.3)$, $p=.8$, smallest eigenvalue & $n=3,4,6,8$: .9921, .9964, .9966, .9970 ($\varepsilon_8=6.4\times10^{-8}$)\\
\bottomrule
\end{tabular}
\begin{tablenotes}[flushleft]\footnotesize
\item The Monte Carlo estimates at $n=200$ and $n=1000$ in Supplementary Table~\ref{tab:supp-eigen-diagnostics} for two factors ($1.046$--$1.062$ and $1.044$--$1.056$ at $p=.8$; $1.063$--$1.197$ and $1.063$--$1.196$ at $p=.95$) continue the computed sequences. For the uniform three-factor law the three eigenvalues within each class $|A|$ coincide to $10^{-15}$, as Proposition~\ref{prop:supp-walsh} requires.
\end{tablenotes}
\end{singlespace}
\end{threeparttable}
\end{table}

For the reported one-factor settings with $1/2<p_{\rm bc}<1$, the computed ratios exceed one for $3\le n\le80$; the ratio is one at $n=1,2$ and at every $n$ for the deterministic one-factor rule. The $1/n$ extrapolations are numerical summaries, not analytic limits. All reported two-factor eigenvalues exceed one. The three-factor interaction ratios are below one at the evaluated sample sizes, with both increases and decreases as $n$ varies. Neither a small finite-sample deficit nor an extrapolation proves the sign of the limiting matrix. The directional covariance ratio must also be distinguished from the reference-to-sampling variance ratio before drawing conclusions about rejection probability.

\section{Null-value study with heterogeneous effects and the corrected test}\label{sec:supp-hetero}
This appendix retains the earlier null studies using the original covariance construction; they are distinct from the paired evaluation in Appendix~K. The $.80$ tables were checked against the supplied aggregate output. The $.95$ tables were supplied as a rounded reconstruction of earlier aggregator summaries rather than an original run file. They are retained as developmental summaries and are not used as independently reproduced validation of the covariance reconstruction.

The results in this appendix use the original covariance construction and the tail conventions specified below. They do not evaluate the new covariance reconstruction or its corrected lattice tail.


\subsection{Design}
This study uses two configurations of the four main design--sample-size combinations: two factors with $n=200$ or $1000$, and five factors with $n=400$ or $2000$. Factor prevalences are those of Appendix~E and factors are independent. Each configuration has 32 scenarios, combining both outcome types and four tested null values: superiority ($b=0$, two-sided at 5\%), non-inferiority ($b=-.20$ for continuous and $b=-.10$ for binary outcomes, one-sided at 2.5\%), and the lower and upper equivalence limits ($\pm.45$ for continuous outcomes; $\pm.21$, $\pm.10$, $\pm.15$ and $\pm.10$ for binary outcomes in the four designs). Equivalence uses two one-sided tests at 5\%. Each scenario is generated with marginal effect equal to the corresponding tested boundary. Both unadjusted and baseline-adjusted scores are evaluated.

In the first configuration $p_{\rm bc}=.80$. For continuous outcomes,
\[
 Y_i(0)=\mu_{S_i}+\varepsilon_i,\qquad
 Y_i(1)=Y_i(0)+\Delta+d_{S_i}+\eta_i,
\]
where $\mu_s$ is as in Appendix~G, $\varepsilon_i\sim N(0,1)$ and $\eta_i\sim N(0,.25^2)$. The direction $v$ is the first factor centred by its $\bm\pi$-weighted mean, and $d_s=.5v_s/\max_t|v_t|$, so $\bm\pi^\top\bm d=0$. For binary outcomes, the control risks are those of Appendix~E.2, with marginal control risk $.60$, and treatment risks are $p_{0s}+\Delta+d_s$. The same centred direction has maximum absolute deviation $.15$, reduced where necessary to keep probabilities in $(.02,.98)$. This reduction gives a maximum deviation of approximately $.109$ at the upper equivalence limits of the two-factor $n=200$ and five-factor $n=400$ designs. Observed binary outcomes are generated with the probabilities of their assigned arms.

The second configuration uses $p_{\rm bc}=.95$. The continuous model has residual standard deviation $.25$, individual-effect standard deviation $.10$ and maximum absolute stratum-effect deviation one. The binary probabilities, including the feasibility reductions, are otherwise unchanged. Thus the continuous comparison changes the coin, residual variability and effect heterogeneity together; it does not isolate the effect of any one of them.

Each scenario uses 100,000 outer trials, 4,999 regenerated sequences per trial and 100,000 three-sequence allocation replicates for covariance estimation. All procedures are compared on common trials. The methods are the reference test (RT); the tie-inclusive corrected test using model-specified deviations, estimated deviations, or an exploratory adjustment for their estimation noise; and SIGA-S and SIGA-R using model-specified deviations for the reference variance. The data-based deviation estimator is that in Appendix~C.2, with zero assigned to a stratum if either arm is absent. Model-specified deviations provide a simulation benchmark rather than an available estimator for a completed trial with unknown effects. All Gaussian tails in this appendix are computed without the continuity correction of Appendix~D. This differs from the unadjusted binary superiority calculation in the original main study.

The exploratory adjustment subtracts
\[
 \sum_s(\widehat{\bm\Psi}_n-\widehat{\bm\Pi}_n)_{ss}\widehat v_s
\]
from $\widehat{\bm d}(b)^\top(\widehat{\bm\Psi}_n-\widehat{\bm\Pi}_n)\widehat{\bm d}(b)$, where $\widehat v_s$ estimates the variance of the within-stratum arm-mean difference. It is a diagonal noise adjustment, not an asserted unbiased correction under the adaptive design. If $\max_s\widehat v_s=O_p(n^{-1})$ with fixed $J$, the subtracted quantity is $O_p(n^{-1})$ and does not change the first-order target. Its finite-sample effect is assessed separately from the data-based estimator covered by the main theorem.

The rejection rates are displayed in percent and their Monte Carlo standard errors in parentheses. With 100,000 trials these errors are approximately 0.07 percentage points at 5\% and 0.05 at 2.5\%. They describe outer-trial uncertainty conditional on the simulated covariance estimates, not all sources of simulation error. Comparisons below are descriptive differences computed from the displayed rates. The tables do not provide the joint rejection counts needed for paired standard errors. The variance diagnostics average trial-level ratios first across trials and then across the four null values; they do not bound the error in each trial.

\subsection{Results under \texorpdfstring{$p_{\rm bc}=.80$}{p=0.80}}
Tables~\ref{tab:supp-hetero-null} and~\ref{tab:supp-hetero-diag} give the complete results. All six procedures are within 0.22 percentage points of nominal in the continuous scenarios. The data-based corrected rates differ from the model-based rates by at most 0.04 points. For unadjusted binary superiority, RT and the three corrected versions have the same displayed rates: 4.26, 4.55, 4.49 and 4.75\% across the four designs. Ordinary Gaussian tails reject 4.85--5.06\%. Away from these unadjusted binary superiority rows, the data-based corrected test differs from RT by at most 0.06 points in the displayed table. Similarity to the reference test is distinct from similarity to nominal size.

The mean model-based reference-to-sampling variance ratios range from 1.0007 to 1.0034. All displayed $\widehat V^*/\widehat V_{{\rm R},n}$ diagnostics are within 0.3\% of one except the unadjusted continuous score with five factors and $n=400$, for which the mean is .9881. Estimation of stratum score means is a possible contributor, but these diagnostics alone do not identify the source of the discrepancy. Equality of rounded rejection rates also does not prove equality of trial-level decisions. Although the unadjusted binary score has unit lattice spacing conditional on the treated count, the regenerated count can vary, so a small rescaling need not leave every rejection decision unchanged.

\subsection{Results under \texorpdfstring{$p_{\rm bc}=.95$}{p=0.95}}
Tables~\ref{tab:supp-hetero-null-stress} and~\ref{tab:supp-hetero-diag-stress} give the second configuration. In continuous scenarios the averaged model-based variance ratios range from 1.0553 to 1.0715. RT rejects 4.25--4.63\% at nominal 5\% and 2.11--2.20\% at nominal 2.5\%. The model-based corrected ranges are 4.93--5.04\% and 2.44--2.56\%; the data-based ranges are 4.85--5.04\% and 2.44--2.50\%. The largest absolute departures from nominal are 0.07, 0.15 and 0.16 points for the model-based, data-based and noise-adjusted corrected versions, respectively. These findings support the correction in these continuous scenarios, rather than establish uniform finite-sample control.

Gaussian approximation is less accurate for the unadjusted continuous score with five factors. At $n=400$, the superiority rates are 4.30\% for RT, 4.99\% for the model-based correction, 4.87\% for the data-based correction, 4.21\% for SIGA-S and 3.63\% for SIGA-R. The mean diagnostic $\widehat V^*/\widehat V_{{\rm R},n}=.9327$ shows conditional-variance overestimation; it does not by itself identify error in the sampling variance. Baseline adjustment changes this diagnostic to .9965, with superiority rates 4.25\% for RT, 4.85\% for the data-based correction, 4.94\% for SIGA-S and 4.21\% for SIGA-R. At $n=2000$ the unadjusted diagnostic is .9837. Excluding the unadjusted five-factor continuous rows, the maximum displayed difference between SIGA-R and RT is 0.23 points, and SIGA-S is within 0.21 points of nominal. A common component of error in the two estimated variances does not cancel exactly in their ratio, so good performance of the corrected test is not a proof of such cancellation.

For binary outcomes, the averaged model-based variance ratios are 1.0017--1.0023. The unadjusted superiority rates for RT and all three corrected versions are, to the displayed precision, 3.67, 4.41, 4.13 and 4.68\%. Thus variance correction with retained ties does not remove the conservativeness in these settings. Ordinary Gaussian tails do not reproduce these finite-sample reference tails, despite close average variance diagnostics. In the two-factor $n=200$ design, the data-based corrected rate differs from RT by up to 0.35 points, or 0.27 after the exploratory noise adjustment. Relative to the model-based correction the corresponding maxima are 0.32 and 0.24 points. The larger estimated rescaling is consistent with sensitivity to estimation noise when the variance correction itself is small, but neither the cause nor general superiority of the noise adjustment is established by these comparisons.

\clearpage
\begin{landscape}\begin{singlespace}\begingroup\footnotesize\setlength{\tabcolsep}{3.3pt}\renewcommand{\arraystretch}{1.0}\setlength{\LTcapwidth}{\linewidth}\setlength{\LTleft}{\fill}\setlength{\LTright}{\fill}
\begin{longtable}{@{}llllrrrrrr@{}}
\caption{Rejection probabilities (\%) at the null values under heterogeneous treatment effects, with Monte Carlo standard errors in parentheses.}
\label{tab:supp-hetero-null}\\
\toprule
Design & Outcome & Null value & Score & RT & CRT (model) & CRT (data) & CRT (data, adj.) & SIGA-S & SIGA-R\\
\midrule\endfirsthead
\multicolumn{10}{c}{\tablename\ \thetable{} -- continued}\\\toprule
Design & Outcome & Null value & Score & RT & CRT (model) & CRT (data) & CRT (data, adj.) & SIGA-S & SIGA-R\\
\midrule\endhead
\midrule\multicolumn{10}{r}{Continued on next page}\\\endfoot
\bottomrule\endlastfoot
$F=2$, $n=200$ & continuous & superiority, $0.00$ & unadjusted & 5.08 (0.07) & 5.13 (0.07) & 5.16 (0.07) & 5.14 (0.07) & 5.08 (0.07) & 5.04 (0.07)\\
$F=2$, $n=200$ & continuous & superiority, $0.00$ & adjusted & 5.08 (0.07) & 5.12 (0.07) & 5.16 (0.07) & 5.14 (0.07) & 5.08 (0.07) & 5.05 (0.07)\\
$F=2$, $n=200$ & continuous & non-inferiority, $-0.20$ & unadjusted & 2.43 (0.05) & 2.45 (0.05) & 2.48 (0.05) & 2.46 (0.05) & 2.40 (0.05) & 2.39 (0.05)\\
$F=2$, $n=200$ & continuous & non-inferiority, $-0.20$ & adjusted & 2.40 (0.05) & 2.41 (0.05) & 2.44 (0.05) & 2.42 (0.05) & 2.36 (0.05) & 2.33 (0.05)\\
$F=2$, $n=200$ & continuous & equivalence (lower limit), $-0.45$ & unadjusted & 5.02 (0.07) & 5.04 (0.07) & 5.07 (0.07) & 5.05 (0.07) & 5.02 (0.07) & 5.00 (0.07)\\
$F=2$, $n=200$ & continuous & equivalence (lower limit), $-0.45$ & adjusted & 5.00 (0.07) & 5.02 (0.07) & 5.04 (0.07) & 5.03 (0.07) & 5.03 (0.07) & 5.00 (0.07)\\
$F=2$, $n=200$ & continuous & equivalence (upper limit), $0.45$ & unadjusted & 5.04 (0.07) & 5.07 (0.07) & 5.10 (0.07) & 5.08 (0.07) & 5.08 (0.07) & 5.05 (0.07)\\
$F=2$, $n=200$ & continuous & equivalence (upper limit), $0.45$ & adjusted & 5.08 (0.07) & 5.11 (0.07) & 5.13 (0.07) & 5.12 (0.07) & 5.10 (0.07) & 5.06 (0.07)\\
$F=2$, $n=200$ & binary & superiority, $0.00$ & unadjusted & 4.26 (0.06) & 4.26 (0.06) & 4.26 (0.06) & 4.26 (0.06) & 4.92 (0.07) & 4.90 (0.07)\\
$F=2$, $n=200$ & binary & superiority, $0.00$ & adjusted & 4.88 (0.07) & 4.89 (0.07) & 4.94 (0.07) & 4.91 (0.07) & 5.01 (0.07) & 5.00 (0.07)\\
$F=2$, $n=200$ & binary & non-inferiority, $-0.10$ & unadjusted & 2.33 (0.05) & 2.33 (0.05) & 2.35 (0.05) & 2.34 (0.05) & 2.39 (0.05) & 2.38 (0.05)\\
$F=2$, $n=200$ & binary & non-inferiority, $-0.10$ & adjusted & 2.36 (0.05) & 2.36 (0.05) & 2.38 (0.05) & 2.37 (0.05) & 2.38 (0.05) & 2.38 (0.05)\\
$F=2$, $n=200$ & binary & equivalence (lower limit), $-0.21$ & unadjusted & 4.79 (0.07) & 4.79 (0.07) & 4.80 (0.07) & 4.80 (0.07) & 4.64 (0.07) & 4.63 (0.07)\\
$F=2$, $n=200$ & binary & equivalence (lower limit), $-0.21$ & adjusted & 4.76 (0.07) & 4.78 (0.07) & 4.80 (0.07) & 4.79 (0.07) & 4.75 (0.07) & 4.75 (0.07)\\
$F=2$, $n=200$ & binary & equivalence (upper limit), $0.21$ & unadjusted & 4.91 (0.07) & 4.91 (0.07) & 4.92 (0.07) & 4.92 (0.07) & 4.92 (0.07) & 4.91 (0.07)\\
$F=2$, $n=200$ & binary & equivalence (upper limit), $0.21$ & adjusted & 4.89 (0.07) & 4.90 (0.07) & 4.93 (0.07) & 4.91 (0.07) & 4.91 (0.07) & 4.90 (0.07)\\
$F=2$, $n=1000$ & continuous & superiority, $0.00$ & unadjusted & 5.00 (0.07) & 5.03 (0.07) & 5.05 (0.07) & 5.04 (0.07) & 5.02 (0.07) & 4.99 (0.07)\\
$F=2$, $n=1000$ & continuous & superiority, $0.00$ & adjusted & 5.02 (0.07) & 5.05 (0.07) & 5.06 (0.07) & 5.05 (0.07) & 5.07 (0.07) & 5.03 (0.07)\\
$F=2$, $n=1000$ & continuous & non-inferiority, $-0.20$ & unadjusted & 2.59 (0.05) & 2.62 (0.05) & 2.62 (0.05) & 2.62 (0.05) & 2.60 (0.05) & 2.58 (0.05)\\
$F=2$, $n=1000$ & continuous & non-inferiority, $-0.20$ & adjusted & 2.60 (0.05) & 2.61 (0.05) & 2.62 (0.05) & 2.62 (0.05) & 2.60 (0.05) & 2.59 (0.05)\\
$F=2$, $n=1000$ & continuous & equivalence (lower limit), $-0.45$ & unadjusted & 4.88 (0.07) & 4.91 (0.07) & 4.91 (0.07) & 4.91 (0.07) & 4.90 (0.07) & 4.88 (0.07)\\
$F=2$, $n=1000$ & continuous & equivalence (lower limit), $-0.45$ & adjusted & 4.85 (0.07) & 4.87 (0.07) & 4.88 (0.07) & 4.88 (0.07) & 4.87 (0.07) & 4.85 (0.07)\\
$F=2$, $n=1000$ & continuous & equivalence (upper limit), $0.45$ & unadjusted & 4.90 (0.07) & 4.93 (0.07) & 4.94 (0.07) & 4.93 (0.07) & 4.90 (0.07) & 4.87 (0.07)\\
$F=2$, $n=1000$ & continuous & equivalence (upper limit), $0.45$ & adjusted & 4.90 (0.07) & 4.93 (0.07) & 4.93 (0.07) & 4.92 (0.07) & 4.91 (0.07) & 4.88 (0.07)\\
$F=2$, $n=1000$ & binary & superiority, $0.00$ & unadjusted & 4.55 (0.07) & 4.55 (0.07) & 4.55 (0.07) & 4.55 (0.07) & 4.87 (0.07) & 4.85 (0.07)\\
$F=2$, $n=1000$ & binary & superiority, $0.00$ & adjusted & 4.81 (0.07) & 4.83 (0.07) & 4.84 (0.07) & 4.83 (0.07) & 4.89 (0.07) & 4.88 (0.07)\\
$F=2$, $n=1000$ & binary & non-inferiority, $-0.10$ & unadjusted & 2.46 (0.05) & 2.47 (0.05) & 2.47 (0.05) & 2.47 (0.05) & 2.52 (0.05) & 2.50 (0.05)\\
$F=2$, $n=1000$ & binary & non-inferiority, $-0.10$ & adjusted & 2.46 (0.05) & 2.46 (0.05) & 2.47 (0.05) & 2.47 (0.05) & 2.48 (0.05) & 2.47 (0.05)\\
$F=2$, $n=1000$ & binary & equivalence (lower limit), $-0.10$ & unadjusted & 4.91 (0.07) & 4.92 (0.07) & 4.92 (0.07) & 4.92 (0.07) & 4.89 (0.07) & 4.87 (0.07)\\
$F=2$, $n=1000$ & binary & equivalence (lower limit), $-0.10$ & adjusted & 4.96 (0.07) & 4.97 (0.07) & 4.97 (0.07) & 4.97 (0.07) & 4.95 (0.07) & 4.93 (0.07)\\
$F=2$, $n=1000$ & binary & equivalence (upper limit), $0.10$ & unadjusted & 5.01 (0.07) & 5.02 (0.07) & 5.03 (0.07) & 5.02 (0.07) & 5.07 (0.07) & 5.06 (0.07)\\
$F=2$, $n=1000$ & binary & equivalence (upper limit), $0.10$ & adjusted & 5.04 (0.07) & 5.05 (0.07) & 5.06 (0.07) & 5.06 (0.07) & 5.09 (0.07) & 5.07 (0.07)\\
$F=5$, $n=400$ & continuous & superiority, $0.00$ & unadjusted & 4.90 (0.07) & 4.92 (0.07) & 4.92 (0.07) & 4.92 (0.07) & 4.80 (0.07) & 4.78 (0.07)\\
$F=5$, $n=400$ & continuous & superiority, $0.00$ & adjusted & 4.91 (0.07) & 4.94 (0.07) & 4.94 (0.07) & 4.93 (0.07) & 4.92 (0.07) & 4.89 (0.07)\\
$F=5$, $n=400$ & continuous & non-inferiority, $-0.20$ & unadjusted & 2.57 (0.05) & 2.58 (0.05) & 2.58 (0.05) & 2.57 (0.05) & 2.46 (0.05) & 2.45 (0.05)\\
$F=5$, $n=400$ & continuous & non-inferiority, $-0.20$ & adjusted & 2.55 (0.05) & 2.56 (0.05) & 2.56 (0.05) & 2.55 (0.05) & 2.53 (0.05) & 2.52 (0.05)\\
$F=5$, $n=400$ & continuous & equivalence (lower limit), $-0.45$ & unadjusted & 4.92 (0.07) & 4.94 (0.07) & 4.94 (0.07) & 4.93 (0.07) & 4.85 (0.07) & 4.83 (0.07)\\
$F=5$, $n=400$ & continuous & equivalence (lower limit), $-0.45$ & adjusted & 4.89 (0.07) & 4.90 (0.07) & 4.91 (0.07) & 4.91 (0.07) & 4.92 (0.07) & 4.91 (0.07)\\
$F=5$, $n=400$ & continuous & equivalence (upper limit), $0.45$ & unadjusted & 5.10 (0.07) & 5.10 (0.07) & 5.12 (0.07) & 5.12 (0.07) & 5.01 (0.07) & 4.99 (0.07)\\
$F=5$, $n=400$ & continuous & equivalence (upper limit), $0.45$ & adjusted & 5.08 (0.07) & 5.11 (0.07) & 5.10 (0.07) & 5.09 (0.07) & 5.13 (0.07) & 5.11 (0.07)\\
$F=5$, $n=400$ & binary & superiority, $0.00$ & unadjusted & 4.49 (0.07) & 4.49 (0.07) & 4.49 (0.07) & 4.49 (0.07) & 5.06 (0.07) & 5.05 (0.07)\\
$F=5$, $n=400$ & binary & superiority, $0.00$ & adjusted & 5.09 (0.07) & 5.10 (0.07) & 5.12 (0.07) & 5.11 (0.07) & 5.09 (0.07) & 5.08 (0.07)\\
$F=5$, $n=400$ & binary & non-inferiority, $-0.10$ & unadjusted & 2.38 (0.05) & 2.38 (0.05) & 2.38 (0.05) & 2.38 (0.05) & 2.37 (0.05) & 2.37 (0.05)\\
$F=5$, $n=400$ & binary & non-inferiority, $-0.10$ & adjusted & 2.38 (0.05) & 2.39 (0.05) & 2.41 (0.05) & 2.40 (0.05) & 2.40 (0.05) & 2.39 (0.05)\\
$F=5$, $n=400$ & binary & equivalence (lower limit), $-0.15$ & unadjusted & 4.74 (0.07) & 4.75 (0.07) & 4.76 (0.07) & 4.75 (0.07) & 4.83 (0.07) & 4.82 (0.07)\\
$F=5$, $n=400$ & binary & equivalence (lower limit), $-0.15$ & adjusted & 4.78 (0.07) & 4.78 (0.07) & 4.80 (0.07) & 4.80 (0.07) & 4.80 (0.07) & 4.79 (0.07)\\
$F=5$, $n=400$ & binary & equivalence (upper limit), $0.15$ & unadjusted & 4.89 (0.07) & 4.89 (0.07) & 4.90 (0.07) & 4.89 (0.07) & 4.92 (0.07) & 4.91 (0.07)\\
$F=5$, $n=400$ & binary & equivalence (upper limit), $0.15$ & adjusted & 4.94 (0.07) & 4.94 (0.07) & 4.97 (0.07) & 4.96 (0.07) & 4.93 (0.07) & 4.92 (0.07)\\
$F=5$, $n=2000$ & continuous & superiority, $0.00$ & unadjusted & 4.96 (0.07) & 4.98 (0.07) & 4.99 (0.07) & 4.99 (0.07) & 4.95 (0.07) & 4.93 (0.07)\\
$F=5$, $n=2000$ & continuous & superiority, $0.00$ & adjusted & 4.98 (0.07) & 5.00 (0.07) & 5.00 (0.07) & 5.00 (0.07) & 5.01 (0.07) & 5.00 (0.07)\\
$F=5$, $n=2000$ & continuous & non-inferiority, $-0.20$ & unadjusted & 2.43 (0.05) & 2.44 (0.05) & 2.45 (0.05) & 2.44 (0.05) & 2.43 (0.05) & 2.42 (0.05)\\
$F=5$, $n=2000$ & continuous & non-inferiority, $-0.20$ & adjusted & 2.45 (0.05) & 2.46 (0.05) & 2.46 (0.05) & 2.46 (0.05) & 2.45 (0.05) & 2.44 (0.05)\\
$F=5$, $n=2000$ & continuous & equivalence (lower limit), $-0.45$ & unadjusted & 4.95 (0.07) & 4.96 (0.07) & 4.97 (0.07) & 4.97 (0.07) & 4.93 (0.07) & 4.91 (0.07)\\
$F=5$, $n=2000$ & continuous & equivalence (lower limit), $-0.45$ & adjusted & 4.95 (0.07) & 4.96 (0.07) & 4.97 (0.07) & 4.97 (0.07) & 4.97 (0.07) & 4.95 (0.07)\\
$F=5$, $n=2000$ & continuous & equivalence (upper limit), $0.45$ & unadjusted & 5.05 (0.07) & 5.07 (0.07) & 5.08 (0.07) & 5.08 (0.07) & 5.04 (0.07) & 5.03 (0.07)\\
$F=5$, $n=2000$ & continuous & equivalence (upper limit), $0.45$ & adjusted & 5.05 (0.07) & 5.07 (0.07) & 5.07 (0.07) & 5.07 (0.07) & 5.08 (0.07) & 5.06 (0.07)\\
$F=5$, $n=2000$ & binary & superiority, $0.00$ & unadjusted & 4.75 (0.07) & 4.75 (0.07) & 4.75 (0.07) & 4.75 (0.07) & 5.00 (0.07) & 4.99 (0.07)\\
$F=5$, $n=2000$ & binary & superiority, $0.00$ & adjusted & 4.99 (0.07) & 4.99 (0.07) & 5.00 (0.07) & 4.99 (0.07) & 5.00 (0.07) & 4.99 (0.07)\\
$F=5$, $n=2000$ & binary & non-inferiority, $-0.10$ & unadjusted & 2.47 (0.05) & 2.48 (0.05) & 2.48 (0.05) & 2.48 (0.05) & 2.46 (0.05) & 2.46 (0.05)\\
$F=5$, $n=2000$ & binary & non-inferiority, $-0.10$ & adjusted & 2.47 (0.05) & 2.47 (0.05) & 2.48 (0.05) & 2.48 (0.05) & 2.48 (0.05) & 2.48 (0.05)\\
$F=5$, $n=2000$ & binary & equivalence (lower limit), $-0.10$ & unadjusted & 4.71 (0.07) & 4.72 (0.07) & 4.72 (0.07) & 4.72 (0.07) & 4.70 (0.07) & 4.69 (0.07)\\
$F=5$, $n=2000$ & binary & equivalence (lower limit), $-0.10$ & adjusted & 4.73 (0.07) & 4.73 (0.07) & 4.74 (0.07) & 4.74 (0.07) & 4.76 (0.07) & 4.75 (0.07)\\
$F=5$, $n=2000$ & binary & equivalence (upper limit), $0.10$ & unadjusted & 5.01 (0.07) & 5.03 (0.07) & 5.04 (0.07) & 5.03 (0.07) & 5.05 (0.07) & 5.04 (0.07)\\
$F=5$, $n=2000$ & binary & equivalence (upper limit), $0.10$ & adjusted & 5.04 (0.07) & 5.04 (0.07) & 5.05 (0.07) & 5.05 (0.07) & 5.05 (0.07) & 5.04 (0.07)\\
\end{longtable}
\endgroup\end{singlespace}\end{landscape}

\begin{table}[!htbp]\centering
\caption{Variance diagnostics in the heterogeneous-effect null study: means over the 100,000 outer trials at the tested null value, averaged over the four null values.}
\label{tab:supp-hetero-diag}
\begin{threeparttable}\begin{singlespace}\footnotesize\setlength{\tabcolsep}{3pt}
\begin{tabular}{@{}lllrrrrrr@{}}\toprule
Design & Outcome & Score & $\widehat\lambda$ (model) & $\widehat\lambda$ (data) & $\widehat\lambda$ (adj.) & $\widehat V_{\rm R}/\widehat V_{\rm S}$ & $\widehat V^{*}/\widehat V_{\rm S}$ & $\widehat V^{*}/\widehat V_{\rm R}$\\\midrule
$F=2$, $n=200$ & continuous & unadjusted & 1.0017 & 1.0027 & 1.0016 & 1.0033 & 1.0019 & 0.9986\\
$F=2$, $n=200$ & continuous & adjusted & 1.0017 & 1.0027 & 1.0017 & 1.0034 & 1.0029 & 0.9995\\
$F=2$, $n=200$ & binary & unadjusted & 1.0007 & 1.0018 & 1.0007 & 1.0013 & 1.0009 & 0.9996\\
$F=2$, $n=200$ & binary & adjusted & 1.0007 & 1.0018 & 1.0007 & 1.0014 & 1.0012 & 0.9998\\
$F=2$, $n=1000$ & continuous & unadjusted & 1.0016 & 1.0018 & 1.0016 & 1.0033 & 1.0030 & 0.9997\\
$F=2$, $n=1000$ & continuous & adjusted & 1.0016 & 1.0018 & 1.0016 & 1.0033 & 1.0031 & 0.9999\\
$F=2$, $n=1000$ & binary & unadjusted & 1.0007 & 1.0009 & 1.0007 & 1.0014 & 1.0013 & 0.9999\\
$F=2$, $n=1000$ & binary & adjusted & 1.0007 & 1.0009 & 1.0007 & 1.0015 & 1.0014 & 0.9999\\
$F=5$, $n=400$ & continuous & unadjusted & 1.0009 & 1.0012 & 1.0010 & 1.0018 & 0.9900 & 0.9881\\
$F=5$, $n=400$ & continuous & adjusted & 1.0010 & 1.0012 & 1.0010 & 1.0019 & 1.0025 & 1.0006\\
$F=5$, $n=400$ & binary & unadjusted & 1.0004 & 1.0010 & 1.0008 & 1.0007 & 0.9984 & 0.9977\\
$F=5$, $n=400$ & binary & adjusted & 1.0004 & 1.0010 & 1.0008 & 1.0007 & 1.0013 & 1.0006\\
$F=5$, $n=2000$ & continuous & unadjusted & 1.0009 & 1.0012 & 1.0010 & 1.0018 & 0.9991 & 0.9973\\
$F=5$, $n=2000$ & continuous & adjusted & 1.0009 & 1.0012 & 1.0010 & 1.0018 & 1.0019 & 1.0002\\
$F=5$, $n=2000$ & binary & unadjusted & 1.0004 & 1.0007 & 1.0005 & 1.0008 & 1.0003 & 0.9995\\
$F=5$, $n=2000$ & binary & adjusted & 1.0004 & 1.0007 & 1.0005 & 1.0008 & 1.0009 & 1.0001\\
\bottomrule\end{tabular}
\begin{tablenotes}[flushleft]\small\item $\widehat V^{*}$ is the sample variance of the 4,999 regenerated statistics within a trial. The displayed ratios are averages of trial-level ratios, subsequently averaged over the four null settings. In ratio columns, $\widehat V_{\rm R}$ uses the model-specified effect deviations; the three columns of $\widehat\lambda$ distinguish the correction versions. These are not per-trial accuracy bounds. The floor in the reported variance calculations was not active.\end{tablenotes}
\end{singlespace}\end{threeparttable}\end{table}

\clearpage
\begin{landscape}\begin{singlespace}\begingroup\footnotesize\setlength{\tabcolsep}{3.3pt}\renewcommand{\arraystretch}{1.0}\setlength{\LTcapwidth}{\linewidth}\setlength{\LTleft}{\fill}\setlength{\LTright}{\fill}
\begin{longtable}{@{}llllrrrrrr@{}}
\caption{Rejection probabilities (\%) at the null values under heterogeneous treatment effects and a strong biased coin ($p_{\rm bc}=.95$), with Monte Carlo standard errors in parentheses.}
\label{tab:supp-hetero-null-stress}\\
\toprule
Design & Outcome & Null value & Score & RT & CRT (model) & CRT (data) & CRT (data, adj.) & SIGA-S & SIGA-R\\
\midrule\endfirsthead
\multicolumn{10}{c}{\tablename\ \thetable{} -- continued}\\\toprule
Design & Outcome & Null value & Score & RT & CRT (model) & CRT (data) & CRT (data, adj.) & SIGA-S & SIGA-R\\
\midrule\endhead
\midrule\multicolumn{10}{r}{Continued on next page}\\\endfoot
\bottomrule\endlastfoot
$F=2$, $n=200$ & continuous & superiority, $0.00$ & unadjusted & 4.43 (0.07) & 5.04 (0.07) & 5.04 (0.07) & 5.00 (0.07) & 4.79 (0.07) & 4.20 (0.06)\\
$F=2$, $n=200$ & continuous & superiority, $0.00$ & adjusted & 4.35 (0.06) & 5.02 (0.07) & 5.02 (0.07) & 4.98 (0.07) & 4.86 (0.07) & 4.26 (0.06)\\
$F=2$, $n=200$ & continuous & non-inferiority, $-0.20$ & unadjusted & 2.11 (0.05) & 2.44 (0.05) & 2.44 (0.05) & 2.42 (0.05) & 2.32 (0.05) & 2.03 (0.04)\\
$F=2$, $n=200$ & continuous & non-inferiority, $-0.20$ & adjusted & 2.14 (0.05) & 2.46 (0.05) & 2.45 (0.05) & 2.43 (0.05) & 2.44 (0.05) & 2.12 (0.05)\\
$F=2$, $n=200$ & continuous & equivalence (lower limit), $-0.45$ & unadjusted & 4.53 (0.07) & 4.99 (0.07) & 5.01 (0.07) & 4.98 (0.07) & 4.85 (0.07) & 4.40 (0.06)\\
$F=2$, $n=200$ & continuous & equivalence (lower limit), $-0.45$ & adjusted & 4.50 (0.07) & 5.01 (0.07) & 5.02 (0.07) & 4.99 (0.07) & 4.90 (0.07) & 4.44 (0.07)\\
$F=2$, $n=200$ & continuous & equivalence (upper limit), $0.45$ & unadjusted & 4.56 (0.07) & 5.01 (0.07) & 5.02 (0.07) & 4.99 (0.07) & 4.93 (0.07) & 4.46 (0.07)\\
$F=2$, $n=200$ & continuous & equivalence (upper limit), $0.45$ & adjusted & 4.54 (0.07) & 5.02 (0.07) & 5.04 (0.07) & 5.01 (0.07) & 4.94 (0.07) & 4.48 (0.07)\\
$F=2$, $n=200$ & binary & superiority, $0.00$ & unadjusted & 3.67 (0.06) & 3.67 (0.06) & 3.67 (0.06) & 3.67 (0.06) & 4.82 (0.07) & 4.80 (0.07)\\
$F=2$, $n=200$ & binary & superiority, $0.00$ & adjusted & 4.31 (0.06) & 4.34 (0.06) & 4.66 (0.07) & 4.58 (0.07) & 4.88 (0.07) & 4.85 (0.07)\\
$F=2$, $n=200$ & binary & non-inferiority, $-0.10$ & unadjusted & 2.27 (0.05) & 2.27 (0.05) & 2.47 (0.05) & 2.39 (0.05) & 2.36 (0.05) & 2.35 (0.05)\\
$F=2$, $n=200$ & binary & non-inferiority, $-0.10$ & adjusted & 2.30 (0.05) & 2.31 (0.05) & 2.57 (0.05) & 2.47 (0.05) & 2.35 (0.05) & 2.35 (0.05)\\
$F=2$, $n=200$ & binary & equivalence (lower limit), $-0.21$ & unadjusted & 4.77 (0.07) & 4.78 (0.07) & 5.08 (0.07) & 4.94 (0.07) & 4.55 (0.07) & 4.53 (0.07)\\
$F=2$, $n=200$ & binary & equivalence (lower limit), $-0.21$ & adjusted & 4.76 (0.07) & 4.77 (0.07) & 5.06 (0.07) & 4.94 (0.07) & 4.68 (0.07) & 4.67 (0.07)\\
$F=2$, $n=200$ & binary & equivalence (upper limit), $0.21$ & unadjusted & 5.04 (0.07) & 5.04 (0.07) & 5.34 (0.07) & 5.21 (0.07) & 5.02 (0.07) & 5.01 (0.07)\\
$F=2$, $n=200$ & binary & equivalence (upper limit), $0.21$ & adjusted & 5.04 (0.07) & 5.05 (0.07) & 5.34 (0.07) & 5.24 (0.07) & 5.04 (0.07) & 5.02 (0.07)\\
$F=2$, $n=1000$ & continuous & superiority, $0.00$ & unadjusted & 4.31 (0.06) & 4.95 (0.07) & 4.95 (0.07) & 4.95 (0.07) & 4.95 (0.07) & 4.33 (0.06)\\
$F=2$, $n=1000$ & continuous & superiority, $0.00$ & adjusted & 4.37 (0.06) & 4.95 (0.07) & 4.94 (0.07) & 4.93 (0.07) & 4.98 (0.07) & 4.38 (0.06)\\
$F=2$, $n=1000$ & continuous & non-inferiority, $-0.20$ & unadjusted & 2.13 (0.05) & 2.44 (0.05) & 2.44 (0.05) & 2.44 (0.05) & 2.41 (0.05) & 2.12 (0.05)\\
$F=2$, $n=1000$ & continuous & non-inferiority, $-0.20$ & adjusted & 2.16 (0.05) & 2.45 (0.05) & 2.45 (0.05) & 2.45 (0.05) & 2.45 (0.05) & 2.14 (0.05)\\
$F=2$, $n=1000$ & continuous & equivalence (lower limit), $-0.45$ & unadjusted & 4.54 (0.07) & 4.98 (0.07) & 4.98 (0.07) & 4.97 (0.07) & 4.94 (0.07) & 4.50 (0.07)\\
$F=2$, $n=1000$ & continuous & equivalence (lower limit), $-0.45$ & adjusted & 4.57 (0.07) & 5.01 (0.07) & 5.01 (0.07) & 5.01 (0.07) & 5.01 (0.07) & 4.57 (0.07)\\
$F=2$, $n=1000$ & continuous & equivalence (upper limit), $0.45$ & unadjusted & 4.59 (0.07) & 5.03 (0.07) & 5.03 (0.07) & 5.03 (0.07) & 5.04 (0.07) & 4.61 (0.07)\\
$F=2$, $n=1000$ & continuous & equivalence (upper limit), $0.45$ & adjusted & 4.63 (0.07) & 5.04 (0.07) & 5.04 (0.07) & 5.04 (0.07) & 5.04 (0.07) & 4.64 (0.07)\\
$F=2$, $n=1000$ & binary & superiority, $0.00$ & unadjusted & 4.41 (0.06) & 4.41 (0.06) & 4.41 (0.06) & 4.41 (0.06) & 4.93 (0.07) & 4.92 (0.07)\\
$F=2$, $n=1000$ & binary & superiority, $0.00$ & adjusted & 4.71 (0.07) & 4.73 (0.07) & 4.83 (0.07) & 4.82 (0.07) & 4.95 (0.07) & 4.94 (0.07)\\
$F=2$, $n=1000$ & binary & non-inferiority, $-0.10$ & unadjusted & 2.37 (0.05) & 2.37 (0.05) & 2.39 (0.05) & 2.38 (0.05) & 2.42 (0.05) & 2.41 (0.05)\\
$F=2$, $n=1000$ & binary & non-inferiority, $-0.10$ & adjusted & 2.39 (0.05) & 2.39 (0.05) & 2.44 (0.05) & 2.42 (0.05) & 2.42 (0.05) & 2.40 (0.05)\\
$F=2$, $n=1000$ & binary & equivalence (lower limit), $-0.10$ & unadjusted & 4.91 (0.07) & 4.91 (0.07) & 4.92 (0.07) & 4.92 (0.07) & 4.98 (0.07) & 4.94 (0.07)\\
$F=2$, $n=1000$ & binary & equivalence (lower limit), $-0.10$ & adjusted & 4.94 (0.07) & 4.95 (0.07) & 4.99 (0.07) & 4.97 (0.07) & 4.98 (0.07) & 4.97 (0.07)\\
$F=2$, $n=1000$ & binary & equivalence (upper limit), $0.10$ & unadjusted & 4.97 (0.07) & 4.98 (0.07) & 4.99 (0.07) & 4.98 (0.07) & 5.05 (0.07) & 5.03 (0.07)\\
$F=2$, $n=1000$ & binary & equivalence (upper limit), $0.10$ & adjusted & 5.00 (0.07) & 5.00 (0.07) & 5.04 (0.07) & 5.03 (0.07) & 5.02 (0.07) & 5.01 (0.07)\\
$F=5$, $n=400$ & continuous & superiority, $0.00$ & unadjusted & 4.30 (0.06) & 4.99 (0.07) & 4.87 (0.07) & 4.86 (0.07) & 4.21 (0.06) & 3.63 (0.06)\\
$F=5$, $n=400$ & continuous & superiority, $0.00$ & adjusted & 4.25 (0.06) & 4.99 (0.07) & 4.85 (0.07) & 4.84 (0.07) & 4.94 (0.07) & 4.21 (0.06)\\
$F=5$, $n=400$ & continuous & non-inferiority, $-0.20$ & unadjusted & 2.18 (0.05) & 2.53 (0.05) & 2.47 (0.05) & 2.46 (0.05) & 2.14 (0.05) & 1.81 (0.04)\\
$F=5$, $n=400$ & continuous & non-inferiority, $-0.20$ & adjusted & 2.20 (0.05) & 2.56 (0.05) & 2.50 (0.05) & 2.49 (0.05) & 2.53 (0.05) & 2.18 (0.05)\\
$F=5$, $n=400$ & continuous & equivalence (lower limit), $-0.45$ & unadjusted & 4.43 (0.07) & 4.93 (0.07) & 4.87 (0.07) & 4.86 (0.07) & 4.35 (0.06) & 3.92 (0.06)\\
$F=5$, $n=400$ & continuous & equivalence (lower limit), $-0.45$ & adjusted & 4.44 (0.07) & 5.00 (0.07) & 4.90 (0.07) & 4.89 (0.07) & 4.94 (0.07) & 4.39 (0.06)\\
$F=5$, $n=400$ & continuous & equivalence (upper limit), $0.45$ & unadjusted & 4.42 (0.06) & 4.93 (0.07) & 4.86 (0.07) & 4.84 (0.07) & 4.35 (0.06) & 3.91 (0.06)\\
$F=5$, $n=400$ & continuous & equivalence (upper limit), $0.45$ & adjusted & 4.41 (0.06) & 5.00 (0.07) & 4.90 (0.07) & 4.89 (0.07) & 4.96 (0.07) & 4.40 (0.06)\\
$F=5$, $n=400$ & binary & superiority, $0.00$ & unadjusted & 4.13 (0.06) & 4.13 (0.06) & 4.13 (0.06) & 4.13 (0.06) & 5.00 (0.07) & 4.98 (0.07)\\
$F=5$, $n=400$ & binary & superiority, $0.00$ & adjusted & 4.95 (0.07) & 4.98 (0.07) & 5.24 (0.07) & 5.16 (0.07) & 5.03 (0.07) & 5.00 (0.07)\\
$F=5$, $n=400$ & binary & non-inferiority, $-0.10$ & unadjusted & 2.28 (0.05) & 2.28 (0.05) & 2.33 (0.05) & 2.31 (0.05) & 2.33 (0.05) & 2.31 (0.05)\\
$F=5$, $n=400$ & binary & non-inferiority, $-0.10$ & adjusted & 2.35 (0.05) & 2.36 (0.05) & 2.44 (0.05) & 2.42 (0.05) & 2.34 (0.05) & 2.33 (0.05)\\
$F=5$, $n=400$ & binary & equivalence (lower limit), $-0.15$ & unadjusted & 4.66 (0.07) & 4.66 (0.07) & 4.74 (0.07) & 4.70 (0.07) & 4.79 (0.07) & 4.76 (0.07)\\
$F=5$, $n=400$ & binary & equivalence (lower limit), $-0.15$ & adjusted & 4.70 (0.07) & 4.71 (0.07) & 4.80 (0.07) & 4.76 (0.07) & 4.73 (0.07) & 4.71 (0.07)\\
$F=5$, $n=400$ & binary & equivalence (upper limit), $0.15$ & unadjusted & 4.86 (0.07) & 4.86 (0.07) & 4.92 (0.07) & 4.89 (0.07) & 4.91 (0.07) & 4.90 (0.07)\\
$F=5$, $n=400$ & binary & equivalence (upper limit), $0.15$ & adjusted & 4.88 (0.07) & 4.89 (0.07) & 5.01 (0.07) & 4.95 (0.07) & 4.91 (0.07) & 4.90 (0.07)\\
$F=5$, $n=2000$ & continuous & superiority, $0.00$ & unadjusted & 4.25 (0.06) & 5.00 (0.07) & 4.99 (0.07) & 4.99 (0.07) & 4.84 (0.07) & 4.09 (0.06)\\
$F=5$, $n=2000$ & continuous & superiority, $0.00$ & adjusted & 4.29 (0.06) & 5.00 (0.07) & 4.98 (0.07) & 4.98 (0.07) & 4.99 (0.07) & 4.26 (0.06)\\
$F=5$, $n=2000$ & continuous & non-inferiority, $-0.20$ & unadjusted & 2.14 (0.05) & 2.51 (0.05) & 2.50 (0.05) & 2.50 (0.05) & 2.41 (0.05) & 2.04 (0.04)\\
$F=5$, $n=2000$ & continuous & non-inferiority, $-0.20$ & adjusted & 2.12 (0.05) & 2.50 (0.05) & 2.50 (0.05) & 2.50 (0.05) & 2.51 (0.05) & 2.12 (0.05)\\
$F=5$, $n=2000$ & continuous & equivalence (lower limit), $-0.45$ & unadjusted & 4.42 (0.06) & 4.97 (0.07) & 4.96 (0.07) & 4.96 (0.07) & 4.82 (0.07) & 4.29 (0.06)\\
$F=5$, $n=2000$ & continuous & equivalence (lower limit), $-0.45$ & adjusted & 4.44 (0.07) & 4.97 (0.07) & 4.97 (0.07) & 4.97 (0.07) & 4.98 (0.07) & 4.42 (0.06)\\
$F=5$, $n=2000$ & continuous & equivalence (upper limit), $0.45$ & unadjusted & 4.49 (0.07) & 5.03 (0.07) & 5.03 (0.07) & 5.03 (0.07) & 4.89 (0.07) & 4.37 (0.06)\\
$F=5$, $n=2000$ & continuous & equivalence (upper limit), $0.45$ & adjusted & 4.50 (0.07) & 5.04 (0.07) & 5.04 (0.07) & 5.04 (0.07) & 5.03 (0.07) & 4.47 (0.07)\\
$F=5$, $n=2000$ & binary & superiority, $0.00$ & unadjusted & 4.68 (0.07) & 4.68 (0.07) & 4.68 (0.07) & 4.68 (0.07) & 5.07 (0.07) & 5.06 (0.07)\\
$F=5$, $n=2000$ & binary & superiority, $0.00$ & adjusted & 5.04 (0.07) & 5.08 (0.07) & 5.13 (0.07) & 5.11 (0.07) & 5.09 (0.07) & 5.06 (0.07)\\
$F=5$, $n=2000$ & binary & non-inferiority, $-0.10$ & unadjusted & 2.55 (0.05) & 2.55 (0.05) & 2.56 (0.05) & 2.55 (0.05) & 2.57 (0.05) & 2.56 (0.05)\\
$F=5$, $n=2000$ & binary & non-inferiority, $-0.10$ & adjusted & 2.59 (0.05) & 2.59 (0.05) & 2.60 (0.05) & 2.60 (0.05) & 2.58 (0.05) & 2.57 (0.05)\\
$F=5$, $n=2000$ & binary & equivalence (lower limit), $-0.10$ & unadjusted & 4.83 (0.07) & 4.85 (0.07) & 4.85 (0.07) & 4.85 (0.07) & 4.84 (0.07) & 4.83 (0.07)\\
$F=5$, $n=2000$ & binary & equivalence (lower limit), $-0.10$ & adjusted & 4.88 (0.07) & 4.90 (0.07) & 4.93 (0.07) & 4.92 (0.07) & 4.90 (0.07) & 4.88 (0.07)\\
$F=5$, $n=2000$ & binary & equivalence (upper limit), $0.10$ & unadjusted & 4.97 (0.07) & 4.98 (0.07) & 4.98 (0.07) & 4.98 (0.07) & 5.02 (0.07) & 5.00 (0.07)\\
$F=5$, $n=2000$ & binary & equivalence (upper limit), $0.10$ & adjusted & 5.01 (0.07) & 5.02 (0.07) & 5.05 (0.07) & 5.04 (0.07) & 5.02 (0.07) & 5.00 (0.07)\\
\end{longtable}
\endgroup\end{singlespace}\end{landscape}

\begin{table}[!htbp]\centering
\caption{Variance diagnostics in the heterogeneous-effect null study under a strong biased coin ($p_{\rm bc}=.95$): means over the 100,000 outer trials at the tested null value, averaged over the four null values.}
\label{tab:supp-hetero-diag-stress}
\begin{threeparttable}\begin{singlespace}\footnotesize\setlength{\tabcolsep}{3pt}
\begin{tabular}{@{}lllrrrrrr@{}}\toprule
Design & Outcome & Score & $\widehat\lambda$ (model) & $\widehat\lambda$ (data) & $\widehat\lambda$ (adj.) & $\widehat V_{\rm R}/\widehat V_{\rm S}$ & $\widehat V^{*}/\widehat V_{\rm S}$ & $\widehat V^{*}/\widehat V_{\rm R}$\\\midrule
$F=2$, $n=200$ & continuous & unadjusted & 1.0281 & 1.0295 & 1.0280 & 1.0571 & 1.0416 & 0.9853\\
$F=2$, $n=200$ & continuous & adjusted & 1.0294 & 1.0308 & 1.0293 & 1.0597 & 1.0494 & 0.9902\\
$F=2$, $n=200$ & binary & unadjusted & 1.0009 & 1.0083 & 1.0008 & 1.0017 & 1.0015 & 0.9998\\
$F=2$, $n=200$ & binary & adjusted & 1.0009 & 1.0083 & 1.0008 & 1.0017 & 1.0018 & 1.0000\\
$F=2$, $n=1000$ & continuous & unadjusted & 1.0273 & 1.0276 & 1.0273 & 1.0553 & 1.0547 & 0.9994\\
$F=2$, $n=1000$ & continuous & adjusted & 1.0275 & 1.0278 & 1.0275 & 1.0558 & 1.0563 & 1.0005\\
$F=2$, $n=1000$ & binary & unadjusted & 1.0009 & 1.0023 & 1.0009 & 1.0018 & 1.0015 & 0.9997\\
$F=2$, $n=1000$ & binary & adjusted & 1.0009 & 1.0023 & 1.0009 & 1.0018 & 1.0015 & 0.9997\\
$F=5$, $n=400$ & continuous & unadjusted & 1.0307 & 1.0261 & 1.0257 & 1.0622 & 0.9908 & 0.9327\\
$F=5$, $n=400$ & continuous & adjusted & 1.0345 & 1.0294 & 1.0288 & 1.0702 & 1.0665 & 0.9965\\
$F=5$, $n=400$ & binary & unadjusted & 1.0010 & 1.0042 & 1.0016 & 1.0020 & 0.9995 & 0.9974\\
$F=5$, $n=400$ & binary & adjusted & 1.0010 & 1.0042 & 1.0016 & 1.0020 & 1.0026 & 1.0006\\
$F=5$, $n=2000$ & continuous & unadjusted & 1.0342 & 1.0339 & 1.0338 & 1.0697 & 1.0522 & 0.9837\\
$F=5$, $n=2000$ & continuous & adjusted & 1.0352 & 1.0349 & 1.0347 & 1.0715 & 1.0706 & 0.9991\\
$F=5$, $n=2000$ & binary & unadjusted & 1.0011 & 1.0020 & 1.0013 & 1.0023 & 1.0017 & 0.9994\\
$F=5$, $n=2000$ & binary & adjusted & 1.0011 & 1.0020 & 1.0013 & 1.0023 & 1.0023 & 1.0000\\
\bottomrule\end{tabular}
\begin{tablenotes}[flushleft]\small\item $\widehat V^{*}$ is the sample variance of the 4,999 regenerated statistics within a trial. The displayed ratios are averages of trial-level ratios, subsequently averaged over the four null settings. In ratio columns, $\widehat V_{\rm R}$ uses the model-specified effect deviations; the three columns of $\widehat\lambda$ distinguish the correction versions. These are not per-trial accuracy bounds. The floor in the reported variance calculations was not active.\end{tablenotes}
\end{singlespace}\end{threeparttable}\end{table}


% BEGIN GENERATED VALIDATION APPENDIX
\section{Paired evaluation of the covariance reconstruction}
This study compares the original and reconstructed covariance estimators on common trials, separately from the earlier high-replication studies in Appendices E--J. It comprises 280 scenarios: 224 broad-validation scenarios, 40 factorial scenarios and 16 sharp-null controls, each with 10,000 outcome trials. The two score analyses use the same trials, giving 2,800,000 outcome trials rather than 5,600,000 independent replications. Each final run retains its first 2,000 pilot trials and adds trials 2000--9999. Thus the reported results extend the pilot; they are not a new independent 10,000-trial experiment. No scenario was removed when increasing replication. The correlated-profile suite available in the code was not run.

\subsection{Design, estimators and reproducibility}
The broad family crosses $(F,n)=(2,200),(2,1000),(5,400),(5,2000)$, continuous and binary outcomes, two fixed effect directions, seven testing roles and two configurations. Factor probabilities are the independent laws used in the main study. The configurations are $p_{\rm bc}=.80$, continuous residual standard deviation one, individual-effect standard deviation $.25$ and maximum stratum-effect deviation $.5$; and $p_{\rm bc}=.95$, residual standard deviation $.25$, individual-effect standard deviation $.10$ and maximum deviation one. Binary deviations have maximum $.15$ in both configurations. The effect directions are the first binary factor and the product of all sign-coded factors, each centred by its profile-probability-weighted mean and scaled to maximum absolute value one. Neither direction is selected using covariance eigenvalues or rejection rates.

Let $x(s)$ be the $0/1$ factor vector, $c=(.5,.4,.3,.2,.1)$ and $a=(.35,-.25,.20,-.15,.10)$, restricted to the first $F$ components. The continuous control mean is the centred version of $c^\top x(s)+.25x_1(s)x_2(s)$. Outcomes follow
\[
 Y_i(0)=\mu_{S_i}+\sigma\varepsilon_i,\qquad
 Y_i(1)=Y_i(0)+\Delta+d_{S_i}+\tau\eta_i,
\]
where the two errors are independent standard normals. Binary control risks are
\[
 p_{0s}=\operatorname{expit}\{a_0+a^\top x(s)+.20x_1(s)x_2(s)\},
 \qquad \sum_s\pi_s p_{0s}=.60,
\]
and treatment risks are $p_{0s}+\Delta+d_s$. If needed, the complete centred deviation vector is multiplied successively by $.9$ before simulation until all treatment risks lie in $(.02,.98)$. This preserves the marginal effect. Requested and actual amplitudes and the feasibility multiplier are retained in every summary. The baseline score regression includes the intercept and factor main effects but omits the interaction.

The seven roles are superiority at its null and at an alternative, non-inferiority at its null and at an alternative, and equivalence at its two boundaries and at an interior alternative. Superiority is two-sided at 5\%; non-inferiority is upper-tailed at 2.5\%; equivalence uses the intersection of two one-sided tests at 5\%. In the order of the four design--sample-size combinations above, the continuous superiority alternatives are $(.430,.190,.300,.135)$, non-inferiority alternatives are $(.230,-.010,.100,-.065)$ with boundary $-.20$, and equivalence alternatives are $(0,.280,.180,.330)$ with margins $.45$. The binary counterparts are $(.20,.10,.14,.07)$, $(.10,0,.05,-.03)$ with boundary $-.10$, and $(0,0,0,.04)$ with margins $(.21,.10,.15,.10)$. These fixed values were inherited from the distributed validation plan, not chosen from the present results.

The factorial family fixes $F=5,n=400$ and $\tau=.25$, crosses $p_{\rm bc}\in\{.80,.95\}$ and $\sigma\in\{.25,1\}$, and uses zero deviations or deviations of amplitude $.5$ or one in either fixed direction. It evaluates superiority at $\Delta=0$ and $.30$, giving 40 scenarios. The controls cross the four designs, both coin probabilities and both outcome types at the sharp null $Y_i(1)=Y_i(0)$; continuous controls use $\sigma=1$, $\tau=0$.

The saved run plans specify 20,000 independent three-sequence allocation replicates for covariance estimation and 999 regenerated assignments per outcome trial. The first sequence estimates $\Gamma_n$ and the covariance of $\mathsf L D_n$; the average of two product covariance estimates gives $\widehat\Psi_n$. Calibration is shared across scenarios with the same allocation design, sample size and profile law. Random streams are keyed separately for calibration, profiles, original assignments, outcomes and reference regeneration. The reconstruction and scenario families were fixed after the developmental diagnosis, before the 2,000-trial pilot. Replication was then increased for every scenario without changing the methods, generating models, master seed, reference-draw count or allocation-covariance estimates. This is post-development evaluation with increased Monte Carlo precision, not preregistration before the initial diagnosis.

All primary analyses estimate $d_s(b)$ from observed arm means using Appendix~C.2, assigning zero if either arm is absent. The analysis interface receives no generating treatment effect, potential outcome, or true outcome variance. Known design-stage profile probabilities are common inputs to both constructions. The raw and reconstructed methods use the same score, effect estimate, product covariance, outcome trials and reference assignments. The optional diagonal noise adjustment is reported separately as exploratory. The computational output names \texttt{original} and \texttt{refined} correspond respectively to the original and reconstructed matrices in the paper; they are not new method names.

Thirteen computational columns are retained: RT, and the original/reconstructed versions of normal SIGA-S, normal SIGA-R, simulated data-based CRT, lattice-enabled SIGA-R, lattice-enabled CRT approximation, and exploratory noise-adjusted simulated CRT. The lattice is used only for unadjusted binary superiority at $b=0$. Elsewhere the lattice-enabled R and CRT approximations use normal R and S tails, respectively. The corrected lattice calculation applies the rescaled rejection region and the original-scale tie clause to probabilities computed using the reference variance. Exact ties are handled numerically with an absolute tolerance of $10^{-12}$. The lattice and normal-only columns are all retained rather than selecting the better tail approximation after simulation.

For each scenario and method, the rejection proportion is accompanied by its binomial Monte Carlo standard error and Wilson interval. For methods $A,B$, the paired difference is $\bar D$ with $D_i=\mathbf1(A\text{ rejects})-\mathbf1(B\text{ rejects})$ and estimated standard error $s_D/\sqrt{10000}$. The summaries retain both directional discordance counts. These are scenario-specific descriptive intervals, not simultaneous guarantees. The outer-trial standard errors condition on the common estimated allocation covariances. The reference distribution itself is estimated with 999 assignments, so differences from RT contain inner Monte Carlo error. The new summaries do not contain separate construction-time benchmarks.

The saved plans specify master seed \texttt{2026100319} and identify the analysis source by its SHA256 signature. The extension records identify the source and target plan hashes, reused trial indices and unchanged settings. The plans, complete aggregate outputs and source code accompany the computational record. Their hashes and reported arithmetic were checked; the trial-level checkpoint files and actual allocation-covariance arrays were not part of the supplied aggregate package. The two scores and methods share trials, and scenarios sharing an allocation design share covariance estimates; counts of improved scenarios below are descriptive summaries, not independent binomial comparisons.

\subsection{Results and interpretation}
All 280 scenarios completed 10,000 trials. The full computational record retains 7,280 rejection-rate rows, 7,840 paired comparisons and 1,504 variance-diagnostic rows. Reported counts, marginal and paired standard errors, and the between-stratum, within-stratum and cross-term decomposition of reference variance are arithmetically consistent. No zero reference variance, variance floor or rank fallback was recorded. Machine-readable results contain every method, scenario and score, including alternatives and comparisons that did not improve.

The main-article table reports discrepancies between mean estimated and mean empirical reference variances; Table~\ref{tab:supp-pilot-paired} reports rejection-rate differences from the Monte Carlo reference test. In the unadjusted factorial settings, the median absolute relative variance discrepancy decreased from 1.67\% to 0.47\%, and the maximum from 8.63\% to 1.78\%. It decreased in 38 of 40 configurations. With baseline adjustment the corresponding medians were 0.43\% and 0.49\%, and maxima 1.87\% and 1.80\%. Thus the largest gains occurred for the unadjusted scores motivating the reconstruction, rather than uniformly across scores. Across the unadjusted factorial settings, the median absolute rejection-rate difference from RT decreased from 0.20 to 0.06 percentage points and the maximum from 0.78 to 0.37 points.

Table~\ref{tab:supp-pilot-target} gives all null roles for the motivating five-factor continuous configuration. At superiority, the mean trial-level empirical-reference-to-estimated-variance ratio changed from 0.9315 (Monte Carlo standard error 0.0009) to 0.9949 (0.0010). The ratio of mean estimated sampling variance to empirical variance of the observed statistic changed from 1.0914 to 1.0173; this comparison is meaningful because the generating effect equals the tested value. The SIGA-R rejection rate increased by 0.66 percentage points (paired standard error 0.08), while the simulated corrected-test rate changed by 0.03 points. Approximation of the original test and correction of its weak-null size remain distinct targets.

Improved mean variance estimation did not always improve reference-tail approximation. In the unadjusted factorial power setting with $p_{\rm bc}=.80$, $\sigma=1$, first-factor deviation amplitude $.5$ and $\Delta=.30$, RT power was 80.42\%, original SIGA-R power 80.27\% and reconstructed SIGA-R power 80.59\%. Their paired differences from RT were -0.15 (0.12) and +0.17 (0.13) percentage points. The absolute relative variance discrepancy decreased from 1.25\% to 0.09\%. This is the same illustrative configuration examined in the pilot, not a newly selected favourable case. All other increases and decreases in discrepancy remain in the complete summaries.

In all 32 control scenario--score combinations, the original and reconstructed simulated CRT had identical rejection decisions, as shown by zero discordance counts. The largest change in normal SIGA-S rejection rate was 0.12 percentage points. This does not establish sharp-null exactness: the two CRT versions can differ from RT. Table~\ref{tab:supp-pilot-null} reports all null-rate ranges. A 10,000-trial estimate has Monte Carlo standard error about 0.22 percentage points at a true rejection probability of 5\%, and about 0.16 at 2.5\%; paired standard errors are used for method differences. These are conditional on the shared allocation estimates and are not uniform error guarantees across the scenario collection.

For unadjusted binary superiority, the original-construction lattice approximation reduced the absolute difference from RT in 38 of 40 scenarios compared with ordinary Gaussian tails. In the two-factor, $n=200$, $p_{\rm bc}=.95$ binary control, RT rejected 3.70\%, normal SIGA-R 4.40\% and the lattice approximation 3.14\%. The reconstructed lattice result was also 3.14\%. The mixture therefore did not reproduce every finite-sample reference tail even when it retained the possible score values. Its weights and variance do not determine exact probabilities conditional on each realised profile sequence. Both normal-only and lattice-enabled columns are retained throughout.

These simulations provide no additional empirical claim for correlated profiles, no new reconstruction-time benchmark, and no proof of universal finite-sample dominance. The earlier high-replication results are not reinterpreted as evaluations of the reconstructed covariance. The 10,000-trial results replace the pilot percentages in this appendix, while the pilot and extension records remain archived separately.

\clearpage
\begin{table}[!htbp]\centering
\caption{Absolute differences from the Monte Carlo reference rejection rate in the 10,000-trial evaluation (percentage points).}
\label{tab:supp-pilot-paired}
\begin{threeparttable}\small
\begin{tabular}{@{}lccrrrr@{}}
\toprule
Family & $F$ & Score & \multicolumn{2}{c}{Median} & \multicolumn{2}{c}{Maximum}\\
 & & & Original & Reconstructed & Original & Reconstructed\\
\midrule
Factorial & 5 & U & 0.20 & 0.06 & 0.78 & 0.37\\
Factorial & 5 & A & 0.06 & 0.06 & 0.29 & 0.27\\
Broad & 2 & U & 0.11 & 0.10 & 1.24 & 1.18\\
Broad & 2 & A & 0.10 & 0.10 & 1.36 & 1.36\\
Broad & 5 & U & 0.10 & 0.08 & 0.69 & 0.61\\
Broad & 5 & A & 0.06 & 0.07 & 0.46 & 0.46\\
Controls & 2 & U & 0.13 & 0.13 & 0.56 & 0.56\\
Controls & 2 & A & 0.09 & 0.09 & 0.46 & 0.46\\
Controls & 5 & U & 0.10 & 0.10 & 0.38 & 0.33\\
Controls & 5 & A & 0.09 & 0.09 & 0.29 & 0.30\\
\bottomrule\end{tabular}
\begin{tablenotes}[flushleft]\small
\item[] Both null and power scenarios are retained. U and A denote unadjusted and adjusted scores. Each comparison uses common trials and reference assignments. The approximation uses lattice tails only for unadjusted binary superiority at $b=0$, and normal tails otherwise. These are magnitudes of estimated differences, not confidence bounds. Scenario-specific paired standard errors and discordant counts are in the computational record.
\end{tablenotes}\end{threeparttable}\end{table}

\begin{table}[!htbp]\centering
\caption{Unadjusted continuous null results for the motivating five-factor setting ($n=400$, $p_{\rm bc}=.95$, first-factor deviations; 10,000 trials).}
\label{tab:supp-pilot-target}
\begin{threeparttable}\small
\begin{tabular}{@{}lrrrr@{}}
\toprule
Null boundary & RT (\%) & Original R--RT & Recon. R--RT & Recon. CRT (\%)\\
\midrule
Superiority & 4.21 & -0.69 (0.10) & -0.03 (0.09) & 4.83\\
Non-inferiority & 2.08 & -0.44 (0.08) & -0.11 (0.07) & 2.43\\
Equivalence, lower & 4.37 & -0.50 (0.08) & +0.02 (0.08) & 5.01\\
Equivalence, upper & 4.70 & -0.59 (0.10) & +0.06 (0.09) & 5.12\\
\bottomrule\end{tabular}
\begin{tablenotes}[flushleft]\small
\item[] Differences are percentage points with paired Monte Carlo standard errors. This setting was identified during development, not selected as the best result of the 10,000-trial evaluation. The nominal level is 2.5\% for non-inferiority and 5\% otherwise. CRT uses estimated stratum effects. Equivalence uses both one-sided components. All configurations are retained in the computational record.
\end{tablenotes}\end{threeparttable}\end{table}

\begin{table}[!htbp]\centering
\caption{Ranges of null rejection rates (\%) across scenario--score combinations, using 10,000 trials per scenario.}
\label{tab:supp-pilot-null}
\begin{threeparttable}\small
\begin{tabular}{@{}llrrrrr@{}}
\toprule
Family/outcome & Level & RT & CRT original & CRT recon. & S recon. & R recon.\\
\midrule
Factorial, continuous & 5 & 4.23--5.48 & 4.69--5.44 & 4.69--5.44 & 4.60--5.48 & 4.22--5.52\\
Broad, continuous & 5 & 3.76--5.34 & 4.46--5.41 & 4.46--5.41 & 4.40--5.40 & 3.60--5.33\\
Broad, continuous & 2.5 & 1.92--2.63 & 2.22--2.66 & 2.22--2.68 & 2.27--2.75 & 1.97--2.59\\
Broad, binary & 5 & 3.45--5.35 & 3.45--5.40 & 3.45--5.40 & 4.45--5.51 & 2.91--5.47\\
Broad, binary & 2.5 & 1.89--2.65 & 2.03--2.69 & 2.03--2.69 & 1.97--2.72 & 1.75--2.67\\
Controls, continuous & 5 & 4.70--5.37 & 4.72--5.44 & 4.72--5.44 & 4.61--5.30 & 4.38--5.25\\
Controls, binary & 5 & 3.70--5.53 & 3.71--5.62 & 3.71--5.62 & 4.75--5.55 & 3.14--5.41\\
\bottomrule\end{tabular}
\begin{tablenotes}[flushleft]\small
\item[] CRT uses simulated reference assignments and estimated effects; S uses normal sampling tails. R uses the fixed rule: lattice probabilities for unadjusted binary superiority at $b=0$, normal tails otherwise. R targets RT, not nominal size. Ranges combine both scores and effect directions where applicable and are not simultaneous confidence intervals. Equivalence records rejection of the intersection--union test at a boundary, not one component.
\end{tablenotes}\end{threeparttable}\end{table}
% END GENERATED VALIDATION APPENDIX
\clearpage

\begin{spacing}{1.1}
\bibliographystyle{apalike}
\bibliography{main}
\end{spacing}

\end{document}
